\documentclass[11pt,a4paper]{article}

% No journal template: this is a Supporting Information document, not a manuscript.
\usepackage[T1]{fontenc}
\usepackage[utf8]{inputenc}
\usepackage{lmodern}
\usepackage[a4paper,margin=2.2cm]{geometry}
\usepackage{amsmath}
\usepackage{amssymb}
\usepackage{array}
\usepackage{longtable}
\usepackage{microtype}

% hyperref goes LAST among packages, and hidelinks keeps the printed page clean:
% the SI is read on paper as often as on screen, and coloured link boxes on every
% cross-reference would be noise.  It is here for the PDF metadata (a file with an
% empty /Title and /Author is a file a reader cannot cite from a PDF library) and
% for the internal links the cross-references already imply.
\usepackage[hidelinks]{hyperref}
\hypersetup{
  pdftitle={Supporting Information --- From Loss Accounting to Feedback Loops: A Coupling Framework for Electrochemical Aging in Lithium-Based Batteries},
  pdfauthor={Huy Hoang Le; Kim-Anh Nguyen},
  pdfsubject={Supporting Information to a Batteries and Supercaps Perspective},
  pdfkeywords={coupling; degradation; feedback loops; knee; lithium-ion batteries; supporting information},
}

\setlength{\parindent}{0pt}
\setlength{\parskip}{5pt}
\renewcommand{\arraystretch}{1.15}
\setlength{\LTpre}{6pt}
\setlength{\LTpost}{6pt}
\allowdisplaybreaks

% Standard remedy for the occasional overfull hbox in justified text: a third pass is
% allowed extra stretch, so TeX can loosen interword space instead of running into the
% margin.  This is a typographic parameter, not a content change, and it is the
% documented cure for "Overfull \hbox ... in paragraph" where feasible breakpoints
% exist but the paragraph cannot be set within tolerance.
\setlength{\emergencystretch}{3em}

\title{Supporting Information\\[4pt]
\large From Loss Accounting to Feedback Loops: A Coupling Framework for
Electrochemical Aging in Lithium-Based Batteries}
\author{Huy Hoang Le \and Kim-Anh Nguyen}
\date{}

\begin{document}
\maketitle
\vspace{-1.2em}
\begin{center}
\small \textit{Batteries \& Supercaps} (Wiley-VCH) --- Perspective --- Supporting Information
\end{center}
\medskip



\subsection*{What this document contains}


A lithium-ion cell usually does not fail gradually. It loses capacity slowly for hundreds of cycles and then turns over and drops. That turnover is the knee. Its position differs from cell to cell and from protocol to protocol, and it is usually explained as a fast fade. A fast fade, however, cannot produce a sharp transition on its own. A sharp transition needs a process whose output feeds its own input.


Four such feedback loops are examined in the accompanying Perspective. Each is written as a chain that returns to its own start:


\begin{itemize}
  \item Loop A. Plating causes dead lithium, dead lithium carries fresh interphase, fresh interphase raises the local overpotential, and the higher overpotential drives more plating. The loop closes when the interphase resistance exceeds the plating threshold.
  \item Loop B. Transition-metal dissolution poisons the anode, which destabilizes the interphase, which consumes more lithium inventory and drives further dissolution. Acid generation is autocatalytic, which is what closes the loop.
  \item Loop C. Fracture exposes fresh surface, interphase grows on that surface, the interphase occupies pore volume and becomes load-bearing, and the changed local constraint drives further fracture.
  \item Loop D. Electrolyte depletion raises transport limitation, which concentrates current and raises local overpotential, which drives more plating and more interphase, which consumes more electrolyte.
\end{itemize}


The document addresses whether a knee follows from the loop, how large each coupling is, which measurement identifies which mechanism, and in what order the measurements apply. Each part is tabulated against the corresponding question.


\begin{enumerate}
  \item S1, whether a knee necessarily follows from Loop A. A two-variable model is built that contains the loop, and then the loop is removed. If the knee survives that removal, the mechanism claim is empty. It does not survive: the knee disappears rather than moving later, so the knee is a consequence of the loop's structure rather than of the fade being fast.
  \item S2, how large each coupling is compared with the effect it drives. A coupling is only worth modeling if its margin is of order unity. Each loop is reduced to a ratio of two quantities that can be estimated independently. Every ratio lands between roughly 1 and 100, which is what makes the mechanisms comparable rather than separable, and it is also why none of them can be dropped.
  \item S3, which measurement identifies which mechanism. A capacity-fade curve integrates all loss mechanisms at once and cannot, by itself, attribute the loss. Twenty established techniques are tabulated against the questions that attribution requires. The order in which they should be applied is the contribution, not the list.
  \item S4, in what order the measurements should be applied to one cell. The tree fixes the order: identify the loss term, locate the electrode, assign the mechanism, then perturb the loop and observe whether the knee moves. It terminates in a perturbation test, because coupling is the one claim that a static measurement cannot settle.
\end{enumerate}


\subsection*{Contents}


\begingroup\small
\setlength{\tabcolsep}{3pt}
\renewcommand{\arraystretch}{1.12}
\sloppy
\begin{tabular}{@{}>{\raggedright\arraybackslash}p{0.108\textwidth}>{\raggedright\arraybackslash}p{0.820\textwidth}@{}}
\hline
\textbf{Section} & \textbf{Contents} \\ \hline
\textbf{S1} & The Loop-A model: state variables, cycle-update equations, the plating threshold, the two self-limitation terms, parameters, numerical method, and results (baseline, the ablation, sensitivity sweep, secondary knee criteria), two applications of the framework to published fade curves (303 cells and 19 fitted exponents), and an explicit list of what the model does not claim. \\
\hline
Table S1 & Mechanism \ensuremath{\times} measurement matrix: 20 techniques against the three questions, with what each can and cannot separate (in S3). \\
\hline
Table S2 & Parameters of the reduced model: 12 parameters with values, units, and provenance. \\
\hline
Table S3 & Baseline landmarks, the ablation, and the secondary knee criteria. \\
\hline
Table S4 & Sensitivity of the knee to the loop gain: eight values spanning 400\ensuremath{\times}. \\
\hline
Table S5 & Operating-condition map: knee cycle count against C-rate and interphase-growth factor, 20 cells. \\
\hline
Table S6 & The framework's knee-to-EOL ratio against the published range: 303 cells from 17 studies (in S1.13). \\
\hline
\textbf{S2} & The derivations behind the four coupling ratios: interphase ohmic drop and its conductivity band; the plating threshold and the thickness it implies; interphase volume and pore filling; Coulombic-efficiency arithmetic and the cycle life it fixes; the Loop-B catalytic shortfall; the Loop-D salt budget; the Loop-C stress margin. Plus the dissolution unit-ambiguity worked example and the voltage-fade separability arithmetic. \\
\hline
\textbf{S3} & Table S1, a 20-row matrix of measurement against mechanism, giving for each technique what it constrains and what it does not, and which loop it addresses; and a 20-item reporting checklist in three tiers. \\
\hline
\textbf{S4} & The diagnostic decision tree, branches B0--B7: node inventory, edge semantics, and a walkthrough of each branch naming the confound it resolves. The tree is not a figure; it exists only as this protocol. \\
\hline
\end{tabular}
\endgroup



\subsection*{Conventions}


\begin{itemize}
  \item Status tags. \texttt{[DERIVED]} = computed here, arithmetic shown. \texttt{[FROM LITERATU}\hspace{0pt}\texttt{RE: DOI]} = supplied by a published source, DOI given. \texttt{[HYPOTHESIS]} = proposed, not established. \texttt{[ASSUMED]} = an order-of-magnitude figure used as input and not traceable to a specific measurement. No number appears without one of these tags.
  \item Units. Areal quantities in mAh cm\ensuremath{^{-}}\ensuremath{^{2}}, mol cm\ensuremath{^{-}}\ensuremath{^{2}}, \ensuremath{\mu}L cm\ensuremath{^{-}}\ensuremath{^{2}}, \ensuremath{\mu}m; resistance in \ensuremath{\Omega} cm\ensuremath{^{2}}; current density in mA cm\ensuremath{^{-}}\ensuremath{^{2}} in the text, but A cm\ensuremath{^{-}}\ensuremath{^{2}} inside the threshold relation, where the conversion is shown explicitly each time; stress in GPa; conductivity in S cm\ensuremath{^{-}}\ensuremath{^{1}}.
  \item Cycle counting. Cycle 0 is the state after formation. All "cycle N" values are post-cycle-N states.
  \item Knee. Primary definition, used unless stated otherwise: cumulative lithium-inventory loss reaching 20 \% of the initial capacity. Four secondary criteria are also reported, because the 20 \% convention is arbitrary; the spread is disclosed rather than hidden.
  \item Precision. A "\textasciitilde{}" means the input is uncertain by at least a factor of 2. Where three significant figures are given, the input is a physical constant and the precision is real; it is the model it enters that is coarse.
  \item Judgement calls are flagged as \texttt{[HYPOTHESIS]} or \texttt{[ASSUMED]}. A reader should never have to guess which numbers are established and which are assumed.
\end{itemize}


\subsection*{Symbols and abbreviations}


\begingroup\small
\setlength{\tabcolsep}{3pt}
\renewcommand{\arraystretch}{1.12}
\sloppy
\begin{longtable}{@{}>{\raggedright\arraybackslash}p{0.163\textwidth}>{\raggedright\arraybackslash}p{0.765\textwidth}@{}}
\hline
\textbf{Term} & \textbf{Meaning} \\ \hline
\endhead
CE & Coulombic efficiency, the fraction of charge passed that is recovered on discharge. \\
CEI & Cathode electrolyte interphase. \\
cryo-EM & Cryogenic electron microscopy, used to image beam-sensitive lithium deposits without damaging them. \\
dead lithium & Metallic lithium that has plated but lost electrical contact, so it can no longer be cycled. \\
DEMS & Differential electrochemical mass spectrometry, which measures gas evolved during cycling. \\
DRT & Distribution of relaxation times, a deconvolution of impedance spectra that avoids assigning arcs to a circuit model in advance. \\
EIS & Electrochemical impedance spectroscopy. \\
HPC & High-precision coulometry, the measurement of Coulombic efficiency to a few parts per million. \\
ICP-MS & Inductively coupled plasma mass spectrometry, used to quantify dissolved transition metal. \\
interphase & The general term for a passivating film between an electrode and the electrolyte: the SEI on the anode, the CEI on the cathode. Where a quantity is said to be per unit interphase, the film is treated as a slab of uniform resistivity. \\
knee & The cycle at which capacity begins to fall sharply rather than linearly. Defined operationally in Conventions above. \\
KL & Kinetic loss: the part of capacity fade caused by increasing polarization rather than by lost inventory or active material. \\
LAM & Loss of active material: the electrode can no longer store the lithium it once could. \\
LLI & Loss of lithium inventory: cyclable lithium has been consumed and is no longer available. \\
Loop A--D & The four feedback loops defined in the opening section of this document. \\
M1--M4 & The four mechanisms the loops are built from: M1 interphase growth, M2 dead-lithium formation, M3 transition-metal dissolution and surface reconstruction, M4 particle fracture. \\
MAS-NMR & Magic-angle-spinning nuclear magnetic resonance, used to resolve lithium environments in the bulk. \\
N/P & Negative-to-positive capacity ratio of a cell. \\
pp & Percentage points, the arithmetic difference between two percentages. \\
Q1--Q3 & The three attribution questions of S3.1: which loss term moved (Q1), which electrode is responsible (Q2), which mechanism family is responsible (Q3). \\
R* & Plating threshold: the interphase resistance at which the anode potential reaches 0 V versus Li/Li\ensuremath{^{+}} at a given current density. \\
SEI & Solid electrolyte interphase, the passivating film that forms on the anode. \\
SOC & State of charge. \\
SRL & Surface reconstruction layer, the disordered rock-salt surface phase left by transition-metal dissolution. \\
STEM-EELS & Scanning transmission electron microscopy with electron energy-loss spectroscopy, a spatially resolved probe of composition and bonding. \\
TGC & Titration gas chromatography, which measures the metallic lithium fraction of a cell's inactive lithium. \\
TL & Thermodynamic, or potential, loss: the part of energy fade caused by a shift in the cell's voltage curve rather than by capacity loss. \\
VRP & Voltage relaxation profile, measured in the rest period after charge. \\
XRD & X-ray diffraction, used here operando to track lattice parameters during cycling. \\
\hline
\end{longtable}
\endgroup



\subsubsection*{Mathematical symbols}


Greek letters and constants used in S1 and S2, which are otherwise introduced only at their first equation. Symbols local to one derivation are defined there.


\begingroup\small
\setlength{\tabcolsep}{3pt}
\renewcommand{\arraystretch}{1.12}
\sloppy
\begin{longtable}{@{}>{\raggedright\arraybackslash}p{0.154\textwidth}>{\raggedright\arraybackslash}p{0.641\textwidth}>{\raggedright\arraybackslash}p{0.122\textwidth}@{}}
\hline
\textbf{Symbol} & \textbf{Meaning} & \textbf{First used} \\ \hline
\endhead
\emph{R}* & Critical anode area-specific resistance for plating onset, (\ensuremath{\Delta} \ensuremath{-} \ensuremath{\eta}\textsubscript{ct})/\emph{j}. Rate-dependent through \emph{j}; \ensuremath{\Delta} itself is a material property (S1.16). & S1.4, S2.2 \\
\emph{R}\textsubscript{SEI} & Area-specific resistance of the anode interphase. & S1.2, S2.1 \\
\ensuremath{\Delta} & Margin between the anode operating potential and 0 V vs Li/Li\ensuremath{^{+}}. & S1.4, S2.2 \\
\ensuremath{\eta}\textsubscript{ct} & Charge-transfer overpotential at the anode. & S1.4, S2.2 \\
\emph{j} & Local current density, set by nominal capacity and C-rate. & S1.2, S2.2 \\
\ensuremath{\sigma} & Ionic conductivity of the interphase, taken as a band of 10\ensuremath{^{-}}\ensuremath{^{9}}--10\ensuremath{^{-}}\ensuremath{^{7}} S cm\ensuremath{^{-}}\ensuremath{^{1}}. & S2.1 \\
\emph{d} & Interphase thickness; \emph{d}* is the thickness at which \emph{R}\textsubscript{SEI} reaches \emph{R}*. & S2.1, S2.2 \\
\emph{b}\textsubscript{fresh} & Loop gain: resistance added per unit of fresh interphase area. The only unanchored parameter. & S1.7 \\
\emph{a}\textsubscript{SEI} & Interphase area created per unit of lithium inventory lost. & S1.5 \\
\emph{f}\textsubscript{dead} & Fraction of plated lithium that becomes electrically isolated. Measured; 0.094 without FEC. & S1.5 \\
\emph{f} & Inefficiency, 1 \ensuremath{-} CE, as a fraction of charge per cycle. & S2.4 \\
\emph{Q}\ensuremath{_{0}} & Nominal areal capacity, taken as 3.0 mAh cm\ensuremath{^{-}}\ensuremath{^{2}}. & S2.3 \\
\emph{V}\textsubscript{m} & Molar volume of the interphase phase. & S2.3 \\
\emph{E}, \ensuremath{\nu} & Young's modulus and Poisson ratio of the cathode secondary particle. & S2.7 \\
\emph{K}\textsubscript{Ic} & Mode-I fracture toughness of the cathode secondary particle. & S2.7 \\
\ensuremath{\Omega} & Partial molar volume of lithium in the host. & S2.7 \\
\emph{c}\textsubscript{LLI} & Areal density of lost lithium inventory. & S1.3 \\
\hline
\end{longtable}
\endgroup



\section*{S1. The reduced Loop-A model: equations, parameters, ablation, and limits}


\subsection*{S1.1 Scope of the model}


Loop A is a verbal argument: plating creates dead lithium, dead lithium carries fresh interphase, fresh interphase raises the local overpotential, and the higher overpotential drives more plating. Read as a sentence, it is plausible. Read as a mechanism, it makes a much stronger claim that a sentence cannot settle: that a capacity knee follows necessarily from the loop, and that the knee is therefore a property of the loop's structure rather than a symptom of the fade being fast.


Those two readings make different predictions, and only one of them can be tested. If the knee is merely what a fast fade looks like, then any cell that fades fast enough will show one, and removing the loop should at most move the knee later. If the knee is a threshold crossing inside the loop, then removing the loop should remove the knee, not delay it. This section builds the smallest model that contains the loop, runs it, and then removes the loop.


The model is a scaling argument, not a battery model. It is built to make that comparison possible, including the possibility that the answer is no.


The formulation has three properties that constrain the numbers in this section.


\begin{enumerate}
  \item Nothing is fitted to a dataset. No parameter was obtained by regression against measured capacity fade. Two parameters are taken from measurement (\emph{f}\textsubscript{dead}, via the cited DOI) and one is anchored to bulk transport (\texttt{a\_SEI}); the rest are order-of-magnitude placeholders chosen to be physically defensible. They are labelled individually in Table S2.
  \item No absolute number here is a prediction of any real cell's cycle life. The numbers that are load-bearing are comparative and structural: the knee exists in the coupled case and does not exist in the severed case; the pre-knee fade is approximately linear and the post-knee fade is superlinear; the knee onset moves monotonically with operating condition. These four statements survive the parameter choices, and Table S4 (the 400\ensuremath{\times} sweep) is the evidence that they do.
  \item The model represents Loop A only. Loops B, C, and D are not represented; the corresponding parts of Figure 1 panel (a) are qualitative. Loop D in particular is absent from the reduced model: no electrolyte inventory, no salt transport, so nothing here bears on it. The model also has no spatial dimension, no porosity field and no electrode thickness, so it cannot in principle support Loop D. Section S1.10 states this and the other limitations explicitly.
\end{enumerate}


\par\noindent\rule{\linewidth}{0.4pt}\par


\subsection*{S1.2 State variables}


A model of a feedback loop has to carry, at minimum, the quantity that drives the loop and the quantity that the loop consumes. Here the driver is the anode interphase resistance, because that is what pushes the anode potential toward the plating threshold, and what the loop consumes is cyclable lithium. The model therefore tracks exactly two state variables per unit electrode area.


If the loss channels were independent, the natural model would carry one variable per mechanism and sum them. Carrying two coupled variables instead is the modeling statement that they are not independent, and it is the assumption the ablation in S1.5 is designed to test.


\begingroup\small
\setlength{\tabcolsep}{3pt}
\renewcommand{\arraystretch}{1.12}
\sloppy
\begin{tabular}{@{}>{\raggedright\arraybackslash}p{0.095\textwidth}>{\raggedright\arraybackslash}p{0.251\textwidth}>{\raggedright\arraybackslash}p{0.119\textwidth}>{\raggedright\arraybackslash}p{0.440\textwidth}@{}}
\hline
\textbf{Symbol} & \textbf{Meaning} & \textbf{Unit} & \textbf{Physical interpretation} \\ \hline
\emph{R} & anode interphase area-specific ionic resistance & \ensuremath{\Omega}\ensuremath{\cdot}cm\ensuremath{^{2}} & lumped resistance of the SEI plus any in-series interphase on the anode; the quantity that pushes the anode potential toward 0 V vs Li/Li\ensuremath{^{+}} at fixed current \\
\hline
\emph{LLI} & cumulative lost lithium inventory & mAh\ensuremath{\cdot}cm\ensuremath{^{-}}\ensuremath{^{2}} & charge that can no longer be cycled; the capacity-loss channel the model resolves \\
\hline
\end{tabular}
\endgroup



Two derived quantities are tracked but are not state variables: the plating fraction \emph{f}(\emph{n}) (a diagnostic of how far into the loop the cell is) and the dead-lithium accumulation (the portion of plated Li that becomes electrically isolated). Neither feeds back except through \emph{R} and \emph{LLI}.


\textbf{Choice of state variables: \emph{R} and \emph{LLI}, rather than \emph{R} and capacity.} Capacity loss would fold the LLI channel together with any loss of active material, and the model has no LAM. Keeping \emph{LLI} explicit means the model's output is directly comparable to the \emph{LLI} axis of the energy balance, which is the only axis it is entitled to address.


\par\noindent\rule{\linewidth}{0.4pt}\par


\subsection*{S1.3 Update equations}


Discrete cycle update, applied once per cycle. The index \emph{n} \ensuremath{\geq} 1 runs over cycles; cycle 0 is the post-formation state. \texttt{i} is used below for the same index inside the summation term.


Current density is set by the nominal capacity and the C-rate, and is constant during a test:


\begin{equation}j = \frac{c_{\text{rate}}\,Q_0}{1000}\quad[\text{A cm}^{-2}] \tag{S1.E1}\end{equation}


with \emph{Q}\ensuremath{_{0}} in mAh cm\ensuremath{^{-}}\ensuremath{^{2}}, so that \emph{c}\textsubscript{rate} = 1 gives \emph{j} = 3.0 mA cm\ensuremath{^{-}}\ensuremath{^{2}}. This point is easy to get wrong and materially so: if \emph{j} is allowed to fall as capacity falls, the feedback is throttled for a spurious reason and a spurious knee appears. Current density is therefore held at the nominal value for the whole test.


\textbf{Excess anode overpotential.} The anode is driven toward 0 V vs Li/Li\ensuremath{^{+}} by the ohmic drop across the interphase. The quantity that matters is the margin left over after the charge-transfer overpotential:


\begin{equation}\eta(n) = j\,R(n-1) + \eta_{\text{ct}} - \Delta \quad[\text{V}] \tag{S1.E2}\end{equation}


where \ensuremath{\Delta} is the operating margin between the anode's equilibrium potential and 0 V vs Li/Li\ensuremath{^{+}}, and \emph{\ensuremath{\eta}}\textsubscript{ct} the charge-transfer overpotential at the operating rate. Plating begins when \emph{\ensuremath{\eta}} > 0.


\textbf{Plating fraction.} Plating turns on smoothly rather than as a step, because a real electrode has a distribution of local \emph{R} across its area, so different patches cross the threshold at different times:


\begin{equation}f(n) = \begin{cases} f_{\max}\left[1-\exp\left(-k_{\text{pl}}\,\eta(n)\right)\right] & \eta(n) > 0\\[2pt] 0 & \eta(n) \le 0\end{cases} \tag{S1.E3}\end{equation}


\emph{f} is the fraction of the applied charge that goes to Li(0) rather than to intercalation. The saturating form is the first self-limitation term: mass balance guarantees \emph{f} \ensuremath{\leq} \emph{f}\textsubscript{max} < 1, so plating can never consume more charge than is passed. This is the structural fix for that failure mode, and it is why the model has a finite end of life rather than a numerical divergence.


\textbf{Plated charge.}


\begin{equation}p_{\text{plate}}(n) = f(n)\,Q_{\text{applied}},\qquad Q_{\text{applied}} = c_{\text{rate}}Q_0 \quad[\text{mAh cm}^{-2}] \tag{S1.E4}\end{equation}


\textbf{Dead-lithium increment.} A measured fraction of plated lithium becomes electrically isolated:


\begin{equation}d_{\text{dead}}(n) = f_{\text{dead}}\,p_{\text{plate}}(n) \tag{S1.E5}\end{equation}


\textbf{Interphase resistance growth.} Two contributions, plus pore-filling saturation:


\begin{equation}R(n) = R(n-1) + \left[\frac{a_{\text{SEI}}}{\sqrt{n}}\,T_{\text{factor}} + b_{\text{fresh}}\,d_{\text{dead}}(n)\right]\cdot h(n) \tag{S1.E6}\end{equation}


\begin{equation}h(n) = \max\!\left(0,\;1-\frac{R(n-1)-R_0}{R_{\text{cap}}}\right) \tag{S1.E7}\end{equation}


The first term in brackets is the diffusion-limited baseline: interphase thickness grows as \ensuremath{\surd}\emph{t}, so its resistance grows as \ensuremath{\surd}\emph{n}. It is present whether or not the loop is active, and it is the reason the severed control cell still fades. The second term is the loop-closing arrow: fresh, unpassivated interphase forms on each cycle's newly created dead-lithium surface, and \emph{more dead Li} therefore means \emph{more interphase}, which means \emph{higher R}, which means \emph{more plating} next cycle. The coefficient \emph{b}\textsubscript{fresh} is the loop gain. \emph{h}(\emph{n}) is the second self-limitation term: as the anode pore volume fills with interphase, further growth is geometrically blocked, so \emph{R} asymptotes to \emph{R}\ensuremath{_{0}} + \emph{R}\textsubscript{cap.} Without \emph{h} the loop would be unbounded even with \emph{f} \ensuremath{\leq} \emph{f}\textsubscript{max}, because \emph{R} could still grow without limit.


\emph{T}\textsubscript{factor} is an Arrhenius proxy that multiplies the SEI-growth prefactor; \emph{T}\textsubscript{factor} = 2 corresponds to roughly +20 K. It is a crude lumping, used only to generate the sensitivity map in Table S4.


\textbf{Loss inventory.}


\begin{equation}LLI(n) = LLI(n-1) + c_{LLI}\,\Delta R(n) + d_{\text{dead}}(n) \tag{S1.E8}\end{equation}


with \ensuremath{\Delta}\emph{R}(\emph{n}) = \emph{R}(\emph{n}) \ensuremath{-} \emph{R}(\emph{n}\ensuremath{-}1). Two contributions to lost inventory: charge consumed in growing the interphase (\emph{c}\textsubscript{LLI} per \ensuremath{\Omega}\ensuremath{\cdot}cm\ensuremath{^{2}} of new interphase) and charge stranded as dead lithium. Note that \emph{c}\textsubscript{LLI}\ensuremath{\cdot}\ensuremath{\Delta}\emph{R} is charged on the whole resistance increment, including the \ensuremath{\surd}\emph{n} baseline term, so even the severed cell loses inventory, which is required for the ablation in S1.5 to be a fair test.


\textbf{End of life.} The cell is dead when LLI reaches \emph{Q}\ensuremath{_{0}}. LLI is clipped at \emph{Q}\ensuremath{_{0}} because capacity cannot be negative; the model stops at that cycle. Onset of the knee is a separate, earlier event, defined in S1.6.


\par\noindent\rule{\linewidth}{0.4pt}\par


\subsection*{S1.4 The plating threshold, R*}


The state equations say that plating begins once the excess overpotential turns positive. That condition can be rearranged into a threshold on the interphase resistance, and the rearrangement is worth doing because it separates what the cell's chemistry controls from what the operator controls. Plating begins when \emph{\ensuremath{\eta}} > 0 in (S1.E2):


\begin{equation}\boxed{\;R > R^{*} \equiv \frac{\Delta - \eta_{\text{ct}}}{j}\;}\tag{S1.E9}\end{equation}


This is the single inequality on which the threshold interpretation of the knee rests. It is worth being precise about what it does and does not depend on. \emph{R}* is a function of \emph{j} (rate, hence also of temperature through the kinetics) and \ensuremath{\Delta} (operating SOC window, and weakly temperature through the equilibrium potential). \emph{R}* is not itself a tabulated property value like a diffusion coefficient or a modulus -- it is assembled from rate and margin -- and that is the formal content of calling knee onset a design variable. But \ensuremath{\Delta} is a material property (S1.16): for graphite it is the 0.08-0.10 V stage-1 plateau relative to 0 V vs Li/Li\ensuremath{^{+}}, so the margin is only as adjustable as the anode chemistry permits. The precision matters, because one input is only partly an operating choice. \ensuremath{\Delta} is the anode's headroom above 0 V, and that headroom is set by the anode's equilibrium potential, which is a material property: about 0.08-0.10 V for the stage-1 plateau of graphite (S1.16). A cell chemist can therefore raise \emph{R}* by cycling higher in the SOC window, by lowering the rate, or by changing the anode material itself. The last is the most effective and the least available, since it means a different cell chemistry and not a tuning of the existing one. A material that delays the knee without moving \ensuremath{\Delta} does so by keeping \emph{R}(\emph{n}) below \emph{R}*, which is a different and slower lever.


With the reference parameters, at 1C: \emph{j} = 3.0 mA cm\ensuremath{^{-}}\ensuremath{^{2}} = 0.0030 A cm\ensuremath{^{-}}\ensuremath{^{2}}, so


\begin{equation}R^{*} = \frac{0.090 - 0.030}{0.0030} = 20.0\ \Omega\,\text{cm}^{2} \nonumber\end{equation}


The fresh cell starts at \emph{R}\ensuremath{_{0}} = 8.0 \ensuremath{\Omega}\ensuremath{\cdot}cm\ensuremath{^{2}}, i.e. at 40 \% of the threshold, with 12 \ensuremath{\Omega}\ensuremath{\cdot}cm\ensuremath{^{2}} of headroom. The cell does not plate at cycle 1; it plates when the slow \ensuremath{\surd}\emph{n} baseline plus the loop term carry \emph{R} past 20 \ensuremath{\Omega}\ensuremath{\cdot}cm\ensuremath{^{2}}. Under the reference parameters that happens at cycle 220 (first cycle with \emph{f} > 10\ensuremath{^{-}}\ensuremath{^{5}}), which is well before the capacity knee at cycle 530, an important ordering that reproduces the qualitative observation that plating is detectable long before the capacity curve turns over.


\emph{R}* scales inversely with rate (this table is inline here rather than in the S1.11 table block, because it is read as part of the threshold argument):


\begingroup\footnotesize
\setlength{\tabcolsep}{3pt}
\renewcommand{\arraystretch}{1.12}
\sloppy
\begin{tabular}{@{}>{\raggedright\arraybackslash}p{0.115\textwidth}>{\raggedright\arraybackslash}p{0.144\textwidth}>{\raggedright\arraybackslash}p{0.156\textwidth}>{\raggedright\arraybackslash}p{0.158\textwidth}>{\raggedright\arraybackslash}p{0.156\textwidth}>{\raggedright\arraybackslash}p{0.151\textwidth}@{}}
\hline
\textbf{C-rate} & \textbf{\emph{j} (mA cm\ensuremath{^{-}}\ensuremath{^{2}})} & \textbf{\emph{R}* (\ensuremath{\Omega}\ensuremath{\cdot}cm\ensuremath{^{2}})} & \textbf{\emph{R}*/\emph{R}\ensuremath{_{0}}} & \textbf{knee (cycle)} & \textbf{EOL (cycle)} \\ \hline
0.5 & 1.500 & 40.00 & 5.00 & 2651 & 4570 \\
\hline
1.0 & 3.000 & 20.00 & 2.50 & 530 & 1044 \\
\hline
1.5 & 4.500 & 13.33 & 1.67 & 204 & 483 \\
\hline
2.0 & 6.000 & 10.00 & 1.25 & 104 & 296 \\
\hline
2.5 & 7.500 & 8.00 & 1.00 & 64 & 210 \\
\hline
\end{tabular}
\endgroup



The last row is the limiting case: at 2.5C the threshold \emph{equals the fresh-cell resistance}, so \emph{R}*/\emph{R}\ensuremath{_{0}} = 1 and the cell is at the plating condition from cycle 1. The model then has no delay to absorb, and the knee arrives at cycle 64. This illustrates the rate dependence.


A limitation of this table must be read with it. The departure detector of S1.5 fits the early power law over cycles 30-200, which it can only do if plating has not begun inside that window. At 1C plating starts at cycle 220, so the fit is genuinely pre-loop. At 1.5C, 2.0C and 2.5C plating starts at cycles 47, 10 and 2 respectively, all inside the window, so for those three rows the fitted baseline already contains loop feedback and the reported knee is a property of the detector's calibration rather than a clean measurement. The R* column is unaffected, since it is arithmetic; the knee and EOL columns for 1.5-2.5C should be read as model output, not as detectable knees. This is the same contamination rule that governs the margin sweep in S1.16, where only delta = 0.090-0.120 V survives.


A caveat that belongs next to (S1.E9). \emph{R}* uses a single lumped \ensuremath{\Delta} and a single lumped \emph{R}, so it is an electrode-average statement. A real electrode has a potential distribution across its area, and local \emph{R} can exceed \emph{R}* in some patches long before the average does. The threshold therefore predicts \emph{onset of detectable plating}, not \emph{first plating event}, the two differ by the width of the distribution, which the model does not contain. No stronger claim is supported.


\par\noindent\rule{\linewidth}{0.4pt}\par


\subsection*{S1.5 Ablation: severance of the coupling}


This is the load-bearing test: if the knee survives severing the coupling, the mechanism claim is empty. The ablation is constructed so that it could show that.


\textbf{Severing method.} Two terms in (S1.E6) and (S1.E5) carry the loop: the \emph{b}\textsubscript{fresh} arrow and the \emph{f}\textsubscript{dead} fraction. Severing both (\emph{b}\textsubscript{fresh} = 0 and \emph{f}\textsubscript{dead} = 0) leaves a cell with identical \ensuremath{\surd}\emph{n} interphase growth but no feedback from plated lithium back into resistance. The cell still fades --- the \ensuremath{\surd}\emph{n} term is untouched, so the comparison is against a fading control.


\textbf{Departure detection.} A knee is detected as a \emph{departure from the cell's own early power law} rather than as an absolute capacity. The early fade is fitted to \emph{LLI} = \emph{A}\ensuremath{\cdot}\emph{n}\textsuperscript{p} over cycles 30--200 (before the first plating event at cycle 220 in the coupled case, so the fit window is genuinely pre-loop), and the departure point is the first cycle after 200 at which LLI exceeds the extrapolated baseline by more than 10 \%. The criterion is conservative in two ways: it requires the cell to exceed its own trend rather than a fixed threshold, and it requires a 10 \% excess rather than any excess.


\begingroup\small
\setlength{\tabcolsep}{3pt}
\renewcommand{\arraystretch}{1.12}
\sloppy
\begin{tabular}{@{}>{\raggedright\arraybackslash}p{0.331\textwidth}>{\raggedright\arraybackslash}p{0.344\textwidth}>{\raggedright\arraybackslash}p{0.242\textwidth}@{}}
\hline
\textbf{Quantity} & \textbf{Coupled} & \textbf{Severed} \\ \hline
Early power-law fit (cycles 30--200) & \emph{LLI} = 7.268 \ensuremath{\times} 10\ensuremath{^{-}}\ensuremath{^{3}} \ensuremath{\cdot} \emph{n}\textsuperscript{0.522} & identical \\
\hline
First plating (cycle) & 220 & 220 \\
\hline
Departure from early power law & cycle 285 & \textbf{none (to \emph{n}\textsubscript{max})} \\
\hline
Knee, LLI \ensuremath{\geq} 20 \% \emph{Q}\ensuremath{_{0}} & cycle 530 & none \\
\hline
End of life, LLI = \emph{Q}\ensuremath{_{0}} & cycle 1044 & not reached by 6000 \\
\hline
\emph{R} at end of run (\ensuremath{\Omega}\ensuremath{\cdot}cm\ensuremath{^{2}}) & 75.1 & 60.5 \\
\hline
Terminal LLI (mAh cm\ensuremath{^{-}}\ensuremath{^{2}}) & 3.00 (clipped at \emph{Q}\ensuremath{_{0}}) & 0.525 \\
\hline
\end{tabular}
\endgroup



The severed cell reaches 6000 cycles at 82.5 \% of \emph{Q}\ensuremath{_{0}} retained, and at the coupled cell's EOL cycle (1044) it still retains 91.6 \%. It has faded to 17.5 \% loss over 5.7 times the coupled cell's life, and it never leaves its early power law.


\textbf{The distinction that matters.} Severing the coupling does not merely delay the knee --- it removes it. The two curves are identical to cycle 200 by construction, so the difference is not a difference in fade rate; it is a difference in \emph{whether a departure ever occurs}. That is why the result bears on the thesis rather than on the parameters: the knee and the fade rate are shown to be independent axes, and the knee is a consequence of the loop's coupling structure, not of the fade being fast.


\textbf{Which term carries the loop.} A natural objection is that the ablation removes more than one factor at once and therefore does not identify the mechanism. Table S3 part B separates them:


\begingroup\footnotesize
\setlength{\tabcolsep}{3pt}
\renewcommand{\arraystretch}{1.12}
\sloppy
\begin{tabular}{@{}>{\raggedright\arraybackslash}p{0.332\textwidth}>{\raggedright\arraybackslash}p{0.192\textwidth}>{\raggedright\arraybackslash}p{0.178\textwidth}>{\raggedright\arraybackslash}p{0.095\textwidth}>{\raggedright\arraybackslash}p{0.095\textwidth}@{}}
\hline
\textbf{Variant} & \textbf{Setting} & \textbf{Departure} & \textbf{Knee} & \textbf{EOL} \\ \hline
Coupled & reference & 285 & 530 & 1044 \\
\hline
Feed-forward arrow severed & \texttt{b\_fresh = 0} & 300 & 670 & 1681 \\
\hline
Dead-lithium fraction severed & \texttt{f\_dead = 0} & none & none & 6000 \\
\hline
Fully severed & both zero & none & none & 6000 \\
\hline
\end{tabular}
\endgroup



The result is asymmetric, and the asymmetry is informative, but it does not by itself separate the two couplings, because the two severing operations are not equivalent in scope. Removing the dead-lithium fraction alone, \texttt{f\_dead = 0}, abolishes the knee outright and gives a run identical to the fully severed case. Removing the feed-forward arrow alone, \texttt{b\_fresh = 0}, only delays the knee, from 530 to 670 cycles.


The reason is structural rather than numerical. Setting \texttt{f\_dead = 0} sends \texttt{d\_dead = 0} in (S1.E5), which removes both terms at once: the resistance contribution \texttt{b\_fresh\ensuremath{\cdot}d\_dead} in (S1.E6) and the stranded-inventory contribution in (S1.E8). The cell is then not a loop at all, which is why it matches the fully severed case exactly. Setting \texttt{b\_fresh = 0} leaves dead lithium forming and stranding inventory, so the fade continues and the knee merely arrives later.


The loop therefore requires a nonzero dead-lithium fraction, and the feed-forward arrow sets when the knee arrives rather than whether it exists. The claim is weaker than an identification: \texttt{f\_dead} is necessary for the loop, and \texttt{b\_fresh} is a gain on one of its arrows. This is the structurally correct answer: \emph{f}\textsubscript{dead} is the property that makes the coupling a loop, while \emph{b}\textsubscript{fresh} is a gain on one of its arrows. It also identifies the measurement that would falsify the model. Anything showing that the fraction of plated lithium becoming electrically isolated is negligibly small would break the loop, whereas a revision of the interphase resistance per unit dead lithium would only move the knee in time.


\textbf{Robustness across \emph{b}\textsubscript{fresh} (claim C2 below).} The severing test is repeated at \emph{b}\textsubscript{fresh} = 10, 30, 60, 100 and 200: the coupled departure ranges 254--294 cycles and the severed case shows no departure at any value. The ablation result is not an artefact of the reference parameter choice.


\textbf{Robustness across rate.} The severing test also runs at 0.5C and 2.0C. At 0.5C the coupled departure is 2380 and the severed departure is none. At 2.0C neither case departs: the coupled cell reaches its 20 \%-LLI knee at cycle 104 before the pre-loop baseline fit can be established, so the departure detector correctly declines to make a claim. That row is reported as invalid rather than omitted, because a detector returning \textquotedblleft{}none\textquotedblright{} for both cases should also be shown to return \textquotedblleft{}some\textquotedblright{} for one of them, and the 0.5C and 1.0C rows do.


\par\noindent\rule{\linewidth}{0.4pt}\par


\subsection*{S1.6 Secondary knee definitions}


The 20 \%-of-\emph{Q}\ensuremath{_{0}} definition of the knee is arbitrary, and the model's conclusion would be weaker than it looks if the knee existed only under that convention. Four independent secondary criteria were therefore computed on the coupled baseline, alongside the primary:


\begingroup\small
\setlength{\tabcolsep}{3pt}
\renewcommand{\arraystretch}{1.12}
\sloppy
\begin{tabular}{@{}>{\raggedright\arraybackslash}p{0.308\textwidth}>{\raggedright\arraybackslash}p{0.442\textwidth}>{\raggedright\arraybackslash}p{0.166\textwidth}@{}}
\hline
\textbf{Criterion} & \textbf{Definition} & \textbf{Knee (cycle)} \\ \hline
Cumulative LLI & LLI \ensuremath{\geq} 0.20 \emph{Q}\ensuremath{_{0}} (primary, used throughout) & 530 \\
\hline
Second derivative, LLI & cycle of maximum d\ensuremath{^{2}}(LLI)/d\emph{n}\ensuremath{^{2}} over the run & 470 \\
\hline
Retention curvature & cycle of maximum curvature of retained capacity vs \emph{n} & 470 \\
\hline
Kneedle chord method & maximum distance from the chord joining first and last point & 479 \\
\hline
Local-exponent acceleration & cycle of fastest growth in local d log(LLI)/d log \emph{n} & 278 \\
\hline
\end{tabular}
\endgroup



Three of the geometric criteria cluster at 470--530; the local-exponent criterion fires earlier, at 278, which is close to the departure point at 285, and that agreement is informative rather than contradictory: the local-exponent method detects \emph{when acceleration begins}, while the others detect \emph{when the acceleration has become large}. Both are legitimate readings of \textquotedblleft{}onset\textquotedblright{}, and they differ by a factor of \textasciitilde{}1.7. The knee is therefore located within 278--530 cycles depending on definition, a band of about a factor of 1.9. The 20 \% convention is adopted because it is the most conservative member of the cluster; suppressing the spread would leave a reader who prefers the second-derivative criterion to conclude that the definition had been chosen to suit the argument.


\textbf{What this does not license.} No primary literature value should be compared numerically against these cycle numbers. They are outputs of an unfitted model with a placeholder loop gain.


\par\noindent\rule{\linewidth}{0.4pt}\par


\subsection*{S1.7 Sensitivity sweep over 400\ensuremath{\times} in \emph{b}\textsubscript{fresh}}


\emph{b}\textsubscript{fresh} is the one parameter with no external anchor, and it is also the parameter that most strongly sets where the knee lands. It is therefore the model's weakest point. The sensitivity is therefore reported in full rather than concealed.


\emph{b}\textsubscript{fresh} carries units \ensuremath{\Omega}\ensuremath{\cdot}cm\ensuremath{^{2}} per (mAh cm\ensuremath{^{-}}\ensuremath{^{2}}) of dead lithium: resistance per unit of newly created fresh interphase. Dividing by an assumed interphase resistivity \ensuremath{\rho} \ensuremath{\approx} 10\ensuremath{^{8}} \ensuremath{\Omega}\ensuremath{\cdot}cm converts that coefficient into a thickness. The dimensional check is worth stating, because the result is a length per unit \emph{areal charge} and not a length: (\ensuremath{\Omega} cm\ensuremath{^{2}})/(mAh cm\ensuremath{^{-}}\ensuremath{^{2}}) divided by (\ensuremath{\Omega} cm) gives cm\ensuremath{^{3}} mAh\ensuremath{^{-}}\ensuremath{^{1}}, which is the same as cm per (mAh cm\ensuremath{^{-}}\ensuremath{^{2}}). The conversion therefore yields \emph{b}\textsubscript{fresh} = 30 \ensuremath{\rightarrow} 3.0 nm, \emph{b}\textsubscript{fresh} = 60 \ensuremath{\rightarrow} 6.0 nm of re-formed interphase per mAh cm\ensuremath{^{-}}\ensuremath{^{2}} of dead lithium. Three nanometers is physically plausible for a re-formed interphase; that is the entire justification for 30, and it is an argument, not a measurement.


Sweep over 1 \ensuremath{\rightarrow} 400 (a 400\ensuremath{\times} range, 2.6 decades). \texttt{C1} is the claim that a knee exists; \texttt{C2} is that severing abolishes it.


\begingroup\scriptsize
\setlength{\tabcolsep}{3pt}
\renewcommand{\arraystretch}{1.12}
\sloppy
\begin{tabular}{@{}>{\raggedright\arraybackslash}p{0.145\textwidth}>{\raggedright\arraybackslash}p{0.133\textwidth}>{\raggedright\arraybackslash}p{0.061\textwidth}>{\raggedright\arraybackslash}p{0.110\textwidth}>{\raggedright\arraybackslash}p{0.061\textwidth}>{\raggedright\arraybackslash}p{0.117\textwidth}>{\raggedright\arraybackslash}p{0.126\textwidth}>{\raggedright\arraybackslash}p{0.105\textwidth}@{}}
\hline
\textbf{\emph{b}\textsubscript{fresh} (\ensuremath{\Omega}\ensuremath{\cdot}cm\ensuremath{^{2}} per mAh cm\ensuremath{^{-}}\ensuremath{^{2}})} & \textbf{implied thickness (nm)} & \textbf{Knee} & \textbf{Departure} & \textbf{EOL} & \textbf{pre-knee exponent} & \textbf{post-knee exponent} & \textbf{\emph{R}\textsubscript{end} (\ensuremath{\Omega}\ensuremath{\cdot}cm\ensuremath{^{2}})} \\ \hline
1.0 & 0.1 & 663 & 299 & 1630 & 0.96 & 1.83 & 40.8 \\
\hline
3.0 & 0.3 & 650 & 298 & 1542 & 0.96 & 1.90 & 43.9 \\
\hline
10.0 & 1.0 & 610 & 294 & 1328 & 0.95 & 2.10 & 53.7 \\
\hline
30.0 & 3.0 & 530 & 285 & 1044 & 0.91 & 2.43 & 75.1 \\
\hline
60.0 & 6.0 & 462 & 275 & 873 & 0.86 & 2.56 & 94.8 \\
\hline
100.0 & 10.0 & 410 & 267 & 768 & 0.81 & 2.54 & 109.4 \\
\hline
200.0 & 20.0 & 348 & 254 & 667 & 0.71 & 2.37 & 123.5 \\
\hline
400.0 & 40.0 & 300 & 244 & 613 & 0.61 & 2.09 & 127.8 \\
\hline
\end{tabular}
\endgroup



\textbf{Conclusions, stated as narrowly as the data allow:}


\begin{itemize}
  \item C1, a knee exists (8/8 values tested). Over a 400\ensuremath{\times} range in the loop gain the knee moves from cycle 663 to 300, a factor of 2.21. The existence of a knee is therefore insensitive to the parameter that sets its position. This is the operative distinction: the argument requires a threshold plus feedback, not a specific gain.
  \item C2 --- severing abolishes it (5/5 values tested, \emph{b}\textsubscript{fresh} = 10--200; the severing test spans a narrower range than C1; the asymmetry is stated rather than smoothed over).
  \item C3, the qualitative shape survives: pre-knee exponent 0.61--0.96 (always \ensuremath{\leq} 1, i.e. never superlinear before the loop engages), post-knee exponent 1.83--2.56 (always > 1). The reference value 30 was chosen partly because 0.91 is the closest to a linear pre-knee fade; lower values give exponent 0.96 (flatter) and higher values give 0.61 (markedly sublinear). The value is a stated choice, not a derived quantity, and is recorded as such.
\end{itemize}


\textbf{What the sweep cannot fix.} The absolute knee position (300--663 cycles) is still set by an unfitted parameter, so the model cannot say where a real cell's knee will be. What it does establish is that a knee of some position follows from the coupling, and that its position is controlled by an operating variable (\emph{R}*, via rate and \ensuremath{\Delta}) rather than by a tabulated property value.


\par\noindent\rule{\linewidth}{0.4pt}\par


\subsection*{S1.8 Numerical method and discretization}


The update is explicit and first-order: each cycle's state is computed from the previous cycle's state, with no iteration and no implicit solve. There is no time integration inside a cycle, the cycle is the time step. The scheme is therefore an Euler map on a discrete index, and its accuracy question is not \textquotedblleft{}is \ensuremath{\Delta}\emph{n} small enough\textquotedblright{} but \textquotedblleft{}does the discrete \ensuremath{\surd}\emph{n} approximation match the continuous \ensuremath{\surd}\emph{t} law\textquotedblright{}.


The only place this matters is the interphase term \emph{a}\textsubscript{SEI}/\ensuremath{\surd}\emph{n}, which is the discrete counterpart of d\emph{R}/d\emph{t} \ensuremath{\propto} \emph{t}\textsuperscript{\ensuremath{-}1/2}. The exact continuous solution gives, for the continuous case, 2\emph{a}\ensuremath{\surd}\emph{n} by integrating a/\ensuremath{\surd}\emph{i} from 0 to \emph{n}; the discrete sum \ensuremath{\Sigma}\textsubscript{i=1}\textsuperscript{n} \emph{a}/\ensuremath{\surd}\emph{i} is systematically smaller:


\begingroup\small
\setlength{\tabcolsep}{3pt}
\renewcommand{\arraystretch}{1.12}
\sloppy
\begin{tabular}{@{}>{\raggedright\arraybackslash}p{0.150\textwidth}>{\raggedright\arraybackslash}p{0.291\textwidth}>{\raggedright\arraybackslash}p{0.258\textwidth}>{\raggedright\arraybackslash}p{0.205\textwidth}@{}}
\hline
\textbf{\emph{n}} & \textbf{\ensuremath{\Sigma} \emph{a}/\ensuremath{\surd}\emph{i}} & \textbf{2\emph{a}\ensuremath{\surd}\emph{n}} & \textbf{ratio} \\ \hline
100 & 8.3653 & 9.0000 & 0.9295 \\
\hline
400 & 17.3541 & 18.0000 & 0.9641 \\
\hline
1000 & 27.8105 & 28.4605 & 0.9772 \\
\hline
\end{tabular}
\endgroup



The discrete scheme underestimates cumulative interphase growth by 7 \% at 100 cycles, 3.6 \% at 400, and 2.3 \% at 1000. This is a systematic bias in the direction that makes the knee slightly later than the continuous model would put it, i.e. it is conservative with respect to the claim (a later knee is harder to attribute to fast fade). It is not corrected: correcting it would change all quoted cycle numbers, and a 3 \% bias is far inside the 2.21\ensuremath{\times} range of knee positions that S1.7 spans.


\textbf{Stopping conditions.} \texttt{n\_max = 6000}. The coupled baseline stops itself at cycle 1044 by reaching LLI = \emph{Q}\ensuremath{_{0}}. The severed control runs to \emph{n}\textsubscript{max.} The model is fully deterministic: there is no randomness anywhere, so every quoted number is exactly reproducible.


\par\noindent\rule{\linewidth}{0.4pt}\par


\subsection*{S1.9 Full operating-condition map}


If the knee is a design variable rather than a fixed property value, then it must move when the operating condition moves, and it must move in a predictable direction: faster cycling and faster interphase growth should both bring the knee forward. The full map tests that prediction across the operating envelope, and it also marks the region in which the model stops producing a knee at all.


Knee cycle count as a function of C-rate (rows) and the SEI-growth factor \emph{T}\textsubscript{factor} (columns), the multiplier applied to the interphase growth rate:


\begingroup\footnotesize
\setlength{\tabcolsep}{3pt}
\renewcommand{\arraystretch}{1.12}
\sloppy
\begin{tabular}{@{}>{\raggedright\arraybackslash}p{0.326\textwidth}>{\raggedright\arraybackslash}p{0.126\textwidth}>{\raggedright\arraybackslash}p{0.107\textwidth}>{\raggedright\arraybackslash}p{0.107\textwidth}>{\raggedright\arraybackslash}p{0.107\textwidth}>{\raggedright\arraybackslash}p{0.107\textwidth}@{}}
\hline
\textbf{C-rate \ensuremath{\downarrow} / \emph{T}\textsubscript{factor} \ensuremath{\rightarrow}} & \textbf{0.25\ensuremath{\times}} & \textbf{0.5\ensuremath{\times}} & \textbf{1.0\ensuremath{\times}} & \textbf{2.0\ensuremath{\times}} & \textbf{4.0\ensuremath{\times}} \\ \hline
0.5C & none & none & 2651 & 910 & 330 \\
\hline
1.0C & 3937 & 1315 & 530 & 249 & 126 \\
\hline
2.0C & 268 & 157 & 104 & 72 & 49 \\
\hline
3.0C & 55 & 51 & 44 & 37 & 29 \\
\hline
\end{tabular}
\endgroup



Monotonic in both directions wherever a knee exists: higher rate \ensuremath{\rightarrow} earlier knee at every \emph{T}\textsubscript{factor}; faster SEI growth \ensuremath{\rightarrow} earlier knee at every rate. The rate effect is steep: over the S1.4 rate table, 0.5C to 2.5C moves the knee from 2651 to 64, a factor of 41.4, and the temperature effect is shallower (0.5\ensuremath{\times} to 2\ensuremath{\times} \emph{T}\textsubscript{factor} moves it 1315 \ensuremath{\rightarrow} 249, a factor of 5.3 at 1C). The rate-dependence factor is therefore 41.4 (\ensuremath{\approx}40).


Two \textquotedblleft{}none\textquotedblright{} entries at 0.5C \ensuremath{\times} \{0.25\ensuremath{\times}, 0.5\ensuremath{\times}\} deserve a note rather than a silent blank: at those combinations interphase growth is slow enough that \emph{R} does not reach \emph{R}* = 40 \ensuremath{\Omega}\ensuremath{\cdot}cm\ensuremath{^{2}} within 6000 cycles, so the loop never engages and there is no knee to find. That is the model behaving as designed --- no feedback, no knee, and it is the same statement as the ablation, obtained a different way.


\textbf{And where it breaks:} at 3.0C the knee arrives between cycle 29 and 55, at or below the fitting window used by the departure detector (30--200). The model is still internally consistent there, but the departure diagnostic is not defined, so the 2.0C row is marked invalid rather than reported. The S1.5 \textquotedblleft{}none\textquotedblright{} entries are the two 0.5C columns at 0.25\ensuremath{\times} and 0.5\ensuremath{\times}, where no knee forms.


\par\noindent\rule{\linewidth}{0.4pt}\par


\subsection*{S1.10 Claims the model does not support}


A model of this kind is easy to over-read, so the boundary is stated explicitly. Each item below is a claim that this model does not support, and the list is intended to be complete.


\begin{enumerate}
  \item It does not predict any real cell's cycle life. None of 530, 1044, 285, or any other number is a prediction. They are outputs of an unfitted model with placeholder parameters.
  \item It does not validate the coupling mechanism. It shows that if the Loop-A coupling exists with the assumed form, a knee follows and severing the coupling removes it. It is silent on whether the coupling exists in a real cell. That question is answered by measurement, not by this model; S4's branch B7 sets out the measurement.
  \item It does not model Loops B, C, or D. No dissolution, no crossover, no fracture, no electrolyte inventory. The rest of Figure 1 panel (a) --- every mechanism other than M1 and the M2 feedback path, and every edge not on the Loop A circuit --- is qualitative, and the model gives it no quantitative support. Loop D is absent from the model altogether, not merely unvalidated by it.
  \item It has no spatial dimension. No porosity, no electrode thickness, no potential distribution. Consequently it is not suited to Loop D (electrolyte depletion), which is intrinsically a volumetric argument. The Loop D budget in S2.6 is a separate, independent calculation with no model input.
  \item It does not model loss of active material or kinetic polarization loss. The model resolves one of the three capacity-loss channels, lithium inventory, and is silent on the other two, active material and kinetics. It therefore says nothing about how the three partition in a real cell.
  \item The knee definition is arbitrary and the spread is a factor of 1.9 (S1.6). Quoting a single knee cycle number without that band is over-precision.
  \item The threshold is an electrode-average statement (S1.4), not a local one. It cannot resolve patchiness, and it does not predict the first plating event.
  \item \textbf{No parameter is fitted, and \emph{b}\textsubscript{fresh} is not measured.} Its value 30 was chosen to give a physically plausible 3 nm re-formed interphase and a pre-knee exponent near 1. The 400\ensuremath{\times} sweep indicates that this choice is unlikely to change the qualitative conclusions, though it does change every cycle number.
  \item The model is not a \textquotedblleft{}bifurcation\textquotedblright{} in the strict dynamical-systems sense. It has a threshold and a saturating nonlinearity, which produces a sharp but finite transition. The Perspective accompanying these derivations describes the knee as behaving as a bifurcation; in the strict sense that word is a misnomer, and the mechanism here is a threshold with a saturating nonlinearity.
  \item Temperature enters only as a scalar multiplier on SEI growth. Real temperature dependence differs between the SEI-growth, charge-transfer, and transport terms, and the model does not represent that.
\end{enumerate}


\subsection*{S1.11 Tables S2--S5}


\subsubsection*{Table S2. Parameters of the reduced Loop-A model.}


All values are for 1C, \emph{Q}\ensuremath{_{0}} = 3.0 mAh cm\ensuremath{^{-}}\ensuremath{^{2}} unless stated. \textquotedblleft{}Provenance\textquotedblright{} distinguishes measured values (traceable to a cited source) from illustrative placeholders; no parameter is fitted to a fade curve.


\begingroup\footnotesize
\setlength{\tabcolsep}{3pt}
\renewcommand{\arraystretch}{1.12}
\sloppy
\begin{longtable}{@{}>{\raggedright\arraybackslash}p{0.108\textwidth}>{\raggedright\arraybackslash}p{0.069\textwidth}>{\raggedright\arraybackslash}p{0.143\textwidth}>{\raggedright\arraybackslash}p{0.226\textwidth}>{\raggedright\arraybackslash}p{0.347\textwidth}@{}}
\hline
\textbf{Symbol} & \textbf{Value} & \textbf{Unit} & \textbf{Meaning} & \textbf{Provenance} \\ \hline
\endhead
\emph{Q}\ensuremath{_{0}} & 3.0 & mAh cm\ensuremath{^{-}}\ensuremath{^{2}} & nominal areal capacity, 1C reference & Illustrative (typical high-energy Li-ion) \\
\emph{\ensuremath{\eta}}\textsubscript{ct} & 0.030 & V & charge-transfer overpotential at the operating rate & Illustrative, order-of-magnitude \\
\ensuremath{\Delta} & 0.090 & V & margin between anode operating potential and 0 V vs Li/Li\ensuremath{^{+}} & Illustrative; sets \emph{R}* \\
\emph{R}\ensuremath{_{0}} & 8.0 & \ensuremath{\Omega} cm\ensuremath{^{2}} & fresh-cell anode area-specific interphase resistance & Illustrative, order-of-magnitude \\
\emph{a}\textsubscript{SEI} & 0.45 & \ensuremath{\Omega} cm\ensuremath{^{2}} cycle\textsuperscript{\ensuremath{-}1/2} & diffusion-limited interphase growth prefactor & Anchored to \ensuremath{\surd}\emph{t} transport law; magnitude illustrative \\
\emph{b}\textsubscript{fresh} & 30.0 & \ensuremath{\Omega} cm\ensuremath{^{2}} per mAh cm\ensuremath{^{-}}\ensuremath{^{2}} & resistance added per unit of dead Li \ensuremath{\rightarrow} fresh interphase (loop gain) & No literature anchor. Reference value chosen for 3 nm implied thickness (\ensuremath{\rho} = 10\ensuremath{^{8}} \ensuremath{\Omega} cm) and pre-knee exponent \ensuremath{\approx} 1. Sweep in Table S4. \\
\emph{R}\textsubscript{cap} & 120.0 & \ensuremath{\Omega} cm\ensuremath{^{2}} & pore-filling saturation asymptote for interphase resistance & Illustrative (geometry of anode pore volume) \\
\emph{k}\textsubscript{pl} & 12.0 & V\ensuremath{^{-}}\ensuremath{^{1}} & sharpness of the plating turn-on & Illustrative (width of local-\emph{R} distribution) \\
\emph{f}\textsubscript{max} & 0.02 & --- & maximum fraction of applied charge that can plate & Illustrative upper bound enforcing mass balance \\
\emph{f}\textsubscript{dead} & 0.094 & --- & fraction of plated Li becoming electrically isolated & Measured, no-FEC case: \texttt{10.1016/j.}\hspace{0pt}\texttt{xcrp.2020.}\hspace{0pt}\texttt{100139} [FROM LITERATURE]. FEC case 0.033 \ensuremath{\pm} 0.006. \\
\emph{c}\textsubscript{LLI} & 0.010 & mAh cm\ensuremath{^{-}}\ensuremath{^{2}} per \ensuremath{\Omega} cm\ensuremath{^{2}} & LLI charged per unit of interphase resistance grown & Illustrative; converts \ensuremath{\Delta}\emph{R} into lost inventory \\
\emph{n}\textsubscript{max} & 6000 & cycles & integration limit & Numeric choice; only reached by the severed control \\
\hline
\end{longtable}
\endgroup



Derived quantities at 1C: \emph{j} = 3.000 mA cm\ensuremath{^{-}}\ensuremath{^{2}}; \emph{R}* = 20.00 \ensuremath{\Omega} cm\ensuremath{^{2}}; headroom \ensuremath{\Delta} \ensuremath{-} \emph{\ensuremath{\eta}}\textsubscript{ct} = 60 mV; threshold \emph{R}*/\emph{R}\ensuremath{_{0}} = 2.50.


\emph{Notes.} [a] The \ensuremath{\surd}\emph{n} form of the \emph{a}\textsubscript{SEI} term is the discrete counterpart of the \ensuremath{\surd}\emph{t} law for diffusion-limited interphase growth; S1.8 quantifies the discretization bias. [b] \emph{b}\textsubscript{fresh} is the only parameter that both (i) lacks an external constraint and (ii) strongly controls the knee position; it is the model's principal vulnerability and is disclosed as such. [c] \emph{f}\textsubscript{dead} is the only parameter obtained from a direct measurement, and it comes from a Li-metal cell; its use in a graphite-based Loop A is an assumption, labeled \texttt{[HYPOTHESIS]}.


\subsubsection*{Table S3. Baseline landmarks and the ablation, 1C, reference parameters.}


\textbf{(A) Baseline trajectory.}


\begingroup\small
\setlength{\tabcolsep}{3pt}
\renewcommand{\arraystretch}{1.12}
\sloppy
\begin{longtable}{@{}>{\raggedright\arraybackslash}p{0.350\textwidth}>{\raggedright\arraybackslash}p{0.347\textwidth}>{\raggedright\arraybackslash}p{0.219\textwidth}@{}}
\hline
\textbf{Quantity} & \textbf{Coupled} & \textbf{Severed} \\ \hline
\endhead
Early power-law fit, cycles 30--200 & \emph{LLI} = 7.268 \ensuremath{\times} 10\ensuremath{^{-}}\ensuremath{^{3}} \ensuremath{\cdot} \emph{n}\textsuperscript{0.522} & identical \\
First plating (cycle) & 220 & 220 \\
Departure from early power law (cycle) & 285 & none to 6000 \\
Knee, LLI \ensuremath{\geq} 20 \% \emph{Q}\ensuremath{_{0}} (cycle) & 530 & none \\
End of life, LLI = \emph{Q}\ensuremath{_{0}} (cycle) & 1044 & not reached \\
\emph{R}(0) (\ensuremath{\Omega} cm\ensuremath{^{2}}) & 8.00 & 8.00 \\
\emph{R}(400) (\ensuremath{\Omega} cm\ensuremath{^{2}}) & 26.85 & 24.16 \\
\emph{R}\textsubscript{max} (\ensuremath{\Omega} cm\ensuremath{^{2}}) & 75.12 & 60.51 \\
Pre-knee exponent (10--90 \% of knee) & 0.910 & --- \\
Post-knee exponent (knee \ensuremath{\rightarrow} 1.8\ensuremath{\times} knee) & 2.434 & --- \\
Retained capacity at cycle 1044 & 0 \% & 91.6 \% \\
Terminal LLI (mAh cm\ensuremath{^{-}}\ensuremath{^{2}}) & 3.00 (clipped) & 0.525 (at 6000) \\
\hline
\end{longtable}
\endgroup



\textbf{(B) Which term carries the loop.}


\begingroup\footnotesize
\setlength{\tabcolsep}{3pt}
\renewcommand{\arraystretch}{1.12}
\sloppy
\begin{tabular}{@{}>{\raggedright\arraybackslash}p{0.326\textwidth}>{\raggedright\arraybackslash}p{0.203\textwidth}>{\raggedright\arraybackslash}p{0.176\textwidth}>{\raggedright\arraybackslash}p{0.094\textwidth}>{\raggedright\arraybackslash}p{0.094\textwidth}@{}}
\hline
\textbf{Variant} & \textbf{Setting} & \textbf{Departure} & \textbf{Knee} & \textbf{EOL} \\ \hline
coupled & reference & 285 & 530 & 1044 \\
\hline
feed-forward arrow severed & \emph{b}\textsubscript{fresh} = 0 & 300 & 670 & 1681 \\
\hline
dead-lithium fraction severed & \emph{f}\textsubscript{dead} = 0 & none & none & 6000 \\
\hline
fully severed & both zero & none & none & 6000 \\
\hline
\end{tabular}
\endgroup



Setting \emph{f}\textsubscript{dead} = 0 forces \emph{d}\textsubscript{dead} = 0 in (S1.E5) and therefore removes both the resistance term \emph{b}\textsubscript{fresh}\ensuremath{\cdot}\emph{d}\textsubscript{dead} in (S1.E6) and the inventory term in (S1.E8), so that row is identical to the fully severed case. This is why the dead-lithium fraction, not the feed-forward arrow, is what makes the coupling a loop.


(C) Secondary knee criteria on the coupled baseline.


\begingroup\small
\setlength{\tabcolsep}{3pt}
\renewcommand{\arraystretch}{1.12}
\sloppy
\begin{tabular}{@{}>{\raggedright\arraybackslash}p{0.606\textwidth}>{\raggedright\arraybackslash}p{0.322\textwidth}@{}}
\hline
\textbf{Criterion} & \textbf{Knee (cycle)} \\ \hline
Cumulative LLI \ensuremath{\geq} 20 \% \emph{Q}\ensuremath{_{0}} & 530 \\
\hline
Max d\ensuremath{^{2}}(LLI)/d\emph{n}\ensuremath{^{2}} & 470 \\
\hline
Max retention curvature & 470 \\
\hline
Kneedle chord method & 479 \\
\hline
Local-exponent acceleration & 278 \\
\hline
\end{tabular}
\endgroup



\emph{Notes.} [a] The departure test fits \emph{A}\ensuremath{\cdot}\emph{n}\textsuperscript{p} over cycles 30--200 (pre-loop: first plating is at 220) and flags the first cycle after 200 with >10 \% excess over the extrapolation. [b] \textquotedblleft{}None\textquotedblright{} means the criterion is not met within \emph{n}\textsubscript{max} = 6000; it is a genuine null, not a failure to converge. [c] The 20 \%-LLI convention is the most conservative of the five criteria; the spread (278--530, a factor of 1.9) should be quoted alongside any single number.


\subsubsection*{Table S4. Sensitivity of the knee to the unanchored loop gain \emph{b}\textsubscript{fresh.}}


Sweep over a 400\ensuremath{\times} range (2.6 decades). Implied interphase thickness assumes \ensuremath{\rho} = 10\ensuremath{^{8}} \ensuremath{\Omega} cm.


\begingroup\scriptsize
\setlength{\tabcolsep}{3pt}
\renewcommand{\arraystretch}{1.12}
\sloppy
\begin{tabular}{@{}>{\raggedright\arraybackslash}p{0.122\textwidth}>{\raggedright\arraybackslash}p{0.145\textwidth}>{\raggedright\arraybackslash}p{0.067\textwidth}>{\raggedright\arraybackslash}p{0.122\textwidth}>{\raggedright\arraybackslash}p{0.067\textwidth}>{\raggedright\arraybackslash}p{0.111\textwidth}>{\raggedright\arraybackslash}p{0.122\textwidth}>{\raggedright\arraybackslash}p{0.100\textwidth}@{}}
\hline
\textbf{\emph{b}\textsubscript{fresh}} & \textbf{Implied thickness} & \textbf{Knee} & \textbf{Departure} & \textbf{EOL} & \textbf{Pre-exp.} & \textbf{Post-exp.} & \textbf{\emph{R}\textsubscript{end}} \\ \hline
1.0 & 0.1 nm & 663 & 299 & 1630 & 0.96 & 1.83 & 40.8 \\
\hline
3.0 & 0.3 nm & 650 & 298 & 1542 & 0.96 & 1.90 & 43.9 \\
\hline
10.0 & 1.0 nm & 610 & 294 & 1328 & 0.95 & 2.10 & 53.7 \\
\hline
30.0 & 3.0 nm & 530 & 285 & 1044 & 0.91 & 2.43 & 75.1 \\
\hline
60.0 & 6.0 nm & 462 & 275 & 873 & 0.86 & 2.56 & 94.8 \\
\hline
100.0 & 10.0 nm & 410 & 267 & 768 & 0.81 & 2.54 & 109.4 \\
\hline
200.0 & 20.0 nm & 348 & 254 & 667 & 0.71 & 2.37 & 123.5 \\
\hline
400.0 & 40.0 nm & 300 & 244 & 613 & 0.61 & 2.09 & 127.8 \\
\hline
\end{tabular}
\endgroup



\emph{Notes.} [a] Bold row is the reference value. [b] Knees exist at 8/8 values (C1); severing abolishes the knee at 5/5 values tested (C2, \emph{b}\textsubscript{fresh} = 10--200, a narrower range than C1). [c] Knee range 300--663 is a factor of 2.21, i.e. existence is robust while position is not. [d] Units of \emph{b}\textsubscript{fresh}: \ensuremath{\Omega} cm\ensuremath{^{2}} per mAh cm\ensuremath{^{-}}\ensuremath{^{2}}; \emph{R}\textsubscript{end} in \ensuremath{\Omega} cm\ensuremath{^{2}}. [e] Implied thickness = \emph{b}\textsubscript{fresh}/\ensuremath{\rho}, converted to nm.


\subsubsection*{Table S5. Knee cycle count vs C-rate and SEI-growth factor (\emph{T}\textsubscript{factor}).}


\begingroup\footnotesize
\setlength{\tabcolsep}{3pt}
\renewcommand{\arraystretch}{1.12}
\sloppy
\begin{tabular}{@{}>{\raggedright\arraybackslash}p{0.326\textwidth}>{\raggedright\arraybackslash}p{0.126\textwidth}>{\raggedright\arraybackslash}p{0.107\textwidth}>{\raggedright\arraybackslash}p{0.107\textwidth}>{\raggedright\arraybackslash}p{0.107\textwidth}>{\raggedright\arraybackslash}p{0.107\textwidth}@{}}
\hline
\textbf{C-rate \ensuremath{\downarrow} / \emph{T}\textsubscript{factor} \ensuremath{\rightarrow}} & \textbf{0.25\ensuremath{\times}} & \textbf{0.5\ensuremath{\times}} & \textbf{1.0\ensuremath{\times}} & \textbf{2.0\ensuremath{\times}} & \textbf{4.0\ensuremath{\times}} \\ \hline
0.5C & none & none & 2651 & 910 & 330 \\
\hline
1.0C & 3937 & 1315 & 530 & 249 & 126 \\
\hline
2.0C & 268 & 157 & 104 & 72 & 49 \\
\hline
3.0C & 55 & 51 & 44 & 37 & 29 \\
\hline
\end{tabular}
\endgroup



\emph{Notes.} [a] \textquotedblleft{}none\textquotedblright{} = \emph{R} does not reach \emph{R}* within 6000 cycles, so no plating and no knee; this is the ablation's conclusion obtained independently. [b] \emph{T}\textsubscript{factor} scales the SEI-growth prefactor; 2\ensuremath{\times} \ensuremath{\approx} +20 K. [c] At 3C the knee precedes the departure-fit window (cycles 30--200), so the departure diagnostic is not defined there even though the model remains internally consistent. [d] Knee cycle count at 1C, \emph{T}\textsubscript{factor} = 1 is 530; over the S1.4 rate table, the 0.5C\ensuremath{\rightarrow}2.5C span is 2651/64 = 41.4.


\par\noindent\rule{\linewidth}{0.4pt}\par


\subsection*{S1.12 Summary of the model}


The reduced model is a two-state, cycle-resolved scaling argument, not a fit: anode interphase resistance \emph{R} grows as \ensuremath{\surd}\emph{n} plus a term proportional to freshly created dead lithium, and lithium inventory loss \emph{LLI} accumulates from interphase growth and from that dead lithium. Plating ignites when the ohmic drop carries the anode past 0 V vs Li/Li\ensuremath{^{+}}, i.e. when \emph{R} > \emph{R}* = (\ensuremath{\Delta} \ensuremath{-} \emph{\ensuremath{\eta}}\textsubscript{ct})/\emph{j} --- a threshold that depends on rate and operating window rather than on any tabulated property value. Two self-limiting terms (mass balance, \emph{f} \ensuremath{\leq} \emph{f}\textsubscript{max}, and pore-filling saturation, \emph{R} \ensuremath{\rightarrow} \emph{R}\ensuremath{_{0}} + \emph{R}\textsubscript{cap}) keep every state variable physical and give the transition a finite post-knee slope.


The ablation separates the two couplings. Removing the dead-lithium fraction (\emph{f}\textsubscript{dead} = 0) abolishes the knee outright, leaving a cell that fades at the same early rate but never departs from its early power law out to 6000 cycles, because \emph{d}\textsubscript{dead} = 0 removes both the resistance term and the stranded-inventory term at once. Removing only the feed-forward arrow (\emph{b}\textsubscript{fresh} = 0) leaves dead lithium forming and stranding inventory, and the knee is delayed from 530 to 670 cycles rather than removed. The knee is therefore a structural consequence of the coupling rather than a kinetic one, and the parameter that makes the coupling a loop is the dead-lithium fraction, not the loop gain.


\subsection*{S1.13 The framework applied to published fade curves}


The Perspective states that the framework is "testable on any published fade curve that has been decomposed", and never performs the test. This section performs the part of it that does not require refitting anyone's published data.


\textbf{What can be compared, and what cannot.} S1.6 warns that no primary literature cycle number should be compared numerically with the model's, because the model is unfitted and its loop gain is a placeholder. That warning rules out reading \emph{530} against any published cycle count. It does not rule out comparing a \textbf{dimensionless ratio}, because a common scale factor cancels. The one ratio the framework makes a prediction about is the distance from the knee to the end of life: the framework is worth having only if there is life left after the knee. Here the model is unhelpful in a way worth stating plainly. Its primary knee convention is LLI \ensuremath{\geq} 0.20 Q\ensuremath{_{0}}, which is 80 \% retention by construction, and 80 \% retention is also the published definition of end of life. Under the primary convention used here the model therefore has zero interval: its knee is its end of life. The meaningful interval only appears when the knee is located by the onset criteria in S1.6, which place it at 278--285 cycles and leave 530 - 285 \ensuremath{\approx} 245 cycles, a ratio of 1.86. So the quantity compared below is definition-dependent from the model side as well as the data side, and that is the substance of the result.


\textbf{Data.} Knee point and end of life for 303 cells from 17 independent published studies, tabulated from the original figures by Attia et al. (J. Electrochem. Soc. 169, 060517, 2022; doi:10.1149/1945-7111/ac6d13) and taken from their public repository (github.com/tinosulzer/kneepoint-review). End of life is 80 \% of nominal capacity throughout, which makes it the one quantity defined identically across all 17 studies. The values are the original authors' readings of published figures, not independent measurements, and no cell was re-analyzed here.


The comparison is definition-limited before it is model-limited. 83 of the 303 cells (27 \%) have a knee point after their end of life. This is not an error in the compilation: the 17 studies detect the knee by different algorithms on different fades, and on a curve that has already passed 80 \% retention the detected inflection can fall later than the 80 \% crossing. The dispersion is therefore a property of the literature the framework is meant to organize. Any claim that the knee is a single well-defined object has to survive this, and the framework currently does not address it.


\textbf{The ratio.} On the 220 cells with knee \ensuremath{\leq} end of life, the published relation is close to linear:


\begingroup\small
\setlength{\tabcolsep}{3pt}
\renewcommand{\arraystretch}{1.12}
\sloppy
\begin{tabular}{@{}>{\raggedright\arraybackslash}p{0.411\textwidth}>{\raggedright\arraybackslash}p{0.517\textwidth}@{}}
\hline
\textbf{quantity} & \textbf{value} \\ \hline
published median EOL / knee & 1.26 \\
\hline
published interquartile range & 1.14--1.31 \\
\hline
least-squares fit & EOL = 1.20 \ensuremath{\times} knee + 22 (R\ensuremath{^{2}} = 0.936, n = 220) \\
\hline
per-study medians & 1.03--1.46 (16 of 17 studies), 2.62 (petzl, n = 3) \\
\hline
\end{tabular}
\endgroup



The model's own ratio depends entirely on which of the five knee definitions in S1.6 is used:


\begingroup\small
\setlength{\tabcolsep}{3pt}
\renewcommand{\arraystretch}{1.12}
\sloppy
\begin{tabular}{@{}>{\raggedright\arraybackslash}p{0.310\textwidth}>{\raggedright\arraybackslash}p{0.179\textwidth}>{\raggedright\arraybackslash}p{0.258\textwidth}>{\raggedright\arraybackslash}p{0.157\textwidth}@{}}
\hline
\textbf{knee definition} & \textbf{model knee (cycle)} & \textbf{model EOL, 80 \% retention (cycle)} & \textbf{ratio} \\ \hline
departure / local exponent & 278--285 & 530 & 1.86--1.91 \\
\hline
second derivative / Kneedle & 470--479 & 530 & 1.11--1.13 \\
\hline
LLI \ensuremath{\geq} 0.20 \emph{Q}\ensuremath{_{0}} (manuscript's "knee") & 530 & 530 & 1.00 \\
\hline
\end{tabular}
\endgroup



\textbf{Table S6.} The framework's ratio against the published range.


\begingroup\small
\setlength{\tabcolsep}{3pt}
\renewcommand{\arraystretch}{1.12}
\sloppy
\begin{tabular}{@{}>{\raggedright\arraybackslash}p{0.483\textwidth}>{\raggedright\arraybackslash}p{0.446\textwidth}@{}}
\hline
\textbf{} & \textbf{value} \\ \hline
cells / studies in the compilation & 303 / 17 \\
\hline
cells with knee \ensuremath{\leq} end of life & 220 (73 \%) \\
\hline
cells with knee > end of life & 83 (27 \%) \\
\hline
published median EOL / knee & 1.26 \\
\hline
published interquartile range & 1.14--1.31 \\
\hline
published least-squares fit & EOL = 1.20 \ensuremath{\times} knee + 22 (R\ensuremath{^{2}} = 0.936) \\
\hline
model ratio, onset knee definitions & 1.86--1.91 \\
\hline
model ratio, late knee definitions & 1.00--1.13 \\
\hline
\end{tabular}
\endgroup



Result: the test is inconclusive, and this is reported rather than resolved. The model's ratio spans 1.00--1.91 depending on a definitional choice that the literature also does not make consistently. The published interquartile range, 1.14--1.31, falls inside that span. Depending on which definition is adopted the framework is either consistent with the data (using the late knee definitions) or sits above it (using the onset definitions). A test whose outcome is fixed by the definition of its own central object has no discriminating power, and presenting either branch as a confirmation would be the same error corrected elsewhere.


One thing the comparison does establish. The framework's claim that the knee is a useful \emph{early-warning} object requires a measurable interval between the knee and end of life. In the published data that interval is about a quarter of the knee value: the median cell reaches 80 \% retention at 1.26 \ensuremath{\times} its knee cycle, not at the \textasciitilde{}2 \ensuremath{\times} the model produces under its onset definitions. The interval exists, which supports the framing, but it is short in most published cells. Any use of the knee as a warning signal is therefore operating on roughly a quarter of the knee-to-knee interval, and a proposed application that assumes more should say where its larger interval comes from.


\textbf{What this does not establish.} It does not test whether the coupling explains the observed knees. That requires the decomposition the Perspective asks for --- LLI, LAM and KL resolved as functions of cycle for individual published cells --- and only then can the framework's diagnostic question ("which loop is engaged?") be put to data. The published knee/EOL tabulation carries no decomposition, so the most informative available test is deferred. This is the largest open item and is stated as such in the accompanying Perspective.


\subsection*{S1.14 A test that does not depend on where the knee is placed}


S1.13 failed to discriminate because its outcome was fixed by the choice among five knee definitions. The model makes a second, sharper prediction that does not require locating the knee at all: the shape of the fade. Before the knee, lithium inventory accumulates close to linearly with cycle number; after it, the feedback makes the accumulation accelerate. In the reference curve the local exponent of LLI growth is 0.91 before the knee and 2.43 after (S1.7). An exponent can be fitted to any published capacity curve over a chosen window, so the prediction can be tested without the definitional choice that defeated S1.13.


\textbf{Data.} Harris et al. (2017), as compiled by Attia et al. (J. Electrochem. Soc. 169, 060517, 2022; \texttt{code/data/}\hspace{0pt}\texttt{variability/}\hspace{0pt}\texttt{harris2017.}\hspace{0pt}\texttt{xlsx}, sheet \texttt{ZP2A}). Twenty-four cells cycled to 10 C. Of these, 22 end below 95 \% retention (every cell dips below 95 \% at some point), 21 reach 80 \%, and 19 meet the inclusion criterion below. The distinction between ending below a threshold and touching it matters here, because a single low cycle should not place a cell in the accelerating group. This is cell-level capacity against cycle number, which is what the test needs; it is not a decomposition.


\textbf{Method.} For each cell the retained-capacity trace is smoothed with a running median to suppress cycle-to-cycle scatter, converted to fractional loss (\texttt{1 - capacity}), and \texttt{log(loss)} is fitted against \texttt{log(cycle)} separately before and after the cycle at which retention crosses 80 \%. The exponent of loss is the quantity the model states; fitting \texttt{log(retention)} instead is meaningless because early retention sits at 0.9995, where the logarithm is flat, and a first version of this analysis that did so returned -0.05 for every cell. Cells whose fitted exponent exceeded 8 in magnitude were rejected as fits to residual scatter rather than to a trend, which removed 3 of 22.


\textbf{Result.} The direction of the prediction is supported in most cells and contradicted in a minority that the framework cannot accommodate:


\begingroup\small
\setlength{\tabcolsep}{3pt}
\renewcommand{\arraystretch}{1.12}
\sloppy
\begin{tabular}{@{}>{\raggedright\arraybackslash}p{0.615\textwidth}>{\raggedright\arraybackslash}p{0.313\textwidth}@{}}
\hline
\textbf{} & \textbf{value} \\ \hline
cells analyzed & 19 \\
\hline
median exponent before 80 \% retention & 1.04 \\
\hline
median exponent after & 1.43 \\
\hline
median change & +0.44 \\
\hline
cells where the exponent rises & 13 of 19 (68 \%) \\
\hline
cells where it is flat (within \ensuremath{\pm}0.1) & 2 \\
\hline
cells where it falls & 4 of 19 (21 \%) \\
\hline
\end{tabular}
\endgroup



The rising majority is consistent with the framework: fade accelerates after the knee, as a feedback loop requires. The four falling cells are the informative result. They do not scatter; they settle at about 0.79 retention and then lose less than 0.003 of retention over the next hundred cycles, which is an order of magnitude below the cycle-to-cycle variation of the cells that accelerate. A framework whose central mechanism is a loop that can only amplify has no route to a cell that decelerates after its knee. Those cells are evidence that something other than the Loop-A feedback dominates their later life, and the framework as written does not predict them.


\textbf{Magnitudes are not compared.} The measured medians, 1.04 and 1.43, are of the same order as the model's 0.91 and 2.43, but the model is unfitted and its loop gain is a placeholder (S1.6), so the numerical proximity is not agreement and is not claimed as such. Only the sign of the change is a fair test, and only the sign is used.


\textbf{What this still does not do.} It does not decompose capacity into LLI, LAM and KL. A rising loss exponent is consistent with a feedback loop and equally consistent with any other accelerating mechanism, including ones already in the literature. The test is therefore a necessary condition and not evidence that the coupling is the cause. Performing the decomposition would require low-rate OCV traces at several ages for individual cells with known half-cell curves; the public compilations that carry knee and end-of-life do not include them, and the open datasets examined here (Harris 2017; Ecker 2014; the Battery Life raw archive, whose cells span only 100 cycles and reach 98 \% retention) do not provide them either. The diagnostic question the Perspective actually poses, which loop is engaged, therefore remains untested on published data, and this section is reported as a partial test rather than as the demonstration the Perspective promises.


\subsection*{S1.15 The Section 2 prescription run on a published dataset}


Section 5 of the Perspective applies the accounting prescription to one public dataset and reports both what it recovers and where it fails. This note gives the method and the full numbers.


\textbf{Data.} Harris et al., \emph{Journal of Power Sources} 342, 589--597 (2017); the file \texttt{raw/harris2017}\hspace{0pt}\texttt{.xlsx} in \texttt{research/13-}\hspace{0pt}\texttt{fade-}\hspace{0pt}\texttt{application}, sheet \texttt{ZP2A}, columns are normalized capacity retention per cycle for 24 commercial pouch cells cycled at 10 C to 100 \% depth of discharge over 593 cycles. This is retention only. It carries no LLI/LAM/KL decomposition, so the analysis tests the shape the framework predicts, not the channel attribution it would prefer. That limitation is stated in the Perspective and is not repaired here.


\textbf{Method.} The per-cycle loss is taken as the negative finite difference of retention, \emph{f}(\emph{N}) = \ensuremath{-}[\emph{Q}(\emph{N}) \ensuremath{-} \emph{Q}(\emph{N}\ensuremath{-}1)]/\emph{Q}\_(N\ensuremath{-}1). This is fitted against cycle number as \emph{f}(\emph{N}) = \emph{f}0 + \ensuremath{\alpha}\emph{N} by ordinary least squares over the cycles where retention lies in a stated window. \ensuremath{\alpha} is the coupling strength of Section 2: the rate at which efficiency loss grows. The window is a parameter, not a convenience, and the sensitivity to it is reported below rather than hidden. Script: \texttt{research/13-}\hspace{0pt}\texttt{fade-}\hspace{0pt}\texttt{application/}\hspace{0pt}\texttt{worked\_}\hspace{0pt}\texttt{example.py}; output \texttt{worked\_}\hspace{0pt}\texttt{example\_}\hspace{0pt}\texttt{results.json}.


\textbf{Result, primary window 90-100 \% retention.} 24 cells analyzed. \ensuremath{\alpha} > 0 in 22 of 24, median 5.57\ensuremath{\times}10\ensuremath{^{-}}\ensuremath{^{7}} per cycle, range \ensuremath{-}3.33\ensuremath{\times}10\ensuremath{^{-}}\ensuremath{^{6}} to 4.00\ensuremath{\times}10\ensuremath{^{-}}\ensuremath{^{6}}. The two negative cells are not a coherent group: one ends at 84.8 \% retention, the other at 3.5 \%, so in the second the fit window is short and the slope is not meaningful. Neither undermines the sign of the majority.


\textbf{A constant \emph{f} is not enough.} A single constant \emph{f} fitted to the whole trace leaves a median RMSE of 0.0420 in retention. A quadratic in \emph{N} leaves 0.0082, better in 24 of 24 cells, and its coefficient is negative in 24 of 24. Negative curvature in retention is the signature Section 2 predicts whenever \emph{f} itself rises with cycle number, so the curvature is present in every cell of this set.


\textbf{Three independent estimates.} (i) The finite-difference fit above. (ii) A global quadratic fit to retention, with an effective coupling taken as minus twice the quadratic coefficient: positive in 24 of 24 cells, median 1.18\ensuremath{\times}10\ensuremath{^{-}}\ensuremath{^{6}}. This estimate does not use \emph{f} at all. (iii) Comparing a linear fit to the first and second halves of the trace: the second-half slope is steeper in 24 of 24 cells. The three agree on sign in every case, so the finding does not depend on the fitting scheme.


\textbf{Window sensitivity, which is the honest limit.} The sign of \ensuremath{\alpha} is not a property of the cell alone:


\begingroup\small
\setlength{\tabcolsep}{3pt}
\renewcommand{\arraystretch}{1.12}
\sloppy
\begin{tabular}{@{}>{\raggedright\arraybackslash}p{0.324\textwidth}>{\raggedright\arraybackslash}p{0.137\textwidth}>{\raggedright\arraybackslash}p{0.110\textwidth}>{\raggedright\arraybackslash}p{0.333\textwidth}@{}}
\hline
\textbf{Window (retention)} & \textbf{cells} & \textbf{\ensuremath{\alpha} > 0} & \textbf{median \ensuremath{\alpha} (per cycle)} \\ \hline
95-100 \% & 24 & 6 & \ensuremath{-}1.43\ensuremath{\times}10\ensuremath{^{-}}\ensuremath{^{7}} \\
\hline
90-100 \% & 24 & 22 & 5.56\ensuremath{\times}10\ensuremath{^{-}}\ensuremath{^{7}} \\
\hline
85-100 \% & 24 & 23 & 8.34\ensuremath{\times}10\ensuremath{^{-}}\ensuremath{^{7}} \\
\hline
80-100 \% & 24 & 22 & 9.88\ensuremath{\times}10\ensuremath{^{-}}\ensuremath{^{7}} \\
\hline
\end{tabular}
\endgroup



Over the narrowest window the curvature is smaller than the cycle-to-cycle scatter of a 10 C protocol, so the sign inverts. This is not a defect of the dataset; it is what any efficiency-curvature measurement faces when the window is chosen after seeing the data. A cross-study comparison of \ensuremath{\alpha} is therefore only meaningful if the window is fixed in advance, which no protocol in the literature does. Window bounds belong in the reporting checklist of Section 8 alongside rate, temperature and stack pressure. They also belong in prediction P6 of Section 7, and the reason is worth stating plainly, because the first version of P6 did not have them and was therefore refuted by this very table. As worded then, P6 said \ensuremath{\alpha} is positive and is refuted if it changes sign while retention stays above 90 \%. Every cell in the 95-100 \% row satisfies that condition, and 18 of 24 have \ensuremath{\alpha} < 0 there, so the prediction was false on the data it cited as support. The amended P6 makes the window part of the prediction rather than an adjustable setting: concave retention over a pre-registered window no narrower than 90 \%, with the 95-100 \% inversion reported in the same paragraph as the limit. What P6 now claims is falsifiable and, on this dataset, not falsified.


\textbf{What this does and does not establish.} It establishes that the quantity Section 2 asks for is measurable on published data at no experimental cost, and that it is positive in nearly every cell of a well-known 24-cell set. It does not establish which channel drives the curvature, because retention alone cannot separate LLI from LAM. It also does not establish that \ensuremath{\alpha} is a materials constant. Among the 22 cells with \ensuremath{\alpha} > 0 the fitted value spans 8.1\ensuremath{\times}10\ensuremath{^{-}}\ensuremath{^{9}} to 4.0\ensuremath{\times}10\ensuremath{^{-}}\ensuremath{^{6}} per cycle, a factor of about 500, across cells of one chemistry from one supplier run under one protocol. At that spread \ensuremath{\alpha} behaves as a cell-level property at least as much as a chemistry-level one. That is consistent with the Perspective's framing, in which operating condition sets the rate at which a margin is consumed. It is not evidence against a role for chemistry: a factor of 500 between two cells of the same nominal chemistry says the chemistry label is too coarse to predict \ensuremath{\alpha}, not that chemistry is irrelevant. S1.16 shows that the most influential parameter in the model is in fact electrode chemistry, which is the sharper caveat.


\subsection*{S1.16 Direct sweep of the operating margin}


Section 5 of the Perspective states that the knee is a design variable because each loop closes on an operating margin. The reduced Loop-A model contains exactly one parameter that is that margin: \ensuremath{\Delta} in (S1.E2), the gap between the anode's equilibrium potential and 0 V vs Li/Li\ensuremath{^{+}}. The plating threshold follows from it as \emph{R}* = (\ensuremath{\Delta} \ensuremath{-} \ensuremath{\eta}\textsubscript{ct})/\emph{j}.


Every sweep reported so far in S1.7 and S1.9 varies something else. S1.7 varies \emph{b}\textsubscript{fresh}, the loop gain. S1.9 varies C-rate and the SEI-growth factor. Neither is the margin. C-rate is not a clean substitute for it: raising the rate lowers \emph{R}* through \emph{j}, but it also raises \emph{Q}\textsubscript{applied} in (S1.E4), so a C-rate sweep moves the plated charge per cycle at the same time. This note closes that gap. \ensuremath{\Delta} is varied at fixed C-rate (1 C), fixed \emph{Q}\ensuremath{_{0}} (3.0 mAh cm\ensuremath{^{-}}\ensuremath{^{2}}) and fixed \emph{b}\textsubscript{fresh} (30), so \emph{j} and \emph{Q}\textsubscript{applied} are constant and the margin is the only thing that moves. Script: \texttt{research/06-}\hspace{0pt}\texttt{proposal/\_}\hspace{0pt}\texttt{support/}\hspace{0pt}\texttt{margin\_sweep.}\hspace{0pt}\texttt{py}.


\textbf{Range, and why it is narrow.} The departure detector of S1.5 fits its pre-loop baseline over cycles 30-200. If plating begins before cycle 200, that baseline is no longer pre-loop and the resulting knee is not comparable with the reference case. Above \ensuremath{\Delta} = 0.090 V plating starts after cycle 200, so the clean window runs from 0.090 to 0.120 V, a factor of 1.33. Below 0.090 V the rows are reported but marked contaminated. This is a limit of the analysis, not a choice, and the narrow window weakens what the sweep can claim.


\begingroup\footnotesize
\setlength{\tabcolsep}{3pt}
\renewcommand{\arraystretch}{1.12}
\sloppy
\begin{tabular}{@{}>{\raggedright\arraybackslash}p{0.116\textwidth}>{\raggedright\arraybackslash}p{0.163\textwidth}>{\raggedright\arraybackslash}p{0.098\textwidth}>{\raggedright\arraybackslash}p{0.098\textwidth}>{\raggedright\arraybackslash}p{0.174\textwidth}>{\raggedright\arraybackslash}p{0.232\textwidth}@{}}
\hline
\textbf{\ensuremath{\Delta} (V)} & \textbf{\emph{R}* (\ensuremath{\Omega} cm\ensuremath{^{2}})} & \textbf{knee} & \textbf{EOL} & \textbf{plating at} & \textbf{baseline} \\ \hline
0.120 & 30.0 & 1134 & 1717 & 771 & clean \\
\hline
0.115 & 28.3 & 1007 & 1577 & 651 & clean \\
\hline
0.110 & 26.7 & 891 & 1449 & 543 & clean \\
\hline
0.105 & 25.0 & 786 & 1333 & 447 & clean \\
\hline
0.100 & 23.3 & 692 & 1227 & 361 & clean \\
\hline
0.095 & 21.7 & 606 & 1131 & 285 & clean \\
\hline
0.090 & 20.0 & 530 & 1044 & 220 & clean \\
\hline
0.085 & 18.3 & 463 & 967 & 164 & contaminated \\
\hline
0.080 & 16.7 & 404 & 898 & 116 & contaminated \\
\hline
0.070 & 13.3 & 308 & 785 & 47 & contaminated \\
\hline
0.060 & 10.0 & 239 & 699 & 10 & contaminated \\
\hline
\end{tabular}
\endgroup



\textbf{Result.} Across the clean window \ensuremath{\Delta} 0.120 \ensuremath{\rightarrow} 0.090 V, a factor of 1.33, the knee moves from 1134 to 530 cycles, a factor of 2.14. The local elasticity d(log knee)/d(log \ensuremath{\Delta}) is 2.5-2.8 and is stable across the window, so the response is not an artefact of one operating point. The structural reason is visible in the threshold: \emph{R}* is linear in \ensuremath{\Delta} with slope 1/\emph{j} = 333 \ensuremath{\Omega} cm\ensuremath{^{2}} V\ensuremath{^{-}}\ensuremath{^{1}}, so a 10 mV change moves the threshold by 3.3 \ensuremath{\Omega} cm\ensuremath{^{2}}, which at the reference \emph{R}* = 20 \ensuremath{\Omega} cm\ensuremath{^{2}} is a 17 \% change in the threshold itself.


Comparison with the gain sweep, on equal footing. The comparison that needs no fitting is the span itself: the loop gain moves the knee by a factor of 2.21 across a 400-fold range (S1.7), while the margin moves it by 2.14 across a 1.33-fold range. A 300-fold larger change in the gain buys about the same shift in the knee.


The elasticities, d(log knee)/d(log x), make the same point in local terms. For the margin it is 2.5-2.8 and rises slowly with \ensuremath{\Delta}. For the gain it is 0.13 over the whole range but reaches 0.17-0.23 locally between 10 and 400. Taking the local gain value, the margin is still roughly a dozen times more effective per e-fold; taking the overall gain value it is about twenty times. Either way the ordering is large and unambiguous, which is the quantitative form of the statement that the knee is set by the operating condition and not by the materials.


\textbf{What this does not settle.} Three things, and the first two cut against the framing of the Perspective rather than for it.


First, \ensuremath{\Delta} is a placeholder, not a measured value, so the sweep establishes the sensitivity of the knee to the margin and not its position in any real cell.


Second, and more awkwardly, \ensuremath{\Delta} is not a process knob. It is the gap between the anode's equilibrium potential and 0 V vs Li/Li\ensuremath{^{+}}, and that equilibrium potential is a property of the electrode material: about 0.08-0.10 V for the stage-1 plateau of graphite. Moving \ensuremath{\Delta} therefore means changing the anode chemistry, not turning a dial on the line. The other margins named in the Perspective are more tractable: C-rate, temperature, voltage window, stack pressure and electrolyte loading are all genuine operating variables. So the sweep shows that the most influential margin in the model is the least accessible one, which weakens rather than strengthens the practical claim. What survives is the conceptual point: the knee is set by a threshold crossing, and a threshold is not a materials constant in the sense the literature assumes.


Third, the model has one loop. Whether the same ordering of sensitivities survives when all four loops compete is not tested, and the Perspective states that ranking as a prediction for that reason.


\section*{S2. Derivations behind the reported quantities}


\subsection*{S2.0 Basis of the derivations}


A coupling only matters if its margin is of order unity. If a feedback arrow delivers a thousandth of the effect it is supposed to drive, it can be dropped and the mechanisms can be treated as independent. If it delivers a thousand times that effect, the loop saturates immediately and one mechanism dominates. The interesting case is the narrow one in between, and deciding which case each of the four loops falls into requires a number for every arrow. This section supplies those numbers.


The arithmetic is coarse by construction. Each derivation is an order-of-magnitude estimate, not a measurement and not a first-principles calculation, and each is stated to the precision its inputs support, usually one significant figure. Quoting "11.19 \ensuremath{\mu}m" would misrepresent the precision; "\textasciitilde{}11 \ensuremath{\mu}m" is the appropriate form. Several derivations carry a stated uncertainty of 2--5\ensuremath{\times}, and the argument is built to survive that uncertainty, because the conclusion is about orders of magnitude and not about a specific value.


No number here is a measured value; each is arithmetic on stated inputs. Second, where this section writes "\textasciitilde{}", the input is uncertain by at least a factor of 2, whereas three significant figures are given only where the input is a physical constant and the precision is real.


Status tags are used throughout: \texttt{[DERIVED]} = computed here, arithmetic shown; \texttt{[FROM LITERATU}\hspace{0pt}\texttt{RE: DOI]} = supplied by a published source; \texttt{[HYPOTHESIS]} = proposed, not established; \texttt{[ASSUMED]} = an order-of-magnitude input that is not traceable to a specific measurement.


\par\noindent\rule{\linewidth}{0.4pt}\par


\subsection*{S2.1 Loop A: the SEI ohmic drop}


Consider the anode of a cell that is about to plate lithium. Whether it plates depends on the anode potential: plating begins when that potential falls to 0 V versus Li/Li\ensuremath{^{+}}. The anode is driven there by the ohmic drop across the interphase, \emph{j\ensuremath{\cdot}R}, and it is held back by the operating margin \ensuremath{\Delta} and by the charge-transfer overpotential. So the interphase resistance controls plating onset, and the question for this derivation is how large that resistance is, and therefore whether Loop A can be active in a real cell.


The sign and the magnitude are answered separately. The sign is robust: a resistive interphase always pushes the anode toward plating, so the direction of the arrow requires no assumption. The magnitude is not robust at all: it spans four orders of magnitude, because the conductivity of the interphase is not known to better than that. The consequence is stated at the end of this section, and it constrains what Loop A can and cannot be used to predict.


\textbf{The derivation.} Treat the interphase as a slab of thickness \emph{d} and ionic conductivity \ensuremath{\sigma} carrying the applied current density \emph{j}. Its area-specific resistance is


\begin{equation}R_{\text{SEI}} = \frac{d}{\sigma}\quad[\Omega\,\text{cm}^{2}] \tag{S2.E1}\end{equation}


and the ohmic potential drop across it is


\begin{equation}\Delta\phi_{\text{SEI}} = j\,R_{\text{SEI}} = \frac{j\,d}{\sigma}\quad[\text{V}] \tag{S2.E2}\end{equation}


\texttt{[DERIVED]}: this is Ohm's law applied to a series slab; there is no modeling choice beyond treating the interphase as uniform.


\textbf{Evaluating it.} Take \emph{d} = 30 nm (a representative SEI thickness for a graphite anode after formation; the surveyed band is \textasciitilde{}10--200 nm, \texttt{[FROM LITERATU}\hspace{0pt}\texttt{RE: 10.1088/}\hspace{0pt}\texttt{1361-6528/}\hspace{0pt}\texttt{ac17fe]}. A round value is used because the spread, not the point, is the result) and \emph{j} = 3.0 mA cm\ensuremath{^{-}}\ensuremath{^{2}} = 3.0 \ensuremath{\times} 10\ensuremath{^{-}}\ensuremath{^{3}} A cm\ensuremath{^{-}}\ensuremath{^{2}} (1C for a 3 mAh cm\ensuremath{^{-}}\ensuremath{^{2}} electrode):


\begingroup\small
\setlength{\tabcolsep}{3pt}
\renewcommand{\arraystretch}{1.12}
\sloppy
\begin{tabular}{@{}>{\raggedright\arraybackslash}p{0.210\textwidth}>{\raggedright\arraybackslash}p{0.284\textwidth}>{\raggedright\arraybackslash}p{0.423\textwidth}@{}}
\hline
\textbf{\ensuremath{\sigma} (S cm\ensuremath{^{-}}\ensuremath{^{1}})} & \textbf{\emph{R}\textsubscript{SEI} (\ensuremath{\Omega} cm\ensuremath{^{2}})} & \textbf{\ensuremath{\Delta}\ensuremath{\phi}\textsubscript{SEI} = \emph{j}\ensuremath{\cdot}\emph{R}\textsubscript{SEI} (mV)} \\ \hline
10\ensuremath{^{-}}\ensuremath{^{7}} & 30 & 90 \\
\hline
10\ensuremath{^{-}}\ensuremath{^{8}} & 300 & 900 \\
\hline
10\ensuremath{^{-}}\ensuremath{^{9}} & 3000 & 9000 \\
\hline
\end{tabular}
\endgroup



\texttt{[DERIVED]}: all three rows. Arithmetic: at \ensuremath{\sigma} = 10\ensuremath{^{-}}\ensuremath{^{8}} S cm\ensuremath{^{-}}\ensuremath{^{1}}, \emph{R} = 30\ensuremath{\times}10\ensuremath{^{-}}\ensuremath{^{7}} cm / 10\ensuremath{^{-}}\ensuremath{^{8}} S cm\ensuremath{^{-}}\ensuremath{^{1}} = 300 \ensuremath{\Omega} cm\ensuremath{^{2}}, and \ensuremath{\Delta}\ensuremath{\phi} = 3\ensuremath{\times}10\ensuremath{^{-}}\ensuremath{^{3}} A cm\ensuremath{^{-}}\ensuremath{^{2}} \ensuremath{\times} 300 \ensuremath{\Omega} cm\ensuremath{^{2}} = 0.9 V.


Why the \ensuremath{\sigma} band is the point, not the value. Pure Li\ensuremath{_{2}}CO\ensuremath{_{3}} is an electronic insulator with ionic conductivity in the 10\ensuremath{^{-}}\ensuremath{^{9}}--10\ensuremath{^{-}}\ensuremath{^{7}} S cm\ensuremath{^{-}}\ensuremath{^{1}} range depending on phase and microstructure, and measured SEI conductivities span at least this band. The choice of \ensuremath{\sigma} inside this band changes the predicted SEI overpotential by a factor of 100, from 90 mV to 9 V. That is not a precision problem to be narrowed; it means the loop's driving term cannot be predicted from first principles, only measured per cell.


The cross-check that makes the band usable. A literal 0.9 V SEI drop (the \ensuremath{\sigma} = 10\ensuremath{^{-}}\ensuremath{^{8}} case) is incompatible with the total anode polarization of a working cell at 1C, 50--150 mV \texttt{[ASSUMED]}. A 0.9 V drop across the interphase alone would exceed the entire measured anode overpotential, which would mean the anode was already plating at cycle 1 in every cell ever measured. It is not. Therefore the effective \ensuremath{\sigma} in a functioning anode must be at or above the 10\ensuremath{^{-}}\ensuremath{^{7}} S cm\ensuremath{^{-}}\ensuremath{^{1}} end, where \ensuremath{\Delta}\ensuremath{\phi}\textsubscript{SEI} \ensuremath{\approx} 90 mV, comparable to but not exceeding the measured polarization. \texttt{[DERIVED]} (the incompatibility argument); \texttt{[HYPOTHESIS]} (that this reflects genuinely higher effective conductivity rather than, e.g., a thinner effective interphase or a partially non-ohmic conduction mechanism).


Reconciling this band with the reduced model. The model of S1 adopts a fresh-cell anode area-specific resistance of \emph{R}\ensuremath{_{0}} = 8.0 \ensuremath{\Omega} cm\ensuremath{^{2}}. At \emph{d} = 30 nm that corresponds to an effective conductivity of 3.8 \ensuremath{\times} 10\ensuremath{^{-}}\ensuremath{^{7}} S cm\ensuremath{^{-}}\ensuremath{^{1}}, marginally above the 10\ensuremath{^{-}}\ensuremath{^{7}} S cm\ensuremath{^{-}}\ensuremath{^{1}} upper end of the band above. The two are not in contradiction, but the difference must be stated rather than left implicit. The incompatibility argument rules out \emph{R}\textsubscript{SEI} \ensuremath{\geq} 300 \ensuremath{\Omega} cm\ensuremath{^{2}}, since such a cell would plate from cycle 1; it does not fix \emph{R}\textsubscript{SEI} below that. A functioning anode with \emph{R}\textsubscript{SEI} \ensuremath{\approx} 8 \ensuremath{\Omega} cm\ensuremath{^{2}} is equally consistent with a thinner effective film (\ensuremath{\approx} 8 nm at 10\ensuremath{^{-}}\ensuremath{^{7}} S cm\ensuremath{^{-}}\ensuremath{^{1}}) or with a more conductive one (\ensuremath{\approx} 4 \ensuremath{\times} 10\ensuremath{^{-}}\ensuremath{^{7}} S cm\ensuremath{^{-}}\ensuremath{^{1}} at 30 nm), and the measurements available do not separate the two. The reduced model therefore takes \emph{R}\ensuremath{_{0}} as a lumped parameter and does not commit to a thickness--conductivity pair. What it requires of S2.1 is only that \emph{R}\ensuremath{_{0}} sit below \emph{R}*, which the band permits.


What the result supports, and what it does not. It supports the sign and the scale: the interphase drop is of the same order as the margin between the anode operating potential and 0 V vs Li/Li\ensuremath{^{+}} (\ensuremath{\Delta} \ensuremath{\approx} 90 mV, from the \ensuremath{\Delta}\ensuremath{\phi}\textsubscript{SEI} = 90 mV case in S2.1), so the interphase resistance is a plausible control on plating onset, which is what Loop A requires. It does not support any prediction of when a given cell will plate from its SEI thickness alone (S1.10, items 1 and 7).


\textbf{Uncertainty.} \ensuremath{\sigma} is unknown to \ensuremath{\pm}2 orders of magnitude, so \ensuremath{\Delta}\ensuremath{\phi}\textsubscript{SEI} is unknown to \ensuremath{\pm}2 orders. The ordering (interphase drop comparable to the plating margin) is robust; the value is not. The quantity is therefore always reported as \textquotedblleft{}comparable to\textquotedblright{} and never as a single number without its range.


\par\noindent\rule{\linewidth}{0.4pt}\par


\subsection*{S2.2 The plating threshold \emph{R}* and the inverse problem}


Section S2.1 gave the interphase resistance as a range four orders of magnitude wide, which is useless as a design number. This section turns the relation around. Instead of asking how resistive a given interphase is, ask what resistance would be \emph{just enough} to start plating at a given current density. That inverted quantity is the plating threshold \emph{R}*, and unlike the resistance itself it depends only on operating variables.


\textbf{The derivation.} Plating requires the anode potential to fall below 0 V vs Li/Li\ensuremath{^{+}}. Starting from an operating anode potential \ensuremath{\Delta} above 0 V and subtracting the charge-transfer overpotential \emph{\ensuremath{\eta}}\textsubscript{ct} and the ohmic drop \emph{j\ensuremath{\cdot}R}, the anode reaches 0 V when


\begin{equation}j\,R + \eta_{\text{ct}} = \Delta \quad\Longrightarrow\quad R^{*} = \frac{\Delta - \eta_{\text{ct}}}{j} \tag{S2.E3}\end{equation}


\texttt{[DERIVED]}. Note the dimensional check: (V \ensuremath{-} V)/(A cm\ensuremath{^{-}}\ensuremath{^{2}}) = \ensuremath{\Omega} cm\ensuremath{^{2}}. At 1C with \ensuremath{\Delta} = 90 mV, \emph{\ensuremath{\eta}}\textsubscript{ct} = 30 mV, \emph{j} = 3.0 \ensuremath{\times} 10\ensuremath{^{-}}\ensuremath{^{3}} A cm\ensuremath{^{-}}\ensuremath{^{2}}: \emph{R}* = 0.060 V / 0.0030 A cm\ensuremath{^{-}}\ensuremath{^{2}} = 20.0 \ensuremath{\Omega} cm\ensuremath{^{2}}. The headroom is \ensuremath{\Delta} \ensuremath{-} \emph{\ensuremath{\eta}}\textsubscript{ct} = 60 mV.


The equation that does the most work in this analysis. Rearranged, it converts an observable --- a resistance at which plating begins --- into an \emph{operating variable}, and shows what it does not depend on:


\begin{equation}R^{*}(j) = \frac{60\ \text{mV}}{j} \tag{S2.E4}\end{equation}


so \emph{R}* falls as 1/\emph{j}: 40.0 \ensuremath{\Omega} cm\ensuremath{^{2}} at 0.5C, 20.0 at 1C, 13.3 at 1.5C, 10.0 at 2C, 8.0 at 2.5C. No tabulated materials constant appears, though \ensuremath{\Delta} itself is set by the anode material (S1.16). A cell plates when its interphase resistance exceeds a rate-dependent threshold, which is why onset is a design variable rather than a property value.


\textbf{The inverse problem: the interphase thickness corresponding to \emph{R}*.} Combining the resistance expression of S2.1 with the threshold gives the \emph{critical thickness}


\begin{equation}d^{*} = \sigma\,R^{*} = \frac{\sigma\,(\Delta - \eta_{\text{ct}})}{j} \tag{S2.E5}\end{equation}


Evaluating at 1C (\emph{R}* = 20 \ensuremath{\Omega} cm\ensuremath{^{2}}) across the same \ensuremath{\sigma} band:


\begingroup\small
\setlength{\tabcolsep}{3pt}
\renewcommand{\arraystretch}{1.12}
\sloppy
\begin{tabular}{@{}>{\raggedright\arraybackslash}p{0.155\textwidth}>{\raggedright\arraybackslash}p{0.219\textwidth}>{\raggedright\arraybackslash}p{0.543\textwidth}@{}}
\hline
\textbf{\ensuremath{\sigma} (S cm\ensuremath{^{-}}\ensuremath{^{1}})} & \textbf{\emph{d}* = \ensuremath{\sigma}\ensuremath{\cdot}\emph{R}* (nm)} & \textbf{Interpretation} \\ \hline
10\ensuremath{^{-}}\ensuremath{^{7}} & 20 & \textasciitilde{}20 nm of interphase is enough to trigger plating at 1C \\
\hline
10\ensuremath{^{-}}\ensuremath{^{8}} & 2 & 2 nm suffices --- below a single monolayer's worth of compact SEI \\
\hline
10\ensuremath{^{-}}\ensuremath{^{9}} & 0.2 & sub-nanometer --- physically implausible for a continuous film \\
\hline
\end{tabular}
\endgroup



\texttt{[DERIVED]}. The result is informative because it converts an abstract threshold into a length: at the upper end of the conductivity band, a mere 20 nm of interphase is sufficient to push a 1C cell into plating, and that is within the 10--200 nm a SEI reaches after formation. The threshold is not a distant failure mode; it is a near-term one. Conversely the \ensuremath{\sigma} = 10\ensuremath{^{-}}\ensuremath{^{9}} row is \emph{self-refuting}: sub-nanometer continuous films do not exist, which is independent evidence that the effective conductivity cannot be at the bottom of the band. This corroborates the polarization cross-check in S2.1 from a completely different direction, and the two together are the support available for the \ensuremath{\sigma}-band argument.


\textbf{Caveat.} \emph{d}* inherits every uncertainty of both \emph{R}* (a lumped single-\ensuremath{\Delta}, single-\emph{R} approximation) and \ensuremath{\sigma} (two orders of magnitude). It is useful as a scaling statement (d* \ensuremath{\propto} 1/j, so doubling the rate halves the interphase thickness needed to plate) and should not be quoted as a number for a specific cell.


\par\noindent\rule{\linewidth}{0.4pt}\par


\subsection*{S2.3 SEI volume and the pore-filling threshold}


This is the derivation that connects Loop A to Loop C: interphase growth does not merely add resistance, it \textbf{consumes anode pore volume}, and there is a finite amount of it. The link is a mass balance.


\textbf{Step 1 --- charge to moles.} Reducing one Li\ensuremath{^{+}} and incorporating it into a typical SEI component consumes one electron per Li (for Li\ensuremath{_{2}}CO\ensuremath{_{3}}, two Li per two electrons, the 1:1 ratio holds). One mAh cm\ensuremath{^{-}}\ensuremath{^{2}} is 10\ensuremath{^{-}}\ensuremath{^{3}} A \ensuremath{\times} 3600 s = 3.6 C cm\ensuremath{^{-}}\ensuremath{^{2}}, so


\begin{equation}n_{\text{SEI}} = \frac{3.6}{96485} = 3.73\times10^{-5}\ \text{mol cm}^{-2}\ \text{per mAh cm}^{-2} \tag{S2.E6}\end{equation}


\texttt{[DERIVED]}. Using \emph{F} = 96485 C mol\ensuremath{^{-}}\ensuremath{^{1}}.


\textbf{Step 2 --- moles to volume.} With a representative SEI molar volume \emph{V}\textsubscript{m} \ensuremath{\approx} 30 cm\ensuremath{^{3}} mol\ensuremath{^{-}}\ensuremath{^{1}}:


\begin{equation}t_{\text{SEI per mAh}} = 3.73\times10^{-5}\ \text{mol cm}^{-2} \times 30\ \text{cm}^{3}\ \text{mol}^{-1} = 1.12\times10^{-3}\ \text{cm} = \mathbf{11.2\ \mu m} \tag{S2.E7}\end{equation}


\texttt{[DERIVED]}: i.e. \textasciitilde{}11 \ensuremath{\mu}m of SEI-equivalent thickness per mAh cm\ensuremath{^{-}}\ensuremath{^{2}} of lost lithium inventory.


\textbf{A labelling caution.} The 30 cm\ensuremath{^{3}} mol\ensuremath{^{-}}\ensuremath{^{1}} above is a representative value for a mixed SEI, not the value for any single compound. The Li\ensuremath{_{2}}CO\ensuremath{_{3}} end-member has \emph{V}\textsubscript{m} = 34.8 cm\ensuremath{^{3}} mol\ensuremath{^{-}}\ensuremath{^{1}} and accordingly gives


\begin{equation}3.73\times10^{-5}\times34.8\times10^{4} = \mathbf{12.9\ \mu m}\ \text{per mAh cm}^{-2} \tag{S2.E8}\end{equation}


and the plausible band \emph{V}\textsubscript{m} = 20--40 cm\ensuremath{^{3}} mol\ensuremath{^{-}}\ensuremath{^{1}} gives 7.4--14.8 \ensuremath{\mu}m. 11 \ensuremath{\mu}m and 12.9 \ensuremath{\mu}m are different numbers answering different questions, representative SEI versus Li\ensuremath{_{2}}CO\ensuremath{_{3}} end-member, and conflating them is how a factor-of-1.17 discrepancy gets mistaken for a bug. The two are therefore quoted against their own molar volumes (11 \ensuremath{\mu}m at 30 cm\ensuremath{^{3}} mol\ensuremath{^{-}}\ensuremath{^{1}}; 12.9 \ensuremath{\mu}m at 34.8 cm\ensuremath{^{3}} mol\ensuremath{^{-}}\ensuremath{^{1}}) and must not be interchanged.


Step 3 --- volume to fraction of available pore space. Take a representative anode pore volume of \textasciitilde{}3 \ensuremath{\mu}m of thickness per unit area. That figure corresponds to a coating about 10 \ensuremath{\mu}m thick at 30 \% porosity, which is a low-loading electrode; against a 3 mAh cm\ensuremath{^{-}}\ensuremath{^{2}} cathode it is N/P \ensuremath{\approx} 0.13. The self-consistent anode loading for such a cathode is \textasciitilde{}77 \ensuremath{\mu}m, with \textasciitilde{}23 \ensuremath{\mu}m of pore volume, and that is the case evaluated in Step 5 below. The LLI required to fill it:


\begin{equation}LLI_{\text{fill}} = \frac{3\ \mu m}{11.2\ \mu m\ \text{per}\ (\text{mAh cm}^{-2})} = \mathbf{0.27\ mAh\ cm^{-2}} \tag{S2.E9}\end{equation}


Against \emph{Q}\ensuremath{_{0}} = 3 mAh cm\ensuremath{^{-}}\ensuremath{^{2}}, that is 8.9 \% of capacity. \texttt{[DERIVED]}. So interphase growth alone reaches the pore-filling limit after losing under 10 \% of the cell's capacity, well before any capacity target of 80 \% retention is threatened, and therefore an order-unity margin, not a 10\ensuremath{^{-}}\ensuremath{^{2}} one.


\textbf{Step 4, the porosity correction.} The 0.27 mAh cm\ensuremath{^{-}}\ensuremath{^{2}} figure treats the interphase as fully dense, which it is not: the SEI is 30--50 \% porous \texttt{[ASSUMED]}. The correction matters, and it must be stated with its convention, because two different porosity conventions appear in this literature and they move the threshold in opposite directions. Stating which one is used is the difference between a derivation and a fit.


\emph{Convention A --- effective solid volume available to block transport.} The fraction of the interphase volume that is solid is (1 \ensuremath{-} \emph{p}), so the thickness per mAh that actually contributes to blocking becomes 11.2 \ensuremath{\times} (1 \ensuremath{-} \emph{p}) \ensuremath{\mu}m. For \emph{p} = 30--50 \% this is 5.5--7.7 \ensuremath{\mu}m per mAh cm\ensuremath{^{-}}\ensuremath{^{2}}, and


\begin{equation}LLI_{\text{fill}} = \frac{3\ \mu m}{5.5\text{--}7.7\ \mu m\ \text{per}\ (\text{mAh cm}^{-2})} \approx \mathbf{0.39\text{--}0.55\ mAh\ cm^{-2}} \tag{S2.E10}\end{equation}


i.e. 13--18 \% of a 3 mAh cm\ensuremath{^{-}}\ensuremath{^{2}} capacity. \texttt{[DERIVED]}.


\emph{Convention B --- geometric void filling.} If instead one asks how much solid SEI is needed to occupy a pore of fixed geometric volume, a porous film reaches that volume with less solid, giving 3(1 \ensuremath{-} \emph{p})/11.2 = 0.13--0.19 mAh cm\ensuremath{^{-}}\ensuremath{^{2}}, i.e. below the dense figure. \texttt{[DERIVED]}.


The two conventions differ by a factor of (1 \ensuremath{-} \emph{p})\ensuremath{^{2}} \ensuremath{\approx} 0.25--0.49, so the number is convention-dependent while the conclusion is not. Under convention A the threshold is 0.39--0.55 mAh cm\ensuremath{^{-}}\ensuremath{^{2}}; under convention B it is 0.13--0.19; under the dense assumption it is 0.27. Every one of these is far below 3 mAh cm\ensuremath{^{-}}\ensuremath{^{2}}, by a factor of 5.5 to 23, so the claim that \textquotedblleft{}interphase growth reaches the pore-filling limit long before a capacity target of 80 \% retention is threatened\textquotedblleft{} survives either convention unchanged. The range can therefore be quoted without resolving the ambiguity.


The range adopted here: \textasciitilde{}0.3--0.6 mAh cm\ensuremath{^{-}}\ensuremath{^{2}}, i.e. \textasciitilde{}10--20 \% of capacity. This is convention A (0.39--0.55) widened at the lower edge toward the dense case (0.27) and at the upper edge to cover a 4 \ensuremath{\mu}m pore volume. Quoting only 0.27 would understate the threshold; quoting only 0.55 would overstate the dense case; quoting convention B's 0.13--0.19 would be defensible arithmetic but would sit below the dense case and invite the question of why porosity reduced the requirement.


\textbf{Step 5, the self-consistent anode.} The 0.27--0.55 mAh cm\ensuremath{^{-}}\ensuremath{^{2}} thresholds above assume \textasciitilde{}3--4 \ensuremath{\mu}m of pore volume, i.e. a low-loading anode. At the \emph{Q}\ensuremath{_{0}} = 3 mAh cm\ensuremath{^{-}}\ensuremath{^{2}} used in the model of S1, the graphite anode is \textasciitilde{}77 \ensuremath{\mu}m thick at 1.6 g cm\ensuremath{^{-}}\ensuremath{^{3}} graphite density and 30 \% porosity, giving \textasciitilde{}23 \ensuremath{\mu}m of pore volume and an areal capacity of \textasciitilde{}2.06 mAh cm\ensuremath{^{-}}\ensuremath{^{2}}, so N/P \ensuremath{\approx} 0.69 against the 3 mAh cm\ensuremath{^{-}}\ensuremath{^{2}} cathode. Filling that pore volume requires


\begin{equation}LLI_{\text{fill}} = \frac{23\ \mu m}{11.2\ \mu m\ \text{per}\ (\text{mAh cm}^{-2})} \approx \mathbf{2.1\ mAh\ cm^{-2}} \tag{S2.11}\end{equation}


\texttt{[DERIVED]}. That is \textasciitilde{}68 \% of \emph{Q}\ensuremath{_{0}}, not \textasciitilde{}9 \%. The correction factor is the anode loading, not the SEI porosity convention, and it is the larger of the two.


The consequence is stated plainly because it changes a conclusion. On the self-consistent anode, pore filling is not reached until most of the capacity is gone, so it is a late-life rather than an early-life mechanism. Loop C therefore does not close early on mechanical grounds at this loading. Whether it closes at all depends on the porosity, and on thinner anodes the margin falls by the ratio of loadings: at a 10 \ensuremath{\mu}m anode the threshold is 0.27 mAh cm\ensuremath{^{-}}\ensuremath{^{2}} and the loop closes at \textasciitilde{}9 \% loss. Both regimes are real; they differ by the electrode design, and the range is reported here rather than a single number.


Coupling to the model. In the reduced model (S1), the coupled cell's LLI crosses 0.27 mAh cm\ensuremath{^{-}}\ensuremath{^{2}} at cycle 386, 144 cycles before its capacity knee at cycle 530. (Use the printed 0.27 rather than the unrounded 0.2680: the cross is then at cycle 386, whereas the unrounded threshold gives 385. A reader re-deriving from the text should reach the printed number, so the rounded value is the one quoted.) Over the quoted range, the crossing moves from 386 (0.27) to 450 (0.39) to 513 (0.55), and the LLI ratio at the knee is 530's LLI divided by the threshold: 2.22\ensuremath{\times} on the dense basis, 1.54\ensuremath{\times} at 0.39, 1.09\ensuremath{\times} at 0.55. So on any basis the pore-filling threshold is reached before the capacity knee, and the model-level statement of \textquotedblleft{}Loop A feeds Loop C\textquotedblright{} is that the interphase has already filled the pore volume by the time the capacity curve turns over, with the margin depending on which threshold convention is used, which is itself the reason not to quote a single number. This is a model result \texttt{[DERIVED]}, not a measurement, and it inherits the model's parameters.


\par\noindent\rule{\linewidth}{0.4pt}\par


\subsection*{S2.4 Coulombic-efficiency arithmetic and compounding laws}


Coulombic efficiency is the most-quoted number in this field and the most often mis-converted into a cycle life. Two different arithmetic laws are in use, and the same CE implies a different cycle life under each, with the gap widening as CE falls. The section settles the arithmetic, then uses it to draw the discriminator between the two channels that consume lithium inventory.


The inventory. Lithium inventory in a 3 mAh cm\ensuremath{^{-}}\ensuremath{^{2}} cell, expressed as a thickness of pure Li:


\begin{equation}m_{\text{Li}} = 3\times10^{-3}\ \text{A h cm}^{-2}\times3600\ \text{C A}^{-1}\text{h}^{-1} / 96485 \times 6.941\ \text{g mol}^{-1} = 7.77\times10^{-4}\ \text{g cm}^{-2} \nonumber\end{equation}


\begin{equation}t_{\text{Li}} = \frac{7.77\times10^{-4}}{0.534} = 1.455\times10^{-3}\ \text{cm} = \mathbf{14.55\ \mu m} \tag{S2.12}\end{equation}


\texttt{[DERIVED]}, with \ensuremath{\rho}\textsubscript{Li} = 0.534 g cm\ensuremath{^{-}}\ensuremath{^{3}}. At CE = 99.9 \%, the absolute loss per cycle is 0.1 \% \ensuremath{\times} 14.55 \ensuremath{\mu}m = 14.55 nm of Li-equivalent per cycle.


Two laws, and the arithmetic of each.


\begin{itemize}
  \item Model A --- constant absolute loss per cycle. Loss is (1 \ensuremath{-} CE)\ensuremath{\cdot}\emph{Q} each cycle, independent of state. In this law the 20 \%-retention life is \emph{N}\ensuremath{_{8}}\ensuremath{_{0}} = 0.20/(1 \ensuremath{-} CE), and full exhaustion is \emph{N} = 1/(1 \ensuremath{-} CE). At 99.9 \%: \emph{N}\ensuremath{_{8}}\ensuremath{_{0}} = 200 cycles, exhaustion = 1000 cycles.
  \item Model B --- compounding (exponential) loss. Retained capacity is the CE raised to the cycle count, so the CE compounds. Then \emph{N}\ensuremath{_{8}}\ensuremath{_{0}} = ln 0.8/ln CE, and exhaustion is \emph{N} = ln 0.01/ln CE. At 99.9 \%: \emph{N}\ensuremath{_{8}}\ensuremath{_{0}} = ln 0.8/ln 0.999 = 223 cycles, exhaustion = \textbf{4603 cycles}.
\end{itemize}


\texttt{[DERIVED]}: all four numbers.


\textbf{The comparison, which is the point:}


\begingroup\scriptsize
\setlength{\tabcolsep}{3pt}
\renewcommand{\arraystretch}{1.12}
\sloppy
\begin{tabular}{@{}>{\raggedright\arraybackslash}p{0.079\textwidth}>{\raggedright\arraybackslash}p{0.172\textwidth}>{\raggedright\arraybackslash}p{0.152\textwidth}>{\raggedright\arraybackslash}p{0.074\textwidth}>{\raggedright\arraybackslash}p{0.156\textwidth}>{\raggedright\arraybackslash}p{0.162\textwidth}>{\raggedright\arraybackslash}p{0.074\textwidth}@{}}
\hline
\textbf{CE} & \textbf{\emph{N}\ensuremath{_{8}}\ensuremath{_{0}}, constant-loss} & \textbf{\emph{N}\ensuremath{_{8}}\ensuremath{_{0}}, compounding} & \textbf{ratio} & \textbf{Exhaustion, constant} & \textbf{Exhaustion, compounding} & \textbf{ratio} \\ \hline
99.90 \% & 200 & 223 & 1.12 & 1000 & 4603 & 4.60 \\
\hline
99.95 \% & 400 & 446 & 1.12 & 2000 & 9208 & 4.60 \\
\hline
99.98 \% & 1000 & 1116 & 1.12 & 5000 & 23024 & 4.60 \\
\hline
99.99 \% & 2000 & 2231 & 1.12 & 10000 & 46049 & 4.60 \\
\hline
\end{tabular}
\endgroup



The conclusions conflict.


\begin{enumerate}
  \item For 80 \%-retention claims the law barely matters, the two models agree within 11.6 \% at every CE in the table. So a paper that says \textquotedblleft{}CE = 99.9 \% gives \textasciitilde{}200 cycles to 80 \% retention\textquotedblright{} is safe under either law, PROVIDED it says 80 \% retention. This is the common case and the risk is low.
  \item For \textquotedblleft{}exhaustion\textquotedblright{} or \textquotedblleft{}infinite-life\textquotedblright{} claims the law dominates by a factor of 4.6. A 99.9 \% cell is exhausted at 1000 cycles under constant loss but survives 4603 under compounding. Anyone quoting a cycle count at a given CE must state which law is assumed, because the number can be wrong by 4.6\ensuremath{\times} without any arithmetic error. A figure of \textquotedblleft{}1000 cycles at 99.9 \% CE\textquotedblright{} is a statement about one assumed law, not a measurement of cycle life.
\end{enumerate}


The CE required for a target life. Inverting Model B for \emph{N} cycles at 80 \% retention gives a required CE equal to the \emph{N}-th root of 0.8:


\begingroup\small
\setlength{\tabcolsep}{3pt}
\renewcommand{\arraystretch}{1.12}
\sloppy
\begin{tabular}{@{}>{\raggedright\arraybackslash}p{0.237\textwidth}>{\raggedright\arraybackslash}p{0.237\textwidth}>{\raggedright\arraybackslash}p{0.237\textwidth}>{\raggedright\arraybackslash}p{0.194\textwidth}@{}}
\hline
\textbf{Target life (cycles)} & \textbf{Required CE, Model B} & \textbf{Required CE, Model A} & \textbf{Difference} \\ \hline
200 & 99.8885 \% & 99.9000 \% & 0.0115 pp \\
\hline
500 & 99.9554 \% & 99.9600 \% & 0.0046 pp \\
\hline
1000 & 99.9777 \% & 99.9800 \% & 0.0023 pp \\
\hline
2000 & 99.9888 \% & 99.9900 \% & 0.0012 pp \\
\hline
\end{tabular}
\endgroup



\texttt{[DERIVED]}. The two laws agree to within 0.012 percentage points, far below any real CE measurement's precision (state-of-the-art precision coulometry reaches \textasciitilde{}0.001 pp, so this difference is at the edge of resolvability but not practically important). The practical message: to reach 1000 cycles at 80 \% retention requires CE \ensuremath{\geq} 99.978 \%, which is a requirement most cells fail. That is a much more useful statement than any cycle-count prediction, because it is a specification against which a design can be checked.


\textbf{Why this matters.} Loop A's parasitic current is the physical origin of CE loss, so the CE requirement is the quantitative bridge between the mechanism (interphase growth) and the performance target (cycle life). It also sets the scale for Loop D: 99.978 \% CE on a 3.0 mAh cm\ensuremath{^{-}}\ensuremath{^{2}} electrode means 0.66 \ensuremath{\mu}Ah cm\ensuremath{^{-}}\ensuremath{^{2}} of charge lost per cycle, which on the 11 \ensuremath{\mu}m/mAh basis of S2.3 is \textbf{\textasciitilde{}7.3 nm of interphase per cycle}. That is the number to compare with Loop B's dissolution budget and Loop D's salt budget.


The discriminator between the two loss channels. Losing a charge \emph{Q} can occur either as dead lithium or as SEI. The two channels differ in the thickness of material they consume per unit of charge:


\begin{itemize}
  \item Dead lithium, at 1:1 Li, consumes 4.85 \ensuremath{\mu}m per mAh cm\ensuremath{^{-}}\ensuremath{^{2}} of inventory, from \ensuremath{\rho}\textsubscript{Li} and \emph{M}\textsubscript{Li.}
  \item SEI formation consumes 11 \ensuremath{\mu}m per mAh cm\ensuremath{^{-}}\ensuremath{^{2}} on the molar-volume basis of S2.3.
\end{itemize}


Both channels scale linearly with the charge lost, so their ratio does not depend on CE: 11/4.85 = \textbf{2.27\ensuremath{\times}}. The two channels are therefore of the same order, with SEI consuming about twice the thickness per unit of inventory lost. A thickness basis ten times smaller would instead give 22\ensuremath{\times}, which would argue the opposite, namely one channel dominant. The 11 \ensuremath{\mu}m/mAh basis above is the one used, and the discriminator is \textasciitilde{}2.3\ensuremath{\times}.


\par\noindent\rule{\linewidth}{0.4pt}\par


\subsection*{S2.5 Loop B: the catalytic versus stoichiometric shortfall (\textasciitilde{}74\ensuremath{\times})}


Loop B claims that dissolved transition metal poisons the anode, raises the local overpotential, and so drives further dissolution. That claim runs into an immediate objection: dissolution is a stoichiometric process, one metal ion per dissolution event, so a small amount of manganese cannot in itself catalyze anything. The resolution is that the constraint is not the amount of manganese but its turnover - how many times each ion can cycle between cathode and anode. This section computes the shortfall between a stoichiometric account and the dissolution rate that Loop B requires.


\textbf{Step 1, the Mn inventory.} For a 20 mg cm\ensuremath{^{-}}\ensuremath{^{2}} cathode at 20 wt \% Mn (NMC-class), with \emph{M}\textsubscript{Mn} = 54.94 g mol\ensuremath{^{-}}\ensuremath{^{1}}:


\begin{equation}n_{\text{Mn}} = \frac{20\times10^{-3}\times0.20}{54.94} = 7.28\times10^{-5}\ \text{mol cm}^{-2} = \mathbf{72.8\ \mu mol\ cm^{-2}} \tag{S2.13}\end{equation}


\texttt{[DERIVED]}.


Step 2 --- how much actually dissolves over a realistic life. Measured dissolution rates in NMC cells are of order 5 ppm of the cathode's Mn per cycle (\texttt{[ASSUMED]}, an order-of-magnitude figure; no single published value is quoted). Over 1000 cycles:


\begin{equation}n_{\text{dissolved}} = 7.28\times10^{-5}\times5\times10^{-6}\times1000 = 3.64\times10^{-7}\ \text{mol cm}^{-2} = \mathbf{0.364\ \mu mol\ cm^{-2}} \tag{S2.14}\end{equation}


As a fraction of inventory: 3.64\ensuremath{\times}10\ensuremath{^{-}}\ensuremath{^{7}}/7.28\ensuremath{\times}10\ensuremath{^{-}}\ensuremath{^{5}} = 0.50 \%. \texttt{[DERIVED]}.


\textbf{Step 3, the shortfall.} A cell targeting 20 \% capacity loss over that 1000 cycles loses 0.20 \ensuremath{\times} 3 mAh cm\ensuremath{^{-}}\ensuremath{^{2}} = 0.6 mAh cm\ensuremath{^{-}}\ensuremath{^{2}} of inventory, which is 0.6 \ensuremath{\times} 3.73\ensuremath{\times}10\ensuremath{^{-}}\ensuremath{^{5}} = 2.24\ensuremath{\times}10\ensuremath{^{-}}\ensuremath{^{5}} mol cm\ensuremath{^{-}}\ensuremath{^{2}}. Compare with the 3.64\ensuremath{\times}10\ensuremath{^{-}}\ensuremath{^{7}} mol cm\ensuremath{^{-}}\ensuremath{^{2}} of Mn that actually dissolved:


\begin{equation}\frac{\text{loss to be explained}}{\text{Mn dissolved}} = \frac{0.6\ \text{mAh cm}^{-2}}{0.50\%\ \text{of Mn}}:\quad \frac{20\%}{0.50\%} = \mathbf{40\times} \tag{S2.15}\end{equation}


\texttt{[DERIVED]}. Dissolved Mn accounts for \textasciitilde{}2.5 \% of the capacity loss it is supposed to cause.


Reconciling 40\ensuremath{\times} with 74\ensuremath{\times}, the denominators differ. The 40\ensuremath{\times} above is a percentage ratio: 20 \% capacity loss against 0.50 \% of the Mn inventory. But Mn is not the only inventory in play, and comparing a fraction of Mn against a fraction of capacity mixes two denominators. On a mole-for-mole basis: the Mn inventory (7.281\ensuremath{\times}10\ensuremath{^{-}}\ensuremath{^{5}} mol cm\ensuremath{^{-}}\ensuremath{^{2}}) is only 54.2 \% of the Li inventory (1.3432\ensuremath{\times}10\ensuremath{^{-}}\ensuremath{^{4}} mol cm\ensuremath{^{-}}\ensuremath{^{2}} at 3.6 mAh cm\ensuremath{^{-}}\ensuremath{^{2}}; Li/Mn = 1.845). The 3.6 mAh cm\ensuremath{^{-}}\ensuremath{^{2}} basis is the high-energy cell this section uses, and it is the basis on which Li/Mn = 1.845 is quoted; the reduced model's reference cell in S1 is 3.0 mAh cm\ensuremath{^{-}}\ensuremath{^{2}}. Each departing Mn\ensuremath{^{2}}\ensuremath{^{+}} destroys one Li\ensuremath{^{+}}, so the dissolved Mn removes 0.50 \% \ensuremath{\times} 0.542 = 0.271 \% of the Li inventory. Against a 20 \% capacity loss the shortfall is 0.20 / 0.00271 = 73.8\ensuremath{\times}, i.e. \textasciitilde{}74\ensuremath{\times}. On the 3.0 mAh cm\ensuremath{^{-}}\ensuremath{^{2}} basis the same arithmetic gives 61.5\ensuremath{\times}, and the conclusion is unchanged either way: the shortfall is two orders of magnitude, so dissolution cannot be stoichiometrically sufficient. The 40\ensuremath{\times} figure is arithmetically correct but is not the physical shortfall; the mole-for-mole ratio is what the turnover argument below must close. \texttt{[DERIVED]}.


How the loop closes the gap is set out below. The \textasciitilde{}74\ensuremath{\times} mole-for-mole gap is the whole reason Loop B must be catalytic rather than stoichiometric. A Mn\ensuremath{^{2}}\ensuremath{^{+}} ion that deposits on the anode is not consumed by doing so: it sits in the SEI, raises its electronic conductivity, and, critically, can be re-released into solution as the SEI is rebuilt, depositing again elsewhere. Each deposited ion can therefore poison more than one site over the cell's life, and the relevant quantity is a turnover number, not a concentration. If each Mn ion cycles through the anode \textasciitilde{}74 times over 1000 cycles, a turnover every \textasciitilde{}14 cycles, which is plausible for an ion exchanging between SEI and electrolyte as the interphase is continuously rebuilt. The budget then closes. \texttt{[HYPOTHESIS]}: the turnover number is not established in the literature; the derivation shows only that a modest turnover suffices, which is a far weaker and much more defensible claim than asserting a specific catalytic mechanism.


\textbf{The supported claim, and its limit.} The supported statement is that \emph{a stoichiometric accounting falls short by \textasciitilde{}74\ensuremath{\times}, so dissolution must act catalytically if it matters at all --- a testable claim rather than an assumption.} The catalytic mechanism is not established, and no turnover number is quoted as measured. The admissible form of Loop B is therefore a conditional: if dissolution causes a substantial share of fade, it does so catalytically, and the required turnover is \textasciitilde{}74. That conditional is falsifiable, and S4's branch B7 sets out the experiment (HF scavenger) that tests it.


\textbf{Sensitivity.} The shortfall scales as 1/(dissolution rate): 2 ppm/cycle gives 184\ensuremath{\times} rather than 74\ensuremath{\times}, and 20 ppm/cycle gives 18\ensuremath{\times}. The \emph{order of magnitude} (tens, not units) is robust across the plausible range, which is all the argument needs. The claim would fail only at \textasciitilde{}200 ppm/cycle, which would mean 20 \% of the cathode's Mn dissolving in 1000 cycles, a catastrophic degradation mode inconsistent with the cells in which Loop B is invoked.


\par\noindent\rule{\linewidth}{0.4pt}\par


\subsection*{S2.6 Loop D: the salt-inventory budget}


A lean cell contains only a few microliters of electrolyte per square centimeter, and every lithium atom that enters the interphase takes electrolyte with it. Loop D claims that this consumption reaches a wall before the capacity target does. Unlike Loops A to C, the argument here is a budget: total what will be consumed over the cell's life and compare it with what is present at the start. The margin turns out to be about a factor of 9, and, more usefully, it depends entirely on cell format.


Step 1, the consumption per unit LLI. Reducing Li\ensuremath{^{+}} to form SEI consumes \textasciitilde{}1 Li per electron, so


\begin{equation}n_{\text{per mAh}} = \frac{3.6\ \text{C cm}^{-2}}{96485\ \text{C mol}^{-1}} = 3.73\times10^{-5}\ \text{mol cm}^{-2}\ \text{per mAh cm}^{-2} \tag{S2.16}\end{equation}


\texttt{[DERIVED]} (the same constant as S2.3; the Faraday arithmetic appears in Loops A/C/D, so one conversion underlies three of the four loops).


For a 4 mAh cm\ensuremath{^{-}}\ensuremath{^{2}} cell (a larger format than S1's 3.0 mAh cm\ensuremath{^{-}}\ensuremath{^{2}} reference, chosen so the budget is stressed) losing 20 \% over its life, LLI \ensuremath{\approx} 0.8 mAh cm\ensuremath{^{-}}\ensuremath{^{2}}, so


\begin{equation}n_{\text{consumed}} = 0.8 \times 3.73\times10^{-5} = 2.98\times10^{-5}\ \text{mol cm}^{-2} = \mathbf{29.8\ \mu mol\ cm^{-2}} \tag{S2.17}\end{equation}


\texttt{[DERIVED]}.


\textbf{Step 2, the inventory present.} At a lean electrolyte loading of 1 g Ah\ensuremath{^{-}}\ensuremath{^{1}} for a 4 mAh cm\ensuremath{^{-}}\ensuremath{^{2}} electrode:


\begin{equation}m_{\text{elec}} = 4\times10^{-3}\ \text{Ah cm}^{-2}\times1\ \text{g Ah}^{-1} = 4\times10^{-3}\ \text{g cm}^{-2} \nonumber\end{equation}


\begin{equation}V_{\text{elec}} = \frac{4\times10^{-3}}{1.2} = 3.33\times10^{-3}\ \text{cm}^{3}\ \text{cm}^{-2} = \mathbf{3.33\ \mu L\ cm^{-2}} \tag{S2.18}\end{equation}


\texttt{[DERIVED]}, with \ensuremath{\rho}\textsubscript{electrolyte} \ensuremath{\approx} 1.2 g cm\ensuremath{^{-}}\ensuremath{^{3}}. Applying 1 M salt and \textasciitilde{}12 M total solvent+anion species:


\begin{equation}n_{\text{Li+}} = 1\ \text{mol L}^{-1}\times3.33\times10^{-6}\ \text{L cm}^{-2} = \mathbf{3.33\ \mu mol\ cm^{-2}} \tag{S2.19}\end{equation}


\begin{equation}n_{\text{total}} = 12\ \text{mol L}^{-1}\times3.33\times10^{-6}\ \text{L cm}^{-2} = \mathbf{40\ \mu mol\ cm^{-2}} \tag{S2.20}\end{equation}


\textbf{Step 3, the budget:}


\begingroup\small
\setlength{\tabcolsep}{3pt}
\renewcommand{\arraystretch}{1.12}
\sloppy
\begin{tabular}{@{}>{\raggedright\arraybackslash}p{0.231\textwidth}>{\raggedright\arraybackslash}p{0.172\textwidth}>{\raggedright\arraybackslash}p{0.257\textwidth}>{\raggedright\arraybackslash}p{0.244\textwidth}@{}}
\hline
\textbf{Inventory} & \textbf{Present (\ensuremath{\mu}mol cm\ensuremath{^{-}}\ensuremath{^{2}})} & \textbf{Consumed for 0.8 mAh cm\ensuremath{^{-}}\ensuremath{^{2}} LLI (\ensuremath{\mu}mol cm\ensuremath{^{-}}\ensuremath{^{2}})} & \textbf{Consumed/present} \\ \hline
Li\ensuremath{^{+}} salt only & 3.33 & 29.8 & 8.95\ensuremath{\times} \\
\hline
Total electrolyte species & 40 & 29.8 & 0.75 \\
\hline
\end{tabular}
\endgroup



\texttt{[DERIVED]}. The Li\ensuremath{^{+}} salt inventory is exceeded almost ninefold by the lithium projected to be lost, and even the total electrolyte inventory is 75 \% consumed. Salt depletion is therefore not a second-order correction at long life; it is a hard limit.


\textbf{The principal caveat.} The 1:1 Li-to-electrolyte stoichiometry is an upper bound: some SEI components consume Li without consuming a separate solvent molecule (LiF from LiPF\ensuremath{_{6}}), and some electrolyte is consumed without consuming Li. The \textasciitilde{}9\ensuremath{\times} ratio is therefore an order-of-magnitude result. The argument is nonetheless robust to the caveat: even a 3\ensuremath{\times} overestimate leaves a ratio of \textasciitilde{}3\ensuremath{\times}, which is still enough to make Loop D first-order. The conclusion does not depend on the stoichiometric detail. \texttt{[DERIVED]} for the ratio; \texttt{[HYPOTHESIS]} for the 1:1 stoichiometry as a model of real SEI chemistry.


\textbf{The format dependence, the practical payoff.} A flooded coin cell with 50 \ensuremath{\mu}L of electrolyte instead of 3.33 \ensuremath{\mu}L cm\ensuremath{^{-}}\ensuremath{^{2}} has 15\ensuremath{\times} more absolute inventory (50/3.33). So the same chemistry that is salt-starved in a lean pouch cell is salt-rich in a flooded coin cell, and coin-cell data cannot be extrapolated to pouch cells on this axis. This explains, with no new chemistry, why \textquotedblleft{}electrolyte volume\textquotedblright{} changes knee position without changing the initial fade slope, and why thick high-areal-capacity electrodes are systematically worse. It also gives S4's branch B5 its discriminator: a Loop-D-limited cell improves when flooded; a genuinely M4-limited cell does not.


\par\noindent\rule{\linewidth}{0.4pt}\par


\subsection*{S2.7 Loop C: diffusion-induced stress}


Cracking a secondary particle exposes fresh surface, interphase grows on it, and the interphase then becomes load-bearing and changes the local mechanical constraint. Loop C claims that the stress driving this process is of the same order as the stress the particle can withstand, so fracture is neither impossible nor guaranteed. The derivation has to be careful about which stress it computes, because two very different numbers are in play, and they are easily conflated.


\textbf{Step 1, the concentration swing.} For layered NMC the molar volume is \emph{V}\textsubscript{m} \ensuremath{\approx} 19.4 \ensuremath{\times} 10\ensuremath{^{-}}\ensuremath{^{6}} m\ensuremath{^{3}} mol\ensuremath{^{-}}\ensuremath{^{1}}, so the maximum Li concentration change is


\begin{equation}\Delta c_{\max} = \frac{1}{V_m} = \frac{1}{19.4\times10^{-6}} = 5.15\times10^{4}\ \text{mol m}^{-3} = \mathbf{0.0515\ mol\ cm^{-3}} \tag{S2.21}\end{equation}


\texttt{[DERIVED]}. A value of 0.048 mol cm\ensuremath{^{-}}\ensuremath{^{3}} is adopted here, a \textasciitilde{}7 \% reduction reflecting the fact that not the full molar volume participates in a given cycle's swing. Both are correct answers to slightly different questions; 0.048 is the figure used downstream.


Step 2, the elastic upper bound, which is NOT the fracture-relevant number. For a spherical particle with a composition gradient under elastic (no-relaxation) conditions, the peak hoop stress is


\begin{equation}\sigma_{\max} = \frac{E\,\Omega\,\Delta c}{3(1-\nu)} \tag{S2.22}\end{equation}


With \emph{E} = 150 GPa, \ensuremath{\Omega} = 3.5 \ensuremath{\times} 10\ensuremath{^{-}}\ensuremath{^{6}} m\ensuremath{^{3}} mol\ensuremath{^{-}}\ensuremath{^{1}}, \ensuremath{\nu} = 0.3, \ensuremath{\Delta}\emph{c} = 0.048 mol cm\ensuremath{^{-}}\ensuremath{^{3}} = 4.8 \ensuremath{\times} 10\ensuremath{^{4}} mol m\ensuremath{^{-}}\ensuremath{^{3}}:


\begin{equation}\sigma_{\max} = \frac{150\times10^{9}\times3.5\times10^{-6}\times4.8\times10^{4}}{3\times0.7} = \mathbf{12.0\ GPa} \tag{S2.23}\end{equation}


\texttt{[DERIVED]}. The 12.0 GPa follows from the adopted \ensuremath{\Delta}c = 0.048 mol cm\ensuremath{^{-}}\ensuremath{^{3}}; the unrounded 0.0515 mol cm\ensuremath{^{-}}\ensuremath{^{3}} would give 12.9 GPa, so the choice of \ensuremath{\Delta}c moves the ceiling by 7 \%. This is an absolute upper bound assuming the particle behaves elastically with the full composition swing simultaneously. This is physically unreachable, because plastic/creep relaxation and finite lithiation rates reduce it by an order of magnitude. It is useful only as a ceiling.


Step 3, the realistic stress, obtained by inverting the same relation. Measured/estimated peak stresses in Ni-rich secondary particles are 0.5--2 GPa. The corresponding gradient fraction is


\begin{equation}\frac{\Delta c_{\text{real}}}{\Delta c_{\max}} = \frac{\sigma_{\text{real}}\,3(1-\nu)}{E\,\Omega\,\Delta c_{\max}} \nonumber\end{equation}


For \ensuremath{\sigma}\textsubscript{real} = 1.5 GPa: \ensuremath{\Delta}\emph{c}\textsubscript{real}/\ensuremath{\Delta}\emph{c}\textsubscript{max} = 1.5\ensuremath{\times}10\ensuremath{^{9}} \ensuremath{\times} 2.1/(150\ensuremath{\times}10\ensuremath{^{9}} \ensuremath{\times} 3.5\ensuremath{\times}10\ensuremath{^{-}}\ensuremath{^{6}} \ensuremath{\times} 4.8\ensuremath{\times}10\ensuremath{^{4}}) = 12.5 \%. \texttt{[DERIVED]}. So a realistic 1.5 GPa stress corresponds to only \textasciitilde{}12 \% of the maximum concentration swing being accommodated elastically at any instant, with the rest relaxed.


Step 4 --- compare with the fracture stress. Using \emph{K}\textsubscript{Ic} = 2 MPa m\textsuperscript{\ensuremath{\tfrac{1}{2}}}, the stress required to propagate a flaw of size \emph{a} is \ensuremath{\sigma}\textsubscript{f} = \emph{K}\textsubscript{Ic}/\ensuremath{\surd}(\ensuremath{\pi}\emph{a}):


\begingroup\small
\setlength{\tabcolsep}{3pt}
\renewcommand{\arraystretch}{1.12}
\sloppy
\begin{tabular}{@{}>{\raggedright\arraybackslash}p{0.198\textwidth}>{\raggedright\arraybackslash}p{0.285\textwidth}>{\raggedright\arraybackslash}p{0.433\textwidth}@{}}
\hline
\textbf{Flaw size \emph{a}} & \textbf{Fracture stress \ensuremath{\sigma}\textsubscript{f}} & \textbf{Margin (\ensuremath{\sigma}\textsubscript{f} / 0.5--1.1 GPa realistic)} \\ \hline
5 \ensuremath{\mu}m & 0.50 GPa & \textasciitilde{}1 \\
\hline
2 \ensuremath{\mu}m & 0.80 GPa & \textasciitilde{}1 \\
\hline
1 \ensuremath{\mu}m & 1.13 GPa & \textasciitilde{}1 \\
\hline
0.2 \ensuremath{\mu}m & 2.52 GPa & \textasciitilde{}2--5 \\
\hline
\end{tabular}
\endgroup



\texttt{[DERIVED]}. For realistic secondary-particle flaw sizes of 1--5 \ensuremath{\mu}m, the fracture stress is comparable to the operating stress, a margin of order 1. Not 10\ensuremath{^{2}} (safe by a wide margin) and not 10\ensuremath{^{-}}\ensuremath{^{2}} (certain to fail), but \emph{at the threshold}. This is the \textquotedblleft{}margins are order unity\textquotedblright{} result, and it is why particle size is a design variable with a real optimum rather than \textquotedblleft{}smaller is better\textquotedblright{}: smaller particles resist fracture (higher \ensuremath{\sigma}\textsubscript{f}) but become surface-dominated and lose energy density.


The two errors this derivation is designed to prevent.


\begin{itemize}
  \item Quoting 12 GPa as the stress. It is an elastic upper bound, not an operating stress, and quoting it would make fracture seem certain. The correct comparison is the 0.5--2 GPa realistic range.
  \item Quoting a margin of 10\ensuremath{^{2}} or 10\ensuremath{^{-}}\ensuremath{^{2}}. The Loop-C margin is \textasciitilde{}1, which is what makes the coupling live rather than safely inactive.
\end{itemize}


\textbf{Uncertainty.} \emph{K}\textsubscript{Ic} = 2 MPa m\textsuperscript{\ensuremath{\tfrac{1}{2}}} for layered oxides is uncertain by \textasciitilde{}2\ensuremath{\times}; flaw sizes span 0.2--5 \ensuremath{\mu}m (a 25\ensuremath{\times} range in \emph{a}, hence 5\ensuremath{\times} in \ensuremath{\sigma}\textsubscript{f}). The margin is therefore \textquotedblleft{}order 1\textquotedblright{} with a spread of roughly 0.5--5, which is the correct granularity for the claim, and stating it that way is more defensible than any single number.


\par\noindent\rule{\linewidth}{0.4pt}\par


\subsection*{S2.8 The dissolution unit ambiguity: a worked example}


Most reports of transition-metal dissolution quote a concentration in ppm, which is a \emph{dose into the electrolyte}, while the coupling argument needs a \emph{fraction of the cathode inventory}. The two are not interchangeable, and the conversion depends on the cell format. This section works the conversion out, as a cautionary example of how the same physical loss can be reported as two numbers that differ by nearly an order of magnitude.


\textbf{The arithmetic.} For an NMC811 cathode (\emph{M} = 97.3 g mol\ensuremath{^{-}}\ensuremath{^{1}}) at 200 mAh g\ensuremath{^{-}}\ensuremath{^{1}} and 1 mAh cm\ensuremath{^{-}}\ensuremath{^{2}}:


\begin{equation}m_{\text{cathode}} = \frac{1\ \text{mAh cm}^{-2}}{200\ \text{mAh g}^{-1}} = 5.0\ \text{mg cm}^{-2} \nonumber\end{equation}


\begin{equation}m_{\text{Mn}} = 5.0\times10^{-3}\times\frac{0.1\times54.94}{97.3} = \mathbf{282\ \mu g\ cm^{-2}} \tag{S2.24}\end{equation}


\texttt{[DERIVED]}. A 1 \% Mn loss is therefore 2.82 \ensuremath{\mu}g cm\ensuremath{^{-}}\ensuremath{^{2}} = 0.051 \ensuremath{\mu}mol cm\ensuremath{^{-}}\ensuremath{^{2}}.


\textbf{Now express that loss as a concentration --- in two different cells:}


\begingroup\small
\setlength{\tabcolsep}{3pt}
\renewcommand{\arraystretch}{1.12}
\sloppy
\begin{tabular}{@{}>{\raggedright\arraybackslash}p{0.172\textwidth}>{\raggedright\arraybackslash}p{0.222\textwidth}>{\raggedright\arraybackslash}p{0.157\textwidth}>{\raggedright\arraybackslash}p{0.353\textwidth}@{}}
\hline
\textbf{Cell format} & \textbf{Electrolyte} & \textbf{1 \% Mn loss \ensuremath{\rightarrow}} & \textbf{100 ppm dose as fraction of Mn inventory} \\ \hline
Lean coin & 6 \ensuremath{\mu}L \ensuremath{\approx} 7.2 mg & \ensuremath{\approx} 390 ppm & \ensuremath{\approx} 0.25 \% \\
\hline
Flooded & 50 \ensuremath{\mu}L \ensuremath{\approx} 60 mg & \ensuremath{\approx} 47 ppm & \ensuremath{\approx} 2.1 \% \\
\hline
\end{tabular}
\endgroup



\texttt{[DERIVED]}. Arithmetic for the first row: 2.82 \ensuremath{\mu}g / 7.2 mg = 2.82\ensuremath{\times}10\ensuremath{^{-}}\ensuremath{^{6}} g / 7.2\ensuremath{\times}10\ensuremath{^{-}}\ensuremath{^{3}} g = 392 ppm. For the second: 2.82 \ensuremath{\mu}g / 60 mg = 47 ppm.


\textbf{Two implications follow:}


\begin{enumerate}
  \item The same 1 \% Mn loss reads as 390 ppm or 47 ppm, an 8.3\ensuremath{\times} ambiguity, purely from cell design. Dissolution published as ppm-in-electrolyte without electrolyte volume and cathode loading is \textbf{not comparable across papers.} This is the reporting gap the checklist (S3) is built to close; it is quantified rather than asserted.
  \item A \textquotedblleft{}100 ppm\textquotedblright{} crosstalk dose means 0.25 \% of the Mn inventory in a lean cell but 2.1 \% in a flooded one, an 8\ensuremath{\times} range in what the same nominal dose delivers. The consequence is that the standard crosstalk experiment is not a gentle probe; in a flooded cell it is a large compressed insult. And the direction of the inference matters: such an experiment measures the cell's sensitivity to a large dose, not the rate at which dissolution occurs in an unperturbed cell. Anyone using a 100 ppm crosstalk result to infer a dissolution rate is reading the experiment backwards.
\end{enumerate}


A claim this section explicitly does NOT make. It does not convert the 100 ppm dose into \textquotedblleft{}cycles' worth of dissolution\textquotedblleft{} at a stated ppm/cycle rate, because that conversion depends on an unstated rate and is not what the source measurement supports. The dose's significance here is as a concentration, not as a time-compressed equivalent, and the SI keeps it that way.


\par\noindent\rule{\linewidth}{0.4pt}\par


\subsection*{S2.9 Voltage fade is arithmetic-separable from capacity fade}


A fourth accounting term, thermodynamic/potential loss (TL), is adopted for energy fade, and it is claimed to be separable from the capacity channels. That claim requires evidence, and the evidence is that the two can move independently by large amounts.


\textbf{The measurement.} In Li-rich/Mn-rich layered oxides after 200 cycles, the midpoint-voltage drop is \textbf{723.5--769.6 mV} while capacity retention remains high, and one sample shows a 491 mV shift. Mean capacity fade over the same interval is far smaller than the energy-density loss those voltage numbers imply, which reaches roughly 30 \% by the arithmetic below. \texttt{[FROM LITERATU}\hspace{0pt}\texttt{RE: 10.1039/}\hspace{0pt}\texttt{D4SC07037B]}.


The arithmetic that makes it a separability result. Energy density is the product of capacity and average voltage. A 200-cycle capacity retention of, say, 90 \% combined with a 770 mV midpoint drop from a \textasciitilde{}3.7 V average operating voltage gives an energy retention of 0.90 \ensuremath{\times} (3.7 \ensuremath{-} 0.77)/3.7 = 0.90 \ensuremath{\times} 0.79 = \textbf{71 \%}, i.e. an additional 19 percentage points of energy loss that the capacity balance does not capture. \texttt{[DERIVED]} from the stated retention and voltage figures; illustrative capacity retention value, since the relevant quantity is the order of magnitude.


Why this forces a fourth term rather than a fudge. Folding a 19 pp energy loss into the kinetic (impedance) channel would double-count it against the EIS measurement in S4 branch B5, which measures a resistive loss that is physically distinct from a thermodynamic potential shift. The two are measured by different instruments and have different temperature dependences. So the LLI/LAM/KL triple is complete for capacity and incomplete for energy, and TL is required, with the explicit scope statement that TL is energy-only and does not enter the capacity balance. \texttt{[HYPOTHESIS]} for the ontology (that TL is the right fourth term); \texttt{[DERIVED]} for the arithmetic showing the two fades are quantitatively separable.


\textbf{The scope limitation.} Three of the four loops (A, B, C) operate on the capacity balance; TL is outside the loop framework entirely. The loop framework is therefore restricted to the capacity channels, and voltage fade remains a named open problem; the four loops are not a complete account. S2.9 and S4's B6 branch both carry this caveat.


\par\noindent\rule{\linewidth}{0.4pt}\par


\subsection*{S2.10 Summary of the ratios}


\begingroup\small
\setlength{\tabcolsep}{3pt}
\renewcommand{\arraystretch}{1.12}
\sloppy
\begin{longtable}{@{}>{\raggedright\arraybackslash}p{0.258\textwidth}>{\raggedright\arraybackslash}p{0.134\textwidth}>{\raggedright\arraybackslash}p{0.282\textwidth}>{\raggedright\arraybackslash}p{0.231\textwidth}@{}}
\hline
\textbf{Quantity} & \textbf{Derivation} & \textbf{Value} & \textbf{Uncertainty} \\ \hline
\endhead
SEI ohmic drop at 1C, \ensuremath{\sigma} = 10\ensuremath{^{-}}\ensuremath{^{7}} S cm\ensuremath{^{-}}\ensuremath{^{1}} & S2.1 & \textasciitilde{}90 mV & \ensuremath{\pm}2 orders (\ensuremath{\sigma}) \\
SEI-drop / anode-polarization cross-check & S2.1 & 0.9 V would exceed measured 50--150 mV & qualitative, robust \\
Plating threshold \emph{R}* at 1C & S2.2 & 20.0 \ensuremath{\Omega} cm\ensuremath{^{2}} & \ensuremath{\pm}(\ensuremath{\Delta}, \emph{\ensuremath{\eta}}\textsubscript{ct} definition) \\
Critical interphase thickness \emph{d}* at 1C & S2.2 & 20 / 2 / 0.2 nm (\ensuremath{\sigma} = 10\ensuremath{^{-}}\ensuremath{^{7}}/10\ensuremath{^{-}}\ensuremath{^{8}}/10\ensuremath{^{-}}\ensuremath{^{9}}) & \ensuremath{\pm}2 orders \\
SEI-equivalent thickness per mAh cm\ensuremath{^{-}}\ensuremath{^{2}} & S2.3 & \textasciitilde{}11.2 \ensuremath{\mu}m (30 cm\ensuremath{^{3}} mol\ensuremath{^{-}}\ensuremath{^{1}}); 12.9 \ensuremath{\mu}m (Li\ensuremath{_{2}}CO\ensuremath{_{3}}) & \textasciitilde{}\ensuremath{\pm}30 \% (\emph{V}\textsubscript{m}) \\
Pore-filling LLI threshold & S2.3 & 0.27 mAh cm\ensuremath{^{-}}\ensuremath{^{2}} dense; 0.3--0.6 (10--20 \%) with porosity & \ensuremath{\pm}2\ensuremath{\times} (pore volume; 2 conventions, S2.3) \\
Model LLI crosses pore-fill threshold & S2.3 & cycle 386, 144 cycles before the knee at 530 (LLI 0.27 vs 0.6 mAh cm\ensuremath{^{-}}\ensuremath{^{2}}, 2.22\ensuremath{\times}) & inherits S1 parameters \\
Li inventory thickness at 3 mAh cm\ensuremath{^{-}}\ensuremath{^{2}} & S2.4 & 14.55 \ensuremath{\mu}m & exact \\
CE required for 1000 cycles at 80 \% & S2.4 & 99.978 \% & law-independent to 0.002 pp \\
CE-law ambiguity in exhaustion life & S2.4 & 4.60\ensuremath{\times} & exact \\
Dead-Li versus SEI discriminator & S2.4 & 2.27\ensuremath{\times} (not 22\ensuremath{\times}) & relies on the 11 \ensuremath{\mu}m/mAh molar-volume basis \\
Loop B shortfall (catalytic requirement) & S2.5 & \textasciitilde{}74\ensuremath{\times} & \ensuremath{\pm}2\ensuremath{\times} (dissolution rate) \\
Loop D Li\ensuremath{^{+}} salt budget & S2.6 & \textasciitilde{}9\ensuremath{\times} overdrawn & \ensuremath{\pm}3\ensuremath{\times} (stoichiometry) \\
Loop D format dependence & S2.6 & 15\ensuremath{\times} (flooded vs lean) & exact \\
Loop C \ensuremath{\Delta}\emph{c}\textsubscript{max} & S2.7 & 0.0515 (adopt 0.048) mol cm\ensuremath{^{-}}\ensuremath{^{3}} & \ensuremath{\pm}5 \% \\
Loop C elastic upper bound & S2.7 & 12.9 GPa (quoted 12) & ceiling only \\
Loop C realistic stress vs fracture stress & S2.7 & 0.5--2 GPa vs 0.5--1.1 GPa \ensuremath{\rightarrow} margin \textasciitilde{}1 & \ensuremath{\pm}2\ensuremath{\times} (\emph{K}\textsubscript{Ic}), 0.5--5 \\
M3 ppm ambiguity & S2.8 & 390 vs 47 ppm (8.3\ensuremath{\times}; the 8.4\ensuremath{\times} below is the inventory fraction, a different ratio) & exact \\
M3 100 ppm dose as inventory fraction & S2.8 & 0.25 \% vs 2.1 \% (8.4\ensuremath{\times}) & exact \\
Voltage-fade energy gap after 200 cycles & S2.9 & 723.5--769.6 mV \ensuremath{\rightarrow} \textasciitilde{}19 pp energy & \ensuremath{\pm}(illustrative retention) \\
\hline
\end{longtable}
\endgroup



\textbf{The claim these support.} Every ratio in the table that is a \emph{coupling margin}, that is, SEI drop against plating margin, pore fill against capacity, and the Loop B, Loop D and Loop C margins, lands between roughly 1 and 100. None is of order 10\ensuremath{^{-}}\ensuremath{^{2}} and none is of order 10\ensuremath{^{4}}. That clustering is the central quantitative observation: the margins are all of order unity, which is why the mechanisms are simultaneously active and why modeling them as independent terms is not a small approximation. It is also why design has implicitly optimized toward several thresholds at once --- enough SEI to passivate but not enough to fill pores, particles small enough to resist fracture but not so small as to be surface-dominated, electrolyte enough to last but not so much as to be a penalty. Coupled degradation is the price of that optimization.


\textbf{What would falsify the claim.} A single coupling whose margin is 10\ensuremath{^{-}}\ensuremath{^{3}} (irrelevant) or 10\ensuremath{^{4}} (overwhelming) would break the symmetry of the argument. The Loop-B \textasciitilde{}74\ensuremath{\times} is the largest ratio in the table and it is the weakest link: if dissolution rates were \textasciitilde{}200 ppm/cycle the required turnover would collapse to \textasciitilde{}1.8 and Loop B would be near-stoichiometric, which would weaken but not destroy the framework (Loops A, C, D are independent of it). This limit is stated explicitly because it bounds the framework.


\section*{S3. Table S1, the mechanism \ensuremath{\times} measurement matrix, and the reporting checklist}


\subsection*{S3.1 Structure of the matrix}


A capacity-fade curve is a single scalar that integrates every loss mechanism at once. It cannot, on its own, answer all of the three questions a degradation paper addresses:


\begin{enumerate}
  \item Loss-term identity. Which of Li inventory, active material, and kinetics moved (a quantity question).
  \item Electrode attribution. Which electrode is responsible (a location question).
  \item Mechanism family. Which mechanism family is responsible (a chemistry question).
\end{enumerate}


The matrix below maps each measurement onto the questions above. It is required because a fade curve reported as evidence for a mechanism typically answers none of them, and because the order in which measurements are applied may matter more than the measurements themselves. Every technique in Table S1 is established; none is novel. Indexing them by question and by loop is what makes the table a diagnostic plan rather than a list.


\textbf{Verdict legend.} \ensuremath{\checkmark} clean separation \ensuremath{\cdot} \ensuremath{\bullet} partial, or requires a differential experiment \ensuremath{\cdot} \ensuremath{\boldsymbol{\times}} cannot separate. A \ensuremath{\boldsymbol{\times}} is information: it tells the reader which claims this measurement cannot support, an element such lists usually omit.


\textbf{Loop column.} Which of the four feedback loops (A: plating/dead Li/SEI; B: dissolution/crossover; C: fracture/stress; D: electrolyte/salt inventory) the measurement is diagnostic of. A blank means the measurement constrains the balance but not a specific loop.


\par\noindent\rule{\linewidth}{0.4pt}\par


\subsection*{S3.2 Table S1. Mechanism \ensuremath{\times} measurement matrix}


\begingroup\scriptsize
\setlength{\tabcolsep}{3pt}
\renewcommand{\arraystretch}{1.12}
\sloppy
\begin{longtable}{@{}>{\raggedright\arraybackslash}p{0.035\textwidth}>{\raggedright\arraybackslash}p{0.180\textwidth}>{\raggedright\arraybackslash}p{0.214\textwidth}>{\raggedright\arraybackslash}p{0.217\textwidth}>{\raggedright\arraybackslash}p{0.134\textwidth}>{\raggedright\arraybackslash}p{0.036\textwidth}>{\raggedright\arraybackslash}p{0.053\textwidth}@{}}
\hline
\textbf{\#} & \textbf{Measurement} & \textbf{What it separates} & \textbf{What it \emph{cannot} separate} & \textbf{Verdict} & \textbf{Q} & \textbf{Loop} \\ \hline
\endhead
1 & dQ/dV, dV/dQ, incremental capacity (ICA/DVA) & Li-inventory loss (peak \emph{shifts}) vs active-material loss (peak \emph{area} shrinks); the canonical LLI/LAM decomposition & Cannot distinguish \emph{where} Li was lost (SEI vs dead Li vs isolated pocket), nor \emph{which} electrode lost AM & \ensuremath{\bullet} for Q1; \ensuremath{\boldsymbol{\times}} for Q2, Q3 & 1 & --- \\
2 & High-precision coulometry (HPC), CE titration & Parasitic current per cycle to \textasciitilde{}0.001 \% precision; reversible vs irreversible inefficiency; the most sensitive early-warning signal & Cannot say whether the parasitic current is SEI (M1) or dead-Li corrosion (M2) --- both appear as CE < 1 & \ensuremath{\checkmark} for Q1 magnitude; \ensuremath{\boldsymbol{\times}} for Q3 & 1 & A \\
3 & Voltage-relaxation profile (VRP) / OCV relaxation & Reversible plating specifically: a voltage plateau during relaxation is the signature of stripping Li\ensuremath{^{0}} & Cannot quantify dead (isolated) Li; a plateau's absence does not mean plating did not occur irreversibly & \ensuremath{\checkmark} for \emph{reversible} M2 only & 1, 3 & A \\
4 & EIS / distribution of relaxation times (DRT) & Electrode-resolved time constants: SEI vs charge-transfer vs diffusion; DRT resolves overlapping arcs without equivalent-circuit guessing & Attribution of a specific arc to M3 vs M1 is model-dependent; reconstruction-layer and CEI impedance overlap in frequency & \ensuremath{\bullet} for Q2; \ensuremath{\boldsymbol{\times}} for Q3 & 2 & A, B, D \\
5 & Three-electrode / multi-reference cell & Directly resolves anode from cathode overpotential and impedance, the cleanest answer to Q2 & Destructive to the cell format; RE placement artefacts; not applicable to commercial cells & \ensuremath{\checkmark} for Q2 & 2 & --- \\
6 & Half-cell reconstruction (coin-cell rebuild) & Rebuilds each electrode against fresh Li counter, separating anode-limited from cathode-limited fade after the fact; the standard post-hoc complement to dV/dQ & Destructive; the rebuild itself perturbs the electrode (new SEI, new pressure, air exposure); cannot distinguish mechanisms within an electrode & \ensuremath{\checkmark} for Q2 & 2 & --- \\
7 & Operando XRD / neutron diffraction & Lattice-parameter evolution, phase fractions (spinel/rock-salt), Li staging in graphite; distinguishes \emph{structural} AM loss from loss of electrochemically active Li & Does not see amorphous or surface-confined species; poor at the few-nm SRL scale & \ensuremath{\checkmark} for Q1 (structural), \ensuremath{\bullet} for Q3 & 1, 3 & C, B \\
8 & Neutron depth profiling / neutron imaging & Spatially resolved Li density, including inactive Li; bulk-sensitive and non-destructive & Modest spatial resolution; not element-specific for TM & \ensuremath{\bullet} & 1, 2 & --- \\
9 & ICP-MS / ICP-OES (electrolyte or anode digest) & Total dissolved TM, element-resolved, the \emph{only} direct proof of M3 mass loss & Only total amount, not rate or timing; destructive & \ensuremath{\checkmark} for Q3 (M3 vs not-M3); \ensuremath{\boldsymbol{\times}} for timing & 3 & B \\
10 & LA-ICP-OES / \ensuremath{\mu}XRF mapping & Spatially resolved TM deposition on anode and Li distribution & Semi-quantitative; chemical state invisible & \ensuremath{\checkmark} for Q3 location & 2, 3 & B \\
11 & XPS (anode, cathode) & Chemical state of deposited TM, SEI/CEI composition, F 1s speciation & Surface-only; quantification unreliable; air-sensitive & \ensuremath{\bullet} & 3 & B, C \\
12 & DEMS (O\ensuremath{_{2}}, CO\ensuremath{_{2}}, CO) & Lattice oxygen loss (the M3 trigger), solvent oxidation; time- and potential-resolved & O\ensuremath{_{2}} partly consumed by solvent before detection, so O\ensuremath{_{2}} yield \ensuremath{\neq} lattice O lost; cannot identify which cation left & \ensuremath{\checkmark} for the oxygen channel; \ensuremath{\boldsymbol{\times}} for TM & 3 & B \\
13 & Operando \ensuremath{^{1}}\ensuremath{^{9}}F NMR + EPR & HF and Mn\ensuremath{^{2}}\ensuremath{^{+}} concentration \emph{simultaneously in time}, the best existing probe of the acid--dissolution loop & Requires specialized hardware; not for commercial sealed cells & \ensuremath{\checkmark} for Q3 (M3 chemistry) & 3 & B \\
14 & Isothermal microcalorimetry / electrochemical calorimetry & Parasitic heat flux, an independent, non-electrical measure of parasitic current, so it separates \emph{chemical} from \emph{transport} inefficiency & Cannot assign heat to a specific mechanism without complementary data & \ensuremath{\bullet} & 1 & --- \\
15 & Stack-pressure / dilatometry & Irreversible vs reversible expansion; mechanical coupling (M4), and via pressure the M4\ensuremath{\rightarrow}M3 arrow & Indirect; requires pressure-controlled hardware & \ensuremath{\bullet} & 2, 3 & C \\
16 & ARC (accelerating rate calorimetry) & Thermal-onset signature; dead-Li and plating detection via rest-time-dependent onset & A safety probe, not a mechanism separator; destructive & \ensuremath{\boldsymbol{\times}} for Q1--Q3; \ensuremath{\checkmark} for safety context & --- & A \\
17 & Isotope labeling (\ensuremath{^{6}}Li/\ensuremath{^{7}}Li, \ensuremath{^{1}}\ensuremath{^{8}}O) & Provenance tracking, where Li and O actually went; the only method distinguishing \emph{native} from \emph{re-deposited} species & Expensive, low-throughput, rarely applied to full cells & \ensuremath{\checkmark} for provenance, but underused & 3 & A, B \\
18 & Titration gas chromatography (TGC) & Separates unreacted metallic Li\ensuremath{^{0}} (dead Li) from Li\ensuremath{^{+}} in the SEI, the definitive M1-vs-M2 discriminator & Destructive; ex situ; requires cell teardown; no Li\ensuremath{^{0}} to titrate in graphite/Li-ion cells & \ensuremath{\checkmark} --- see S3.3 & 3 & A \\
19 & \ensuremath{^{7}}Li in situ / operando NMR & Dead Li vs SEI Li vs bulk Li, non-destructively; also Li site occupancy for voltage fade & Requires NMR-accessible cell; overlapping shifts need modeling & \ensuremath{\checkmark} & 1, 3 & A \\
20 & Cryo-EM / STEM-EELS & Nanoscale morphology and phase of dead Li and of the SRL & Local only; no statistics; extreme care needed & \ensuremath{\bullet} & 3 & A, C \\
\hline
\end{longtable}
\endgroup



\emph{Notes.} [a] Row 6 (half-cell reconstruction) is included to reach 20 rows: it answers question 2, locating the loss on the anode or the cathode, and is the standard complement to row 1 in the literature. [b] Q = the question from S3.1 that the measurement primarily answers. [c] \textquotedblleft{}Loop\textquotedblright{} is the assignment adopted here, not one taken from the source literature; where a measurement constrains two loops, both are listed. [d] The six rows that carry Loop A (2, 3, 4, 18, 19, 20) are exactly the ones a Loop-A test plan needs; there is no single measurement that resolves Loop A, which is the matrix's principal negative finding.


What the matrix shows that a list would not. Score every row on question 3 (chemical identity --- \emph{which mechanism family}), which is the question the field most wants answered:


\begingroup\small
\setlength{\tabcolsep}{3pt}
\renewcommand{\arraystretch}{1.12}
\sloppy
\begin{tabular}{@{}>{\raggedright\arraybackslash}p{0.544\textwidth}>{\raggedright\arraybackslash}p{0.125\textwidth}>{\raggedright\arraybackslash}p{0.247\textwidth}@{}}
\hline
\textbf{Chemistry capability} & \textbf{Count} & \textbf{Rows} \\ \hline
\ensuremath{\checkmark} clean --- directly identifies chemical species or state & 6 & 9, 10, 13, 17, 18, 19 \\
\hline
\ensuremath{\bullet} partial --- indirect, or only for one sub-case & 5 & 3, 7, 11, 12, 20 \\
\hline
\ensuremath{\boldsymbol{\times}} cannot --- explicitly fails on chemistry & 4 & 1, 2, 4, 16 \\
\hline
, silent, answers Q1/Q2 only, says nothing about chemistry & 5 & 5, 6, 8, 14, 15 \\
\hline
\end{tabular}
\endgroup



So only 6 of 20 techniques cleanly answer the chemistry question, and 9 either only partly answer it or cannot answer it at all. That asymmetry is the empirical case for the decision tree in S4: because no single measurement settles the chemistry question, the productive strategy is ordering measurements so each eliminates a branch, rather than seeking one definitive probe. The matrix provides the evidence and the tree the response.


\par\noindent\rule{\linewidth}{0.4pt}\par


\subsection*{S3.3 The best single diagnostic for SEI versus dead Li}


\textbf{TGC, per Fang et al.} \texttt{[FROM LITERATU}\hspace{0pt}\texttt{RE: 10.1038/}\hspace{0pt}\texttt{s41586-019-}\hspace{0pt}\texttt{1481-z]}, reinforced by operando \ensuremath{^{7}}Li NMR when the cell cannot be destroyed.


The chemistry is the reason it works and should be stated rather than cited as a black box: unreacted metallic Li\ensuremath{^{0}} reacts with a protic reagent (water or an alcohol) to evolve H\ensuremath{_{2}} quantitatively, while Li\ensuremath{^{+}} bound in the SEI (Li\ensuremath{_{2}}CO\ensuremath{_{3}}, LiF, ROCO\ensuremath{_{2}}Li, Li\ensuremath{_{2}}O) does not. Measuring evolved H\ensuremath{_{2}} by GC therefore gives a direct, calibration-free count of Li\ensuremath{^{0}} atoms; subtracting that from total inactive Li yields the SEI-bound Li\ensuremath{^{+}} fraction.


The result is central: in Li-metal cells, unreacted metallic Li\ensuremath{^{0}}, not SEI-bound Li\ensuremath{^{+}}, dominates inactive lithium. That re-partitions the two most-studied mechanism families. The SEI is not primarily a lithium sink but primarily a \emph{transport barrier}, and dead lithium, not interphase, is where most of the stranded inventory sits in these cells.


The two families remain parallel members of the loop rather than one being a downstream product of the other: Loop A closes through M2 \ensuremath{\rightarrow} M1 (dead lithium carries fresh interphase) and through M1 \ensuremath{\rightarrow} M2 (interphase raises the overpotential that drives plating). What TGC settles is which family dominates in a given cell, not the direction of the arrow between them.


\textbf{Practical caveats to state.} TGC requires teardown and is post-mortem, not operando. It is not applicable to graphite/Li-ion cells, where all Li is either intercalated or in the SEI --- there is no Li\ensuremath{^{0}} to titrate. For those cells the equivalent discriminator is high-precision coulometry plus three-electrode reconstruction, which localizes the loss to the anode but does not split SEI from dead Li. That residual gap is stated in S4.11 item 1 and is the most significant acknowledged limit of the tree.


\par\noindent\rule{\linewidth}{0.4pt}\par


\subsection*{S3.4 The 20-item reporting checklist}


The biggest reporting gap in this literature is not a missing technique --- it is that results are not comparable across papers, because the quantities that set a mechanism's rate are not reported. Dissolution in \textquotedblleft{}ppm\textquotedblright{} and fade in \textquotedblleft{}\% per 100 cycles\textquotedblright{} are not physical quantities until the cell's normalization is known; S2.8 quantifies that ambiguity at 8.3\ensuremath{\times} for a single cell-design change. The checklist provides the response, in three tiers.


\subsubsection*{Tier 1: required for an interpretable number}


\begingroup\small
\setlength{\tabcolsep}{3pt}
\renewcommand{\arraystretch}{1.12}
\sloppy
\begin{tabular}{@{}>{\raggedright\arraybackslash}p{0.031\textwidth}>{\raggedright\arraybackslash}p{0.533\textwidth}>{\raggedright\arraybackslash}p{0.352\textwidth}@{}}
\hline
\textbf{\#} & \textbf{Item} & \textbf{Why it is Tier 1} \\ \hline
1 & Cathode: composition, morphology (polycrystalline vs single-crystal), particle-size distribution, loading (mg cm\ensuremath{^{-}}\ensuremath{^{2}}), density/porosity, areal capacity (mAh cm\ensuremath{^{-}}\ensuremath{^{2}}) & Sets M3 and M4 susceptibility; without areal capacity no current density is defined \\
\hline
2 & Anode: type, loading, porosity, areal capacity & Sets N/P and which mechanism's loss is masked \\
\hline
3 & N/P ratio & Determines whether the cell is anode- or cathode-limited, hence which loss is visible \\
\hline
4 & Electrolyte: salt, concentration, solvent ratio, additive identity and wt\%, total mass or volume per cell, and electrolyte-to-active-material ratio & Without this, ppm values are uninterpretable, the 8.3\ensuremath{\times} ambiguity of S2.8 \\
\hline
5 & Water and HF content of the as-assembled cell, with method (coulometric Karl Fischer / acid--base titration) & The solvent spec is not the cell content; both drive Loop B \\
\hline
6 & Cell format and stack pressure & Coin vs pouch vs cylindrical changes pressure history; pressure changes both M4 cracking and M2 plating \\
\hline
7 & Voltage window, C-rate (charge and discharge), temperature, aging protocol (CC-CV vs CC, rest periods, depth of discharge) & Sets \emph{R}* (S2.2); without it a knee position is not comparable \\
\hline
8 & Cycle number at which fade is quoted, plus the full trajectory & A single endpoint hides the knee, which is the object of study \\
\hline
9 & Formation protocol & First-cycle irreversible loss is often a large fraction of total \textquotedblleft{}degradation\textquotedblright{} and is protocol-dependent \\
\hline
\end{tabular}
\endgroup



\subsubsection*{Tier 2: required when a mechanism is claimed}


\begingroup\small
\setlength{\tabcolsep}{3pt}
\renewcommand{\arraystretch}{1.12}
\sloppy
\begin{tabular}{@{}>{\raggedright\arraybackslash}p{0.050\textwidth}>{\raggedright\arraybackslash}p{0.496\textwidth}>{\raggedright\arraybackslash}p{0.370\textwidth}@{}}
\hline
\textbf{\#} & \textbf{Item} & \textbf{Why it is Tier 2} \\ \hline
10 & Loss decomposition (LLI / LAM / polarization) with the method used & \textquotedblleft{}Capacity fade\textquotedblright{} alone does not answer question 1 \\
\hline
11 & Coulombic efficiency at \ensuremath{\geq}4 significant figures, SOC window, number of cycles averaged & CE below 4 s.f. cannot discriminate the Loop-A rates of S2.4 \\
\hline
12 & Reference-cycle protocol before/after aging (rate, window, temperature) & Otherwise fade is confounded with rate-capability loss \\
\hline
13 & Impedance spectra with the equivalent circuit or DRT used, and the frequency range & An arc without its deconvolution is not an assignment \\
\hline
14 & Post-mortem sampling plan: electrode, location (edge vs center), number of samples, how the cell was opened & Air exposure is the classic artefact in TM analysis \\
\hline
15 & Detection limits for ICP-MS/XPS/DEMS & \textquotedblleft{}Below detection\textquotedblright{} is a result only if the limit is stated \\
\hline
\end{tabular}
\endgroup



\subsubsection*{Tier 3: recommended for reuse}


\begingroup\small
\setlength{\tabcolsep}{3pt}
\renewcommand{\arraystretch}{1.12}
\sloppy
\begin{tabular}{@{}>{\raggedright\arraybackslash}p{0.045\textwidth}>{\raggedright\arraybackslash}p{0.431\textwidth}>{\raggedright\arraybackslash}p{0.440\textwidth}@{}}
\hline
\textbf{\#} & \textbf{Item} & \textbf{Why it is Tier 3} \\ \hline
16 & Raw or minimally processed data for reference cycles and impedance spectra (supplementary) & Lets others re-run the decomposition \\
\hline
17 & Cell-level metadata for physics-based modeling: separator thickness, current-collector foils, tab configuration & Needed to reproduce any transport model \\
\hline
18 & The claim-to-measurement mapping: explicitly, which measurement supports which mechanistic claim & The item most often skipped, and the one that would most improve the field --- it forces authors to notice when a claim rests on a fade curve alone \\
\hline
19 & Negative results: mechanisms ruled out, and by what measurement & The tree's value depends on eliminations being reported \\
\hline
20 & Whether fade was path-dependent (rate history, rest, temperature) & Path dependence is the normal case; a single aging protocol underdetermines the mechanism \\
\hline
\end{tabular}
\endgroup



\subsubsection*{Suggested template sentence for a methods section}


\begin{quote}\textquotedblleft{}Cell build: [cathode composition, morphology, loading mg cm\ensuremath{^{-}}\ensuremath{^{2}}, areal capacity mAh cm\ensuremath{^{-}}\ensuremath{^{2}}] \ensuremath{\|} [anode, loading, areal capacity], N/P = [x], electrolyte [salt, conc., solvent ratio, additives wt\%], electrolyte mass [mg] (E/AM = [x] g Ah\ensuremath{^{-}}\ensuremath{^{1}}), separator [type, thickness], cell format [coin/pouch], stack pressure [MPa]. Water [ppm, method]; HF [mg L\ensuremath{^{-}}\ensuremath{^{1}}, method]. Formation: [protocol]. Aging: [protocol, voltage window, C-rate, T]. Reference cycle: [C/20, window, T]. Loss decomposition by [method]; CE precision [\ensuremath{\pm} x \%].\textquotedblleft{}\end{quote}


The checklist is applied to the present work as well as to others. A methods paragraph that omits items 4, 5, or 11 does not meet the standard the checklist sets.


\par\noindent\rule{\linewidth}{0.4pt}\par


\subsection*{S3.5 Coverage of the four loops}


\begingroup\small
\setlength{\tabcolsep}{3pt}
\renewcommand{\arraystretch}{1.12}
\sloppy
\begin{tabular}{@{}>{\raggedright\arraybackslash}p{0.246\textwidth}>{\raggedright\arraybackslash}p{0.358\textwidth}>{\raggedright\arraybackslash}p{0.312\textwidth}@{}}
\hline
\textbf{Loop} & \textbf{Diagnostic measurements (row numbers)} & \textbf{Coverage} \\ \hline
A --- plating \ensuremath{\rightarrow} dead Li \ensuremath{\rightarrow} SEI \ensuremath{\rightarrow} \ensuremath{\eta} \ensuremath{\rightarrow} plating & 2 (HPC/CE), 3 (VRP), 4 (EIS/DRT), 18 (TGC), 19 (\ensuremath{^{7}}Li NMR), 20 (cryo-EM), 16 (ARC, safety) & Strong on \emph{detecting} plating; weak on the feedback arrow (whether dead Li actually raises \emph{R}) \\
\hline
B --- dissolution \ensuremath{\rightarrow} anode poisoning \ensuremath{\rightarrow} \ensuremath{\eta} \ensuremath{\rightarrow} dissolution & 4, 9 (ICP-MS), 10 (LA-ICP), 11 (XPS), 12 (DEMS), 13 (\ensuremath{^{1}}\ensuremath{^{9}}F NMR/EPR), 17 (isotopes) & Strong on \emph{detecting} dissolution; weak on turnover --- no routine measurement gives the catalytic turnover number S2.5 requires \\
\hline
C --- fracture \ensuremath{\rightarrow} fresh surface \ensuremath{\rightarrow} SEI/dissolution & 4, 7 (XRD), 11, 15 (pressure/dilatometry), 20 & Adequate for \emph{detecting} fracture; fresh-surface quantification is indirect only \\
\hline
D --- electrolyte/salt depletion & 4, plus a flooding differential (not a row --- it is an \emph{experiment design}, not a technique) & Weakest loop. No matrix row measures salt inventory directly; Loop D is diagnosed by elimination and by the flooding test \\
\hline
\end{tabular}
\endgroup



\textbf{The gap this exposes.} Table S1 is most complete exactly where the argument is least constrained. The matrix detects mechanisms well (plating, dissolution, fracture) but detects coupling, the operative claim --- poorly. No row measures \textquotedblleft{}the arrow from dead Li to interphase resistance,\textquotedblright{} and no row measures the catalytic turnover of Loop B or the salt inventory of Loop D. That is why the decision tree (S4) terminates in B7, loop identification, which is not a measurement but a \emph{perturbation protocol}: it infers coupling from the sensitivity of the knee to a driving variable, because coupling itself is not directly measurable with current techniques. The checklist's item 18 (claim-to-measurement mapping) and item 19 (negative results) are what make that inference auditable.


A consequence for the scope of the claim. Table S1 is not a complete diagnostic toolkit. It provides a sequencing of established techniques that resolves mechanism, plus one perturbation protocol that tests coupling. A stronger claim would exceed what the matrix supports, as its \ensuremath{\boldsymbol{\times}} marks are there to prevent.


\section*{S4. The diagnostic decision tree --- branches B0--B7}


\subsection*{S4.1 Scope of the tree}


A faded cell arrives at the bench with no label. Deciding which mechanism caused the fade, and then whether that mechanism is coupled to another, means choosing among roughly twenty techniques, each of which constrains a different part of the problem. The order in which they are applied is the substance of this section.


The decision tree below is a sequencing of established measurements, not a new diagnostic. Every technique in it appears in Table S1 and none is novel. The contribution is the order: each step names one measurement, the \emph{discriminating outcome}, and the mechanism that outcome implies, so that a reader starting from a faded cell can reach a mechanism assignment (or an explicit \textquotedblleft{}confounded\textquotedblright{}) without having to hold the whole technique set in mind.


The tree does not attempt the following:


\begin{enumerate}
  \item It does not start from a mechanism. Beginning with the hypothesis \textquotedblleft{}M3\textquotedblright{} invites collecting evidence for M3. The tree starts from the observed symptom and forces the loss decomposition first, because a fade curve is a single scalar integrating every mechanism and cannot be evidence for any of them (Table S1, S3.1).
  \item It does not promise a unique answer. Three branches terminate in \textquotedblleft{}confounded\textquotedblright{} and say what would be needed to resolve them. A tree that resolved every branch would claim more than the measurements support; the residual confounds in S4.11 are an explicit part of the result.
  \item It does not identify coupling. Every branch up to B6 assigns a mechanism. The operative claim is about loops, and mechanism identity does not imply coupling, because two mechanisms can both be present and uncoupled. Branch \textbf{B7} exists to close that gap, and it is a perturbation protocol rather than a measurement.
\end{enumerate}


\textbf{Entry condition.} A cell that has faded measurably (\ensuremath{\geq}5 \% capacity loss or a detectable knee). For a cell that has not yet faded, start at \textbf{D2} (S4.12) instead, the early-warning branch.


\par\noindent\rule{\linewidth}{0.4pt}\par


\subsection*{S4.2 Node inventory and edge semantics}


The tree has 20 nodes: 10 actions and decisions, 10 terminals. The entry node is the faded cell; each decision names the observation that discriminates, and each terminal is either a mechanism assignment or an explicit \textquotedblleft{}confounded\textquotedblright{}.


\begingroup\small
\setlength{\tabcolsep}{3pt}
\renewcommand{\arraystretch}{1.12}
\sloppy
\begin{longtable}{@{}>{\raggedright\arraybackslash}p{0.173\textwidth}>{\raggedright\arraybackslash}p{0.610\textwidth}>{\raggedright\arraybackslash}p{0.133\textwidth}@{}}
\hline
\textbf{ID} & \textbf{Label} & \textbf{Section} \\ \hline
\endhead
B0 & Faded cell (\ensuremath{\geq}5 \% loss or knee) \ensuremath{\rightarrow} slow reference cycle + dV/dQ decomposition & S4.3 \\
D1 & Identification of the dominant loss term & S4.3 \\
B1 & LLI dominates \ensuremath{\rightarrow} HPC/CE + TGC (Li-metal) or \ensuremath{^{7}}Li NMR / 3-electrode reconstruction & S4.4 \\
B2 & Anode overpotential grew \ensuremath{\rightarrow} VRP after charge, then rest-period CE test & S4.5 \\
B3 & LAM dominates \ensuremath{\rightarrow} operando XRD across a full cycle & S4.6 \\
B4 & Structural change \ensuremath{\rightarrow} ICP-MS (electrolyte + anode digest) + stack-pressure differential & S4.7 \\
B5 & Polarization dominates \ensuremath{\rightarrow} EIS + DRT (not equivalent-circuit fitting) & S4.8 \\
B6 & Voltage fade only \ensuremath{\rightarrow} operando XRD + \ensuremath{^{7}}Li MAS-NMR & S4.9 \\
B7 & Loop identification \ensuremath{\rightarrow} knee-sensitivity perturbations & S4.10 \\
T-M1 \ldots{} T-M4 & M1 (SEI) \ensuremath{\cdot} M2 (dead Li) \ensuremath{\cdot} M3 (dissolution) \ensuremath{\cdot} M4 (fracture) & --- \\
T-TL & TL (thermodynamic/potential loss) \ensuremath{\rightarrow} energy balance, not capacity & S4.9 \\
T-LOOP & Loop confirmed \ensuremath{\rightarrow} name it (A/B/C/D) & S4.10 \\
T-NOLOOP & No loop \ensuremath{\rightarrow} evidence against the coupling framing; report it & S4.10 \\
T-CONF1\ldots{}3 & Confounded \ensuremath{\rightarrow} see S4.11 & S4.11 \\
D2 & No knee yet \ensuremath{\rightarrow} CE tracking, VRP, cathode DRT & S4.12 \\
\hline
\end{longtable}
\endgroup



\textbf{Edge labels are the discriminating outcomes, not descriptions of the action:} an edge out of D1 reads \textquotedblleft{}LLI dominates\textquotedblright{}, not \textquotedblleft{}went to B1\textquotedblright{}. This is what makes the tree a decision procedure. Every edge label is a physical observation a reader could actually make.


\textbf{Branch order is not arbitrary.} B1/B3/B5/B6 fan out from D1 in the order M1/M2 \ensuremath{\rightarrow} M4/M3 \ensuremath{\rightarrow} kinetics \ensuremath{\rightarrow} TL, so that a reader moving down the tree moves from quantity to location to chemistry, matching the three questions of S3.1. B7 sits below all of them as a single convergence point: every mechanism terminal feeds into it, because coupling is orthogonal to mechanism identity.


\par\noindent\rule{\linewidth}{0.4pt}\par


\subsection*{S4.3 B0: establishment of the loss terms}


\textbf{Measurement.} Run a slow (\ensuremath{\leq}C/20) reference cycle and compute d\emph{V}/d\emph{Q}. Decompose the fade into lithium inventory loss (LLI), loss of active material (LAM), and kinetic polarization.


Why this is step zero and not step one. The decomposition is the only measurement that addresses question 1 (\emph{which loss term moved}), and every downstream branch is conditioned on it. A tree that began with a chemistry probe would be guessing at which chemistry to probe. The \ensuremath{\leq}C/20 rate matters because at faster rates the kinetic contribution is confounded with the thermodynamic one and the d\emph{V}/d\emph{Q} features broaden into each other.


\textbf{Discriminating outcomes:}


\begin{itemize}
  \item LLI dominates \ensuremath{\rightarrow} \textbf{B1}
  \item LAM dominates \ensuremath{\rightarrow} \textbf{B3}
  \item Polarization dominates with little capacity loss \ensuremath{\rightarrow} \textbf{B5}
  \item Voltage fade with only modest capacity fade \ensuremath{\rightarrow} \textbf{B6}
\end{itemize}


\textbf{Confound resolved at this step.} A capacity fade reported as a single number. d\emph{V}/d\emph{Q} separates a peak shifting (inventory loss) from a peak shrinking (active-material loss), which is the canonical LLI/LAM split and is invisible in the fade curve. Without this step, \textquotedblleft{}the cell lost 20 \% capacity\textquotedblright{} cannot be attributed to any mechanism.


\textbf{What it cannot do.} It cannot say where Li was lost (SEI vs dead Li vs isolated pocket) nor which electrode lost active material. Those are B1/B2 and B3/B4 respectively. This limitation is structural, not a matter of care: d\emph{V}/d\emph{Q} integrates over both electrodes.


\par\noindent\rule{\linewidth}{0.4pt}\par


\subsection*{S4.4 B1: LLI dominance, M1 versus M2}


\textbf{Measurement.} High-precision coulometry for the parasitic current, and --- if the cell can be destroyed and is Li-metal --- titration gas chromatography (TGC). For sealed Li-ion cells, operando \ensuremath{^{7}}Li NMR if available; otherwise electrode-resolved three-electrode reconstruction of the anode overpotential.


\textbf{Discriminating outcomes:}


\begin{itemize}
  \item H\ensuremath{_{2}} yield high (Li\ensuremath{^{0}} dominates inactive Li) \ensuremath{\rightarrow} M2 (dead Li) \texttt{[FROM LITERATU}\hspace{0pt}\texttt{RE: 10.1038/}\hspace{0pt}\texttt{s41586-019-}\hspace{0pt}\texttt{1481-z]}
  \item SEI-bound Li\ensuremath{^{+}} dominates \ensuremath{\rightarrow} \textbf{M1}
  \item Anode overpotential grew but cathode did not \ensuremath{\rightarrow} M1/M2 on the anode \ensuremath{\rightarrow} \textbf{B2} to split them
  \item Cathode overpotential grew \ensuremath{\rightarrow} suspect \textbf{M3} even though the loss appears as LLI \ensuremath{\rightarrow} \textbf{B4}
\end{itemize}


\textbf{Confound resolved at this step.} The field's tendency to equate \textquotedblleft{}LLI\textquotedblright{} with \textquotedblleft{}SEI growth\textquotedblright{}. TGC's result is that in Li-metal cells unreacted Li\ensuremath{^{0}}, not SEI-bound Li\ensuremath{^{+}}, dominates inactive lithium, so an LLI-dominant cell is more likely a dead-lithium cell than an SEI cell. That re-partitions the two most-studied families. Note that Figure 1 panel (a) draws M1 and M2 as parallel sinks, both feeding the shared overpotential bus. The TGC result does not make one a downstream product of the other; it changes which of the two is expected to dominate in a Li-metal cell. The tree is what separates them; the figure does not pre-judge the split.


The outcome that carries the most diagnostic value. A cathode overpotential rise producing an LLI-appearing loss is the signature of M3 masquerading as M1: the cathode's impedance growth forces the anode to a more reducing potential, so the anode plates and loses inventory, and the loss presents on the LLI axis while its cause is at the other electrode. This is an instance of cross-electrode coupling argued for here, and it is a trap that a mechanism-by-mechanism analysis walks directly into.


\textbf{Confound not resolved at this step.} In a sealed Li-ion cell there is no non-destructive equivalent of TGC, and high-precision coulometry gives the total parasitic current but not its partition between direct SEI growth and corrosion of irreversibly plated Li (which becomes SEI). This is the field's principal diagnostic gap --- see S4.11.


\par\noindent\rule{\linewidth}{0.4pt}\par


\subsection*{S4.5 B2: separation of M1 from M2 at the anode}


\textbf{Measurement.} First a voltage-relaxation profile (VRP) immediately after charge; then the rest-period CE test.


\textbf{Discriminating outcomes, and a trap that must be stated in the figure caption:}


\begin{itemize}
  \item A voltage plateau during relaxation (mixed potential from Li\ensuremath{^{0}} stripping) \ensuremath{\rightarrow} reversible plating present \ensuremath{\rightarrow} \textbf{M2}
  \item No plateau \ensuremath{\rightarrow} this rules out only reversible plating. It does not rule out dead Li: an electrically isolated Li\ensuremath{^{0}} deposit cannot strip, so it produces no relaxation signal at all. Do not assign M1 on a negative VRP.
  \item Then the rest-period CE test, which distinguishes the irreversible channels: CE improves after rest \ensuremath{\rightarrow} the cell was plating-limited \ensuremath{\rightarrow} \textbf{M2} (plated Li partly recovered); CE degrades monotonically across rest \ensuremath{\rightarrow} \textbf{M1} (SEI-limited, no recoverable reservoir); CE approximately invariant to rest \ensuremath{\rightarrow} neither channel dominates the anodic loss \ensuremath{\rightarrow} return to \textbf{B4}.
\end{itemize}


\textbf{Confound resolved at this step.} The false inference \textquotedblleft{}no plating signature, therefore SEI.\textquotedblright{} Dead lithium is diagnostically silent in the relaxation measurement, so absence of evidence is genuinely not evidence of absence here. Many mechanism assignments in the literature rest on exactly this invalid inference, and naming it is one of the tree's specific contributions.


\textbf{Residual confound.} These rest-period tests separate reversible from irreversible loss but still cannot split \textquotedblleft{}irreversibly plated Li that then corroded into SEI\textquotedblright{} from \textquotedblleft{}direct SEI growth.\textquotedblleft{} That requires TGC or \ensuremath{^{7}}Li NMR. This is the same gap as S4.4's, approached from the other side, and it is the reason the M1/M2 partition is not conserved over a cell's life.


\par\noindent\rule{\linewidth}{0.4pt}\par


\subsection*{S4.6 B3: LAM dominance, electrode and structure}


\textbf{Measurement.} Operando XRD across a full cycle. The lattice-parameter test is conclusive, because it is the one thing M3 and M4 do not share.


\textbf{Discriminating outcomes:}


\begin{itemize}
  \item Cathode lattice parameters shift irreversibly, or new spinel/rock-salt reflections appear \ensuremath{\rightarrow} chemical change \ensuremath{\rightarrow} M3, not M4 \ensuremath{\rightarrow} B4 to confirm dissolution
  \item Cathode reflections unchanged but capacity lost \ensuremath{\rightarrow} \textbf{M4} (fracture/isolation) or \textbf{M2} (isolated Li pockets); separate by the time law below
  \item Anode staging capacity lost with unchanged lattice \ensuremath{\rightarrow} anode LAM (particle isolation, M4-driven)
  \item Both electrodes retain lattice but capacity is gone \ensuremath{\rightarrow} electrically disconnected material \ensuremath{\rightarrow} \textbf{M4} (contact loss) or \textbf{M2} (isolated Li pockets)
\end{itemize}


\textbf{Confound resolved at this step.} M3 versus M4 --- two mechanisms that both destroy active material and both look identical in a fade curve. Pure mechanical fracture leaves lattice parameters intact; dissolution and reconstruction change them (notably the \emph{c}-parameter). One measurement therefore separates them, and it is an argument that can be stated briefly.


Time-law discriminator, applied before any LAM assignment. Fit the loss as \ensuremath{\Delta}\emph{C} = \emph{a}\ensuremath{\surd}\emph{t} + \emph{b}\ensuremath{\cdot}\emph{t} over a high-quality aging series:


\begin{itemize}
  \item \emph{b}/\emph{a} \ensuremath{\approx} 0 (square-root growth) \ensuremath{\rightarrow} transport-limited interphase growth \ensuremath{\rightarrow} \textbf{M1}
  \item measurable \emph{b} term (linear in \emph{area creation rate}, not time) \ensuremath{\rightarrow} mechanically driven fresh-surface re-formation \ensuremath{\rightarrow} \textbf{M4}
\end{itemize}


The ratio \emph{b}/\emph{a} is a direct if crude measure of the mechanical contribution, and it costs nothing beyond a well-sampled fade curve, which is why it belongs in the tree before any instrument time is spent.


\textbf{Two M4 regimes require separation.} A C-rate sweep sees only \emph{gradient-stress} fracture (rate-dependent). \emph{Coherency-stress} (\textquotedblleft{}electrochemical shock\textquotedblright{}) fracture is nearly rate-independent and occurs in phase-separating materials, so it proceeds even in slow cycling and can mimic a calendar-driven mechanism. Therefore a null result from a C-rate sweep does not rule out M4. Always pair the rate sweep with a calendar control at matched SOC and temperature.


\par\noindent\rule{\linewidth}{0.4pt}\par


\subsection*{S4.7 B4: M3 dissolution versus M4 fracture}


\textbf{Measurement.} ICP-MS on the electrolyte and on the anode digest, read together with B3's lattice test.


\textbf{Discriminating outcomes:}


\begin{itemize}
  \item Significant Mn/Ni/Co detected on the anode or in electrolyte \ensuremath{\rightarrow} M3 confirmed (dissolution happened)
  \item Negligible TM detected, but cracking/expansion seen with lattice intact \ensuremath{\rightarrow} \textbf{M4}
  \item TM detected and cracking seen \ensuremath{\rightarrow} M3 + M4 coupled; use the stack-pressure differential test: raise stack pressure. If lifetime improves with pressure, the M4\ensuremath{\rightarrow}M3 arrow is active.
  \item Then separate dissolution-driven from reconstruction-driven fade: compare single-crystal vs polycrystalline of the same composition. M4 largely disappears; if fade persists, the residual is \textbf{M3 proper}.
\end{itemize}


\textbf{Confound resolved at this step.} \textquotedblleft{}TM was detected\textquotedblright{} is routinely reported as \textquotedblleft{}dissolution caused the fade.\textquotedblright{} \textbf{Confirmed dissolution is not confirmed causation} --- it proves TM has physically left the cathode, which is a mass-loss fact, not a mechanistic one. The single-crystal control is what separates the dissolution channel from the fracture channel that exposes fresh surface for further dissolution.


\textbf{Why the M3+M4 pair matters.} The M4\ensuremath{\rightarrow}M3 arrow is bidirectional and self-reinforcing: crack faces are unpassivated oxide, and the lattice-parameter collapse from TM loss feeds more cracking. That bidirectionality is a loop, and it illustrates why the framework is drawn as arrows rather than as a pathway.


\par\noindent\rule{\linewidth}{0.4pt}\par


\subsection*{S4.8 B5: polarization dominance and the resistance location}


\textbf{Measurement.} EIS with distribution-of-relaxation-times (DRT) deconvolution rather than equivalent-circuit fitting, because the arcs overlap.


\textbf{Discriminating outcomes:}


\begin{itemize}
  \item High-frequency arc grows \ensuremath{\rightarrow} interphase resistance. On the cathode \ensuremath{\rightarrow} CEI growth or M3 reconstruction; on the anode \ensuremath{\rightarrow} M1 SEI growth, or M2 plating onset (operando EIS detects plating down to <0.6 \% of graphite capacity via a rise in graphite SEI resistance \texttt{[FROM LITERATU}\hspace{0pt}\texttt{RE: 10.1016/j.}\hspace{0pt}\texttt{xcrp.2021.}\hspace{0pt}\texttt{100589]})
  \item Mid-frequency arc grows \ensuremath{\rightarrow} charge-transfer limitation (M3 reconstruction, or electrolyte depletion)
  \item Low-frequency / diffusion tail grows \ensuremath{\rightarrow} transport limitation. If it scales with electrode thickness/areal capacity independently of chemistry, suspect Loop D rather than M4 --- the distinguishing test is whether a flooded cell shows the same fade, which a Loop-D-limited cell will not
  \item Resolve anode vs cathode by three-electrode. Positioning matters: the reference electrode must be placed by modeling, not by convenience, or the measured half-cell overpotentials are artefacts.
\end{itemize}


\textbf{Confound resolved at this step.} Equivalent-circuit fitting assigns overlapping arcs arbitrarily, and the assignment of a specific arc to M1 vs M3 is model-dependent. DRT removes the circuit-topology choice, which is why it is specified rather than left optional. The second confound resolved is regional: an impedance rise with no capacity loss is a different object from capacity fade, and the tree separates them at B5 rather than lumping them.


Why the flooded-cell test is the discriminator for Loop D. Loop D is a volumetric budget (S2.6): a lean cell is salt-starved at \textasciitilde{}9\ensuremath{\times} overdraw, a flooded cell has \textasciitilde{}15\ensuremath{\times} more inventory. So a fade that disappears when the cell is flooded is a Loop-D fade, and one that persists is not. This is a \emph{design-of-experiment} discriminator, not a measurement, which is why Table S1 has no row for it.


\par\noindent\rule{\linewidth}{0.4pt}\par


\subsection*{S4.9 B6: the voltage-fade branch}


\textbf{Measurement.} If the symptom is voltage fade with modest capacity fade: operando XRD + \ensuremath{^{7}}Li MAS-NMR, looking for tetrahedral-site occupancy change and spinel/rock-salt phase growth.


\textbf{Discriminating outcomes:}


\begin{itemize}
  \item Tetrahedral occupancy change / spinel growth \ensuremath{\rightarrow} voltage fade \ensuremath{\rightarrow} TL term, not KL
  \item If voltage fade and capacity fade are decoupled in the data, report that explicitly --- it is a result about separability, not a measurement failure
\end{itemize}


\textbf{Confound resolved at this step.} Folding voltage fade into the kinetic-loss channel. Voltage fade is a thermodynamic/potential loss (the fourth accounting channel, TL): the cell stores less energy at the same capacity and the same impedance. Assigning it to KL would double-count it against the EIS measurement in B5, which measures a physically distinct resistive loss with a different temperature dependence. Separating them is what makes the capacity-balance claims well-posed.


\textbf{The scope statement this branch carries.} The reported study shows that voltage fade and capacity fade can be decoupled (a 723.5--769.6 mV midpoint drop after 200 cycles alongside high capacity retention, and a 491 mV shift in one sample) \texttt{[FROM LITERATU}\hspace{0pt}\texttt{RE: 10.1039/}\hspace{0pt}\texttt{D4SC07037B]}, which is why TL is required for energy while the LLI/LAM/KL triple is complete \emph{for capacity}. The tree's B6 branch therefore returns \emph{TL}, not KL. The limitation is that TL sits outside the loop framework: three of the four loops act on the capacity balance, and voltage fade remains a named open problem rather than something the loops explain.


\par\noindent\rule{\linewidth}{0.4pt}\par


\subsection*{S4.10 B7: loop identification}


This is the branch that distinguishes the tree from a list.


\textbf{The problem B7 solves.} Every branch above assigns a mechanism. The Perspective's claim is about coupling. Two mechanisms can both be present, both be measured, and still be uncoupled --- each fading the cell on its own without either driving the other. Mechanism identity therefore does not imply a loop, and it would be wrong to infer one. B7 is the test that separates them, and it is required because no measurement in Table S1 measures coupling directly (S3.5).


\textbf{The diagnostic logic.} A single mechanism gives a smooth, monotone fade. A positive-feedback loop gives an acceleration whose onset moves with the driving variable. So:


\begin{enumerate}
  \item Locate the acceleration. Compute the local fade-rate derivative d(\ensuremath{\Delta}\emph{C})/d\emph{N} and find its inflection: that cycle is the loop's onset, not the mechanism's onset. This is the same departure-from-baseline idea the reduced model uses (S1.5), applied to data rather than a simulation.
  \item Perturb one loop-specific driving variable and test whether the knee moves:
\end{enumerate}


\begingroup\small
\setlength{\tabcolsep}{3pt}
\renewcommand{\arraystretch}{1.12}
\sloppy
\begin{tabular}{@{}>{\raggedright\arraybackslash}p{0.347\textwidth}>{\raggedright\arraybackslash}p{0.295\textwidth}>{\raggedright\arraybackslash}p{0.274\textwidth}@{}}
\hline
\textbf{Loop} & \textbf{Perturbation} & \textbf{Expected if the loop is real} \\ \hline
A --- plating \ensuremath{\rightarrow} dead Li \ensuremath{\rightarrow} SEI \ensuremath{\rightarrow} \ensuremath{\eta} \ensuremath{\rightarrow} plating & Raise temperature, or lower C-rate & Knee moves sharply; a purely M3/M4 knee does not \\
\hline
B --- dissolution \ensuremath{\rightarrow} anode poisoning \ensuremath{\rightarrow} \ensuremath{\eta} \ensuremath{\rightarrow} dissolution & Add an HF scavenger, or switch to a coating & Knee moves; an M4 knee does not \\
\hline
C --- fracture \ensuremath{\rightarrow} fresh surface \ensuremath{\rightarrow} SEI/dissolution & Raise stack pressure & Moves the M4\ensuremath{\rightarrow}M3 leg \\
\hline
D --- electrolyte depletion / porosity clogging & Flood the cell with extra electrolyte & Only a Loop-D knee responds \\
\hline
\end{tabular}
\endgroup



\begin{enumerate}
  \setcounter{enumi}{2}
  \item Read the outcome. A knee whose onset moves under a loop-specific perturbation is evidence for that loop. A knee that does not move under any of the four perturbations is evidence against the loop framing for that cell --- report it.
\end{enumerate}


Why the negative terminal is marked explicitly. The tree carries its own falsification test. The loops predict parameter sensitivity of the knee onset, so an insensitive knee is a counterexample, not a null result. That terminal is what distinguishes a framework from a taxonomy: it hands the reader the experiment that could refute the claim.


\textbf{Quantitative anchor from the reduced model.} S1's ablation gives the expected \emph{sign and magnitude} of the perturbation effect: severing the Loop-A arrow abolishes the knee entirely rather than delaying it (coupled departure at cycle 285, severed: none within 6000), and the knee moves monotonically with rate (2651 \ensuremath{\rightarrow} 64 cycles from 0.5C to 2.5C, a factor of 41.4). A real cell that shows no knee displacement under a 5\ensuremath{\times} rate change is quantitatively inconsistent with a Loop-A knee, though the model's absolute numbers are illustrative and only the \emph{direction and order of sensitivity} should be used.


\textbf{What B7 cannot claim.} The perturbations are not orthogonal in a real cell --- raising temperature changes SEI growth, charge transfer, and dissolution simultaneously; flooding a cell changes pressure as well as inventory. So a knee that moves does not by itself identify which loop moved it, and a full assignment needs at least two perturbations with different fingerprints. B7 is a test for \emph{whether coupling is present}, not a method for \emph{which coupling}, and the distinction is stated explicitly.


\par\noindent\rule{\linewidth}{0.4pt}\par


\subsection*{S4.11 Residual confounds}


These are the places where separability is incomplete. Each is a gray \textquotedblleft{}confounded\textquotedblright{} terminal in the node inventory of S4.2.


\begin{enumerate}
  \item M1 vs M2 in sealed Li-ion cells, the field's principal diagnostic gap. There is no non-destructive equivalent of TGC for a commercial cell. High-precision coulometry gives the total parasitic current but not its partition between direct SEI growth and the corrosion of irreversibly plated Li (which becomes SEI). Note what is available: for Li-metal/anode-free cells TGC gives absolute Li\ensuremath{^{0}} and operando \ensuremath{^{7}}Li NMR is non-destructive; for graphite/Li-ion cells no method currently counts Li\ensuremath{^{0}} against SEI-Li\ensuremath{^{+}} without destroying the cell. Voltage relaxation, high-precision CE and operando EIS all detect plating, not the \emph{dead-Li fraction}. This is the gap the tree cannot close, and it is stated explicitly rather than passed over.
  \item M3 vs M1 at the anode interface. Deposited TM changes SEI composition and makes the interphase more electronically conductive, so \textquotedblleft{}anode SEI degradation\textquotedblright{} in a cell that also has a dissolving cathode is partly M3. Element-resolved XPS/ICP helps but is surface-only or destructive. This is the same failure-mode the TGC result exposes on the Li-metal side, appearing again at the anode of a Li-ion cell.
  \item M3 reconstruction vs CEI growth on impedance. Both appear in the same frequency range. Only depth-resolved structural probes (STEM-EELS, grazing-incidence XRD) can separate them.
  \item M4 vs M2 under pressure, a two-sided confound. Stack pressure simultaneously suppresses cracking and changes plating behavior, so a pressure series moves the two in opposite directions. Worse, the sign of the M2 pressure effect is itself unresolved: beneficial pressure (75--1200 kPa) \texttt{[FROM LITERATU}\hspace{0pt}\texttt{RE: 10.1149/2.}\hspace{0pt}\texttt{0091908jes]} is far below the Li yield stress, and creep is the plausible but undemonstrated explanation. A pressure series must be interpreted with both mechanisms in the model, not as a clean M4 assay. This directly limits B7's Loop-C perturbation, and the figure caption should say so.
  \item M4's regimes. Gradient-stress fracture is rate-dependent; coherency-stress (\textquotedblleft{}electrochemical shock\textquotedblleft{}) fracture is nearly rate-independent. A C-rate sweep is therefore a valid M4 assay only for the first regime and will produce a false negative for the second. Always pair it with a calendar control.
  \item Voltage fade vs capacity fade --- resolvable, and now a stated scope boundary rather than an open problem. The two can be decoupled; TL captures the difference for energy. B6 is read as returning \emph{TL}, not KL. What remains genuinely open is why the layered\ensuremath{\rightarrow}spinel transition and voltage fade are not one-to-one: the extent of the transition does not by itself determine the voltage fade.
\end{enumerate}


\par\noindent\rule{\linewidth}{0.4pt}\par


\subsection*{S4.12 D2: the early-warning branch}


Enter here only for a cell that has not yet faded measurably. The tree's main body assumes \ensuremath{\geq}5 \% loss; D2 is what to do before that threshold.


\textbf{Actions and outcomes:}


\begin{itemize}
  \item Track CE at high precision. A sustained CE below roughly 0.9995 is the earliest usable signal. That is a screening level, looser than the 99.978 \% that S2.4 shows is needed for 1000 cycles at 80 \% retention; it is set by the instrument's own noise floor and by the number of cycles over which CE is averaged, and should be reported as such rather than quoted as a constant. A cell that trips it is screened for the full branch; a cell that passes it is not thereby cleared.
  \item Track voltage relaxation after every charge. A growing plateau precedes measurable capacity loss.
  \item Track cathode impedance (DRT mid-frequency). M3 raises it before capacity moves.
\end{itemize}


\textbf{Why this branch exists.} The practical claim is that knee onset is a design variable and therefore partly predictable. A diagnostic tree that operates only after the knee has arrived would concede that claim. D2 tests whether the logic reaches earlier than the capacity curve; it does so by a modest margin, using CE and relaxation, not by a wide one.


\par\noindent\rule{\linewidth}{0.4pt}\par


\subsection*{S4.13 Branch summary}


\begingroup\small
\setlength{\tabcolsep}{3pt}
\renewcommand{\arraystretch}{1.12}
\sloppy
\begin{tabular}{@{}>{\raggedright\arraybackslash}p{0.107\textwidth}>{\raggedright\arraybackslash}p{0.356\textwidth}>{\raggedright\arraybackslash}p{0.454\textwidth}@{}}
\hline
\textbf{Branch} & \textbf{Discrimination settled} & \textbf{Confound resolved} \\ \hline
\textbf{B0} & Identity of the dominant loss term & A single fade number hiding LLI/LAM/kinetic partitions \\
\hline
\textbf{B1} & Partition of lost Li between SEI and dead Li & Equating LLI with SEI growth \\
\hline
\textbf{B2} & Separation of M1 from M2 at the anode & \textquotedblleft{}No relaxation plateau, therefore SEI\textquotedblright{} (dead Li is diagnostically silent) \\
\hline
\textbf{B3} & Electrode attribution of LAM, and its structural character & M3 vs M4 --- both destroy AM, only one moves the lattice \\
\hline
\textbf{B4} & Attribution between dissolution and fracture & \textquotedblleft{}TM detected\textquotedblright{} mistaken for \textquotedblleft{}dissolution caused the fade\textquotedblright{} \\
\hline
\textbf{B5} & Location of the dominant resistance & Arbitrary arc assignment from equivalent-circuit fitting; local vs regional loss \\
\hline
\textbf{B6} & Partition between voltage and capacity loss & Folding thermodynamic voltage fade into kinetic loss (double-counting TL as KL) \\
\hline
\textbf{B7} & Presence of a \emph{loop} rather than mechanisms alone & Mechanism identity mistaken for coupling \\
\hline
\textbf{D2} & Detection of onset before the capacity curve moves & Waiting for the capacity curve before acting \\
\hline
\end{tabular}
\endgroup



\textbf{In one sentence:} the tree proceeds quantity \ensuremath{\rightarrow} location \ensuremath{\rightarrow} chemistry \ensuremath{\rightarrow} coupling, because each step's confound is invisible to the step before it, and it terminates in a test designed to be able to fail.

\end{document}
